\newMonth{September}
\setrunningtitles{September}{September}

% ---- martyrology/mart09/mart0901.htm
\needspace{10\baselineskip}
\begin{paracol}{2}
\selectlanguage{latin}
\begin{center}{\color{gregoriocolor} Kaléndis Septémbris. 
 Luna\dots\ }\end{center}
\switchcolumn
\selectlanguage{english}
\begin{center}{\color{gregoriocolor} The   1\textsuperscript{st} Day of 
 September. The\dots\ Day of the Moon.}\end{center}
\end{paracol}

\noindent\begin{tabularx}{\linewidth}{*{19}{>{\centering\arraybackslash}X}}
 \textcolor{gregoriocolor}{a} & \textcolor{gregoriocolor}{b} & \textcolor{gregoriocolor}{c} & \textcolor{gregoriocolor}{d} & \textcolor{gregoriocolor}{e} & \textcolor{gregoriocolor}{f} & \textcolor{gregoriocolor}{g} & \textcolor{gregoriocolor}{h} & \textcolor{gregoriocolor}{i} & \textcolor{gregoriocolor}{k} & \textcolor{gregoriocolor}{l} & \textcolor{gregoriocolor}{m} & \textcolor{gregoriocolor}{n} & \textcolor{gregoriocolor}{p} & \textcolor{gregoriocolor}{q} & \textcolor{gregoriocolor}{r} & \textcolor{gregoriocolor}{s} & \textcolor{gregoriocolor}{t} & \textcolor{gregoriocolor}{u} \\
 9 & 10 & 11 & 12 & 13 & 14 & 15 & 16 & 17 & 18 & 19 & 20 & 21 & 22 & 23 & 24 & 25 & 26 & 27 \\
\end{tabularx}
\vspace{0.5\baselineskip}
\noindent\begin{tabularx}{\linewidth}{*{12}{>{\centering\arraybackslash}X}}
 \textcolor{gregoriocolor}{A} & \textcolor{gregoriocolor}{B} & \textcolor{gregoriocolor}{C} & \textcolor{gregoriocolor}{D} & \textcolor{gregoriocolor}{E} & F & \textcolor{gregoriocolor}{F} & \textcolor{gregoriocolor}{G} & \textcolor{gregoriocolor}{H} & \textcolor{gregoriocolor}{M} & \textcolor{gregoriocolor}{N} & \textcolor{gregoriocolor}{P} \\
 28 & 29 & 30 & 1 & 2 & 3 & 3 & 4 & 5 & 6 & 7 & 8 \\
\end{tabularx}

\begin{paracol}{2}
\selectlanguage{latin}
\lettrine[lines=2]{I}{n} província Narbonénsi 
 sancti Ægídii, Abbátis et Confessóris, cujus nómine est appellátum óppidum, 
 quod póstea crevit in loco, ubi ipse monastérium eréxerat et mortális vitæ 
 cursum absólverat.
\switchcolumn
\selectlanguage{english}
\lettrine[lines=2]{I}{n} the province of Narbonne, St. 
 Giles, abbot and confessor. A town which later arose in the place 
 where he had built his monastery and where he died was named after him.
\switchcolumn*
\selectlanguage{latin}
Sentiáni, in fínibus 
 Apúliæ, pássio sanctórum Donáti et altérius Felícis; qui, sanctórum 
 Bonifátii et Theclæ fílii, a Valeriáno Júdice, sub Maximiáno Imperatóre, 
 jussi sunt, post vária torménta, cápite præcídi hodiérna die, in qua et 
 festívitas aliórum ex duódecim frátribus, quorum natális respectívis diébus 
 ágitur, institúta est celebrári. Ipsórum vero duódecim fratrum córpora 
 Benevéntum póstea transláta sunt, ibíque honorífice asserváta.
\switchcolumn
\selectlanguage{english}
At Sentiano, in the district of 
 Apulia, the passion of Saints Donatus and a second Felix who were the sons 
 of Saints Boniface and Thecla. After they had endured various torments 
 under the judge Valerian in the reign of Emperor Maximian, they were 
 condemned to be beheaded on this day. Today also is kept the festival 
 of the others of the Twelve Holy Brethren, whose birthdays are noted in 
 their proper place. The bodies of these Twelve Holy Brethren were 
 later translated to Benevento where they are honourably enshrined.
\switchcolumn*
\selectlanguage{latin}
In Palæstína sanctórum 
 Jósue et Gedeónis.
\switchcolumn
\selectlanguage{english}
In Palestine, the Saints Joshua and 
 Gideon.
\switchcolumn*
\selectlanguage{latin}
Hierosólymis beátæ Annæ 
 Prophetíssæ, cujus sanctitátem sermo Evangélicus prodit.
\switchcolumn
\selectlanguage{english}
At Jerusalem, blessed Anna, 
 prophetess, whose sanctity is revealed in the Gospel.
\switchcolumn*
\selectlanguage{latin}
Cápuæ, via Aquária, 
 sancti Prisci Mártyris, qui fuit unus de antíquis Christi discípulis.
\switchcolumn
\selectlanguage{english}
At Capua, on the Via Aquaria, St. 
 Priscus, martyr, who was formerly one of the disciples of Christ.
\switchcolumn*
\selectlanguage{latin}
Tudérti, in Umbria, 
 sancti Terentiáni, Epíscopi et Mártyris; qui, sub Hadriáno Imperatóre, 
 Lætiáni Procónsulis jussu, equúleo et scorpiónibus cruciátus est, ac demum, 
 abscíssa lingua, cápitis damnátus martyrium complévit.
\switchcolumn
\selectlanguage{english}
At Todi in Umbria, St. Terentian, 
 bishop and martyr. Under Emperor Hadrian, by order of the proconsul 
 Laetian, he was racked, scourged with whips set with metal, and finally 
 having had his tongue cut out, he ended his martyrdom by undergoing capital 
 punishment.
\switchcolumn*
\selectlanguage{latin}
Heracléæ, in Thrácia, 
 sancti Ammónis Diáconi, et sanctárum quadragínta Vírginum, quas ille 
 erudívit in fide, et, sub Licínio tyránno, ad martyrii glóriam secum 
 perdúxit.
\switchcolumn
\selectlanguage{english}
At Heraclea, under the tyrant 
 Licinius, St. Ammon, deacon, and forty holy virgins whom he instructed in 
 the faith and led with him to the glory of martyrdom.
\switchcolumn*
\selectlanguage{latin}
In Hispániæ sanctórum 
 Mártyrum Vincéntii et Læti.
\switchcolumn
\selectlanguage{english}
In Spain, the holy martyrs Vincent 
 and Laetus.
\switchcolumn*
\selectlanguage{latin}
Populónii, in Túscia, 
 sancti Réguli Mártyris, qui ex Africa illuc venit, ibíque, sub Tótila, 
 martyrium consummávit.
\switchcolumn
\selectlanguage{english}
At Piombino in Tuscany, St. Regulus, 
 martyr, who went thither from Africa, and consummated his martyrdom under 
 Totila.
\switchcolumn*
\selectlanguage{latin}
Cápuæ sancti Prisci 
 Epíscopi, qui unus fuit ex illis Sacerdótibus, qui, in persecutióne 
 Wandalórum, ob fidem cathólicam várie afflícti et vetústæ navi impósiti, ex 
 Africa ad Campániæ líttora pervenérunt, et Christiánam religiónem, in iis 
 locis dispérsi diversísque Ecclésiis præfécti, mirífice propagárunt. 
 Ipsíus autem fuérunt sócii Castrénsis, cujus dies natális tértio Idus 
 Februárii recólitur, Támmarus, Rósius, Heráclius, Secundínus, Adjútor, 
 Marcus, Augústus, Elpídius, Cánion et Vindónius.
\switchcolumn
\selectlanguage{english}
At Capua, St. Priscus, bishop. 
 He was one of those priests who were subjected to various trials for the 
 Catholic faith during the persecution of the Vandals. Being put in an 
 old ship on the coast of Africa, they reached the shores of Campania, and 
 separating, they were placed at the head of various churches, and thus 
 greatly extended the Christian religion. The companions of Priscus 
 were Castrensis, whose birthday is mentioned on the 11th of February, 
 Tammarius, Rosius, Heraclius, Secundinus, Adjutor, Mark, Augustus, Elpidius, 
 Canion, and Vindonius.
\switchcolumn*
\selectlanguage{latin}
Apud Sénonas beáti Lupi, 
 Epíscopi et Confessóris; de quo refértur quod quadam die, præsénte Clero, 
 dum sacris altáribus adstáret, lapsa est cælitus gemma in ejus cálicem 
 sanctum.
\switchcolumn
\selectlanguage{english}
At Sens, St. Lupus, bishop and 
 confessor, of whom it is related that on a certain day, while he stood at 
 the holy altar in the presence of the clergy, a gem fell from heaven into 
 the consecrated chalice which he was using.
\switchcolumn*
\selectlanguage{latin}
Rhemis, in Gállia, 
 sancti Xysti, qui fuit primus ejúsdem civitátis Epíscopus.
\switchcolumn
\selectlanguage{english}
At Rheims in France, St. Sixtus, 
 disciple of the blessed apostle Peter, who consecrated him the first bishop 
 of that city. He received the crown of martyrdom under Nero.
\switchcolumn*
\selectlanguage{latin}
Apud Cenómanos, in 
 Gállia, sancti Victórii Epíscopi.
\switchcolumn
\selectlanguage{english}
At Le Mans in France, St. Victorinus, 
 bishop.
\switchcolumn*
\selectlanguage{latin}
Apud Aquínum sancti 
 Constántii Epíscopi, prophetíæ dono multísque virtútibus clari.
\switchcolumn
\selectlanguage{english}
At Aquino, St. Constantius, a bishop 
 renowned for the gift of prophecy and many virtues.
\switchcolumn*
\selectlanguage{latin}
Ad Aquas Duras, in 
 Constantiénsi Germániæ território, sanctæ Verénæ Vírginis.
\switchcolumn
\selectlanguage{english}
In Baden, in the province of 
 Constance, St. Verena, virgin.
\switchcolumn*
\selectlanguage{latin}
\end{paracol}

% \continueMonth{September}

% ---- martyrology/mart09/mart0902.htm
\needspace{10\baselineskip}
\begin{paracol}{2}
\selectlanguage{latin}
\begin{center}{\color{gregoriocolor} Quarto Nonas Septémbris. 
 Luna\dots\ }\end{center}
\switchcolumn
\selectlanguage{english}
\begin{center}{\color{gregoriocolor} The 2\textsuperscript{nd} Day of September. The\dots\ Day 
 of the Moon.}\end{center}
\end{paracol}

\noindent\begin{tabularx}{\linewidth}{*{19}{>{\centering\arraybackslash}X}}
 \textcolor{gregoriocolor}{a} & \textcolor{gregoriocolor}{b} & \textcolor{gregoriocolor}{c} & \textcolor{gregoriocolor}{d} & \textcolor{gregoriocolor}{e} & \textcolor{gregoriocolor}{f} & \textcolor{gregoriocolor}{g} & \textcolor{gregoriocolor}{h} & \textcolor{gregoriocolor}{i} & \textcolor{gregoriocolor}{k} & \textcolor{gregoriocolor}{l} & \textcolor{gregoriocolor}{m} & \textcolor{gregoriocolor}{n} & \textcolor{gregoriocolor}{p} & \textcolor{gregoriocolor}{q} & \textcolor{gregoriocolor}{r} & \textcolor{gregoriocolor}{s} & \textcolor{gregoriocolor}{t} & \textcolor{gregoriocolor}{u} \\
 10 & 11 & 12 & 13 & 14 & 15 & 16 & 17 & 18 & 19 & 20 & 21 & 22 & 23 & 24 & 25 & 26 & 27 & 28 \\
\end{tabularx}
\vspace{0.5\baselineskip}
\noindent\begin{tabularx}{\linewidth}{*{12}{>{\centering\arraybackslash}X}}
 \textcolor{gregoriocolor}{A} & \textcolor{gregoriocolor}{B} & \textcolor{gregoriocolor}{C} & \textcolor{gregoriocolor}{D} & \textcolor{gregoriocolor}{E} & F & \textcolor{gregoriocolor}{F} & \textcolor{gregoriocolor}{G} & \textcolor{gregoriocolor}{H} & \textcolor{gregoriocolor}{M} & \textcolor{gregoriocolor}{N} & \textcolor{gregoriocolor}{P} \\
 29 & 30 & 1 & 2 & 3 & 4 & 4 & 5 & 6 & 7 & 8 & 9 \\
\end{tabularx}

\begin{paracol}{2}
\selectlanguage{latin}
\lettrine[lines=2]{S}{ancti} Stéphani, Regis 
 Hungarórum et Confessóris; qui décimo octávo Kaléndas Septémbris obdormívit 
 in Dómino.
\switchcolumn
\selectlanguage{english}
\lettrine[lines=2]{S}{t.} Stephen, king of Hungary and 
 confessor, who fell asleep in the Lord on the 15th of August.
\switchcolumn*
\selectlanguage{latin}
Romæ sanctæ Máximæ 
 Mártyris, quæ, simul cum sancto Ansáno Christum conféssa, in persecutióne 
 Diocletiáni, dum fústibus cæditur, réddidit spíritum.
\switchcolumn
\selectlanguage{english}
At Rome, the holy martyr Maxima, who 
 confessed Christ with St. Ansanus in the persecution of Diocletian, and 
 yielded up her soul while being beaten with rods.
\switchcolumn*
\selectlanguage{latin}
Pámiæ, in Gállia, 
 sancti Antoníni Mártyris, cujus relíquiæ apud Ecclésiam Palentínam, in 
 Hispánia, magna veneratióne asservántur.
\switchcolumn
\selectlanguage{english}
At Pamiers in France, St. Antoninus, 
 martyr, whose relics are kept with great veneration in the church of 
 Palencia, in Spain.
\switchcolumn*
\selectlanguage{latin}
Item sanctórum Mártyrum 
 Diomédis, Juliáni, Philíppi, Euthychiáni, Hesychii, Leónidæ, Philadélphi, 
 Menalíppi et Pantágapæ; quorum álii igne, álii aqua, ense álii et cruce 
 martyrium complevérunt.
\switchcolumn
\selectlanguage{english}
Also, the holy martyrs, Diomedes, 
 Julian, Philip, Eutychian, Hesychius, Leonides, Philadelphus, Menalippus, 
 and Pantagapas. Their martyrdoms were completed, some by fire, some 
 water, others by the sword or by the cross.
\switchcolumn*
\selectlanguage{latin}
Nicomedíæ sanctórum 
 Mártyrum Zenónis, atque Concórdii et Theodóri, filiórum ejus.
\switchcolumn
\selectlanguage{english}
At Nicomedia, the holy martyrs Zeno, 
 and his sons Concordius and Theodore.
\switchcolumn*
\selectlanguage{latin}
Lugdúni, in Gállia, 
 sancti Elpídii, Epíscopi et Confessóris.
\switchcolumn
\selectlanguage{english}
At Lyons in France, St. Elpidius, 
 bishop and confessor.
\switchcolumn*
\selectlanguage{latin}
In Picéno item sancti 
 Elpídii Abbátis, cujus nómine óppidum est appellátum, quod ejus sacrum 
 corpus se possidére congáudet.
\switchcolumn
\selectlanguage{english}
In Piceno, another St. Elpidius, an 
 abbot. A town bearing his name glories in the possession of his holy 
 body.
\switchcolumn*
\selectlanguage{latin}
In monte Sorácte sancti 
 Nonnósi Abbátis, qui ingéntis molis saxum oratióne sua tránstulit, aliísque 
 miráculis coruscávit.
\switchcolumn
\selectlanguage{english}
On Mount Soracte, Abbot St. Nonnosus, 
 who by his prayers moved a rock of huge proportions, and was renowned for 
 other miracles.
\switchcolumn*
\selectlanguage{latin}
Eódem die Commemorátio 
 sanctórum Mártyrum germanórum Evódii, Hermógenis et Callístæ; de quibus, in 
 Syracusána Sicíliæ urbe martyrium passis, ágitur étiam séptimo Kaléndas Maji.
\switchcolumn
\selectlanguage{english}
On the same day, the commemoration 
 of the holy martyrs Evodius and Hermogenes, brothers, and Callista, their 
 sister. Mention is made of them that they died on the 25th of April in 
 the city of Syracuse in Italy.
\switchcolumn*
\selectlanguage{latin}
Lugdúni, in Gállia, 
 Translátio sanctórum Justi, Epíscopi et Confessóris, ac Viatóris, qui ejus 
 fúerat miníster; quorum natális dies respectíve prídie Idus Octóbris et 
 duodécimo Kaléndas Novémbris recensétur.
\switchcolumn
\selectlanguage{english}
At Lyons in France, the translation 
 of St. Justus, bishop and confessor, and Viator, his servant, whose 
 birthdays occur on the 14th of October and the 21st of October.
\switchcolumn*
\selectlanguage{latin}
\end{paracol}


% ---- martyrology/mart09/mart0903.htm
\needspace{10\baselineskip}
\begin{paracol}{2}
\selectlanguage{latin}
\begin{center}{\color{gregoriocolor} Tértio Nonas Septémbris. 
 Luna\dots\ }\end{center}
\switchcolumn
\selectlanguage{english}
\begin{center}{\color{gregoriocolor} The   3\textsuperscript{rd} Day of 
 September. The\dots\ Day of the Moon.}\end{center}
\end{paracol}

\noindent\begin{tabularx}{\linewidth}{*{19}{>{\centering\arraybackslash}X}}
 \textcolor{gregoriocolor}{a} & \textcolor{gregoriocolor}{b} & \textcolor{gregoriocolor}{c} & \textcolor{gregoriocolor}{d} & \textcolor{gregoriocolor}{e} & \textcolor{gregoriocolor}{f} & \textcolor{gregoriocolor}{g} & \textcolor{gregoriocolor}{h} & \textcolor{gregoriocolor}{i} & \textcolor{gregoriocolor}{k} & \textcolor{gregoriocolor}{l} & \textcolor{gregoriocolor}{m} & \textcolor{gregoriocolor}{n} & \textcolor{gregoriocolor}{p} & \textcolor{gregoriocolor}{q} & \textcolor{gregoriocolor}{r} & \textcolor{gregoriocolor}{s} & \textcolor{gregoriocolor}{t} & \textcolor{gregoriocolor}{u} \\
 11 & 12 & 13 & 14 & 15 & 16 & 17 & 18 & 19 & 20 & 21 & 22 & 23 & 24 & 25 & 26 & 27 & 28 & 29 \\
\end{tabularx}
\vspace{0.5\baselineskip}
\noindent\begin{tabularx}{\linewidth}{*{12}{>{\centering\arraybackslash}X}}
 \textcolor{gregoriocolor}{A} & \textcolor{gregoriocolor}{B} & \textcolor{gregoriocolor}{C} & \textcolor{gregoriocolor}{D} & \textcolor{gregoriocolor}{E} & F & \textcolor{gregoriocolor}{F} & \textcolor{gregoriocolor}{G} & \textcolor{gregoriocolor}{H} & \textcolor{gregoriocolor}{M} & \textcolor{gregoriocolor}{N} & \textcolor{gregoriocolor}{P} \\
 30 & 1 & 2 & 3 & 4 & 5 & 5 & 6 & 7 & 8 & 9 & 10 \\
\end{tabularx}

\begin{paracol}{2}
\selectlanguage{latin}
\lettrine[lines=2]{S}{ancti} Pii Papæ Décimi, 
 cujus natális dies tertiodécimo Kaléndas Septémbris recensétur.
\switchcolumn
\selectlanguage{english}
\lettrine[lines=2]{P}{ope} St. Pius X, whose birthday is 
 mentioned on the 20th of August.
\switchcolumn*
\selectlanguage{latin}
Corínthi natális sanctæ 
 Phœbes, cujus méminit beátus Apóstolus Paulus ad Romános scribens.
\switchcolumn
\selectlanguage{english}
At Corinth the birthday of St. 
 Phoebe, mentioned by the blessed apostle Paul in his Epistle to the Romans.
\switchcolumn*
\selectlanguage{latin}
Cápuæ sanctórum 
 Mártyrum Aristǽi Epíscopi, et Antoníni púeri.
\switchcolumn
\selectlanguage{english}
At Capua, the holy martyrs Aristaeus, 
 bishop, and Antoninus, a young boy.
\switchcolumn*
\selectlanguage{latin}
Eódem die natális 
 sanctórum Mártyrum Aigúlfi, Abbátis Lirinénsis, et Sociórum ipsíus 
 Monachórum, qui, linguis præcísis oculísque effóssis, gládio obtruncáti sunt.
\switchcolumn
\selectlanguage{english}
Also, the birthday of the holy 
 martyrs Aigulphus, abbot of Lerins, and the monks, his companions, who, 
 after their tongues were cut off and their eyes plucked out, were killed 
 with the sword.
\switchcolumn*
\selectlanguage{latin}
Item sanctórum Mártyrum 
 Zenónis et Charitónis; quorum alter in lebétem liquáti plumbi conjéctus est, 
 alter in ignis fornácem immíssus.
\switchcolumn
\selectlanguage{english}
Also, the holy martyrs Zeno and 
 Chariton. The one was cast into a cauldron of melted lead, the other 
 into a burning furnace.
\switchcolumn*
\selectlanguage{latin}
Córdubæ, in Hispánia, 
 sancti Sándali Mártyris.
\switchcolumn
\selectlanguage{english}
At Cordova in Spain, St. Sandal the 
 martyr.
\switchcolumn*
\selectlanguage{latin}
Aquiléjæ sanctárum 
 Vírginum et Mártyrum Euphémiæ, Dorótheæ, Theclæ et Erásmæ, quæ, sub Neróne 
 Imperatóre et Sebásto Præside, post multa supplícia, gládio cæsæ sunt, et a 
 sancto Hermágora sepúltæ.
\switchcolumn
\selectlanguage{english}
At Aquileia, the holy virgins and 
 martyrs Euphemia, Dorothy, Thecla, and Erasma. Under Nero, after 
 enduring many torments, they were slain with the sword and buried by St. 
 Hermagoras.
\switchcolumn*
\selectlanguage{latin}
Nicomedíæ pássio sanctæ 
 Basilíssæ, Vírginis et Mártyris; quæ, annos novem nata, cum in persecutióne 
 Diocletiáni Imperatóris, sub Alexándro Præside, vérbera, ignes ac béstias 
 divína virtúte superásset, ipsum Præsidem ad Christi fidem convértit, ac 
 tandem extra urbem in oratióne spíritum Deo réddidit.
\switchcolumn
\selectlanguage{english}
At Nicomedia, the passion of St. 
 Basilissa, virgin and martyr, in the persecution of Diocletian, under the 
 governor Alexander. At the age of nine years, after having, through 
 the power of God, overcome scourging, fire, and the beasts—by which she 
 converted the governor to the faith of Christ—she at length gave up her soul 
 to God while at prayer outside the city.
\switchcolumn*
\selectlanguage{latin}
Tulli, in Gállia, 
 sancti Mansuéti, Epíscopi et Confessóris.
\switchcolumn
\selectlanguage{english}
At Toul in France, St. Mansuetus, 
 bishop and confessor.
\switchcolumn*
\selectlanguage{latin}
Medioláni deposítio 
 sancti Auxáni Epíscopi.
\switchcolumn
\selectlanguage{english}
At Milan, the death of St. Auxanus, 
 bishop.
\switchcolumn*
\selectlanguage{latin}
Eódem die sancti 
 Simeónis Stylítæ junióris.
\switchcolumn
\selectlanguage{english}
The same day, St. Simon Stylites the 
 Younger.
\switchcolumn*
\selectlanguage{latin}
Romæ Translátio sanctæ 
 Serápiæ, Vírginis et Mártyris; quæ passa est quarto Kaléndas Augústi.
\switchcolumn
\selectlanguage{english}
At Rome, the translation of St. 
 Serapia, virgin and martyr, who suffered on the 29th of July.
\switchcolumn*
\selectlanguage{latin}
Item Romæ Ordinátio 
 incomparábilis viri sancti Gregórii Magni in Summum Pontíficem; qui, onus 
 illud subíre coáctus, e sublimióri throno clarióribus sanctitátis rádiis in 
 Orbe refúlsit.
\switchcolumn
\selectlanguage{english}
Also at Rome, the raising to the 
 Sovereign Pontificate of St. Gregory the Great. This incomparable man, 
 being forced to take that burden upon himself, sent forth from the exalted 
 throne brighter rays of sanctity upon the world.
\switchcolumn*
\selectlanguage{latin}
\end{paracol}


% ---- martyrology/mart09/mart0904.htm
\needspace{10\baselineskip}
\begin{paracol}{2}
\selectlanguage{latin}
\begin{center}{\color{gregoriocolor} Prídie Nonas Septémbris. 
 Luna\dots\ }\end{center}
\switchcolumn
\selectlanguage{english}
\begin{center}{\color{gregoriocolor} The   4\textsuperscript{th} Day of 
 September. The\dots\ Day of the Moon.}\end{center}
\end{paracol}

\noindent\begin{tabularx}{\linewidth}{*{19}{>{\centering\arraybackslash}X}}
 \textcolor{gregoriocolor}{a} & \textcolor{gregoriocolor}{b} & \textcolor{gregoriocolor}{c} & \textcolor{gregoriocolor}{d} & \textcolor{gregoriocolor}{e} & \textcolor{gregoriocolor}{f} & \textcolor{gregoriocolor}{g} & \textcolor{gregoriocolor}{h} & \textcolor{gregoriocolor}{i} & \textcolor{gregoriocolor}{k} & \textcolor{gregoriocolor}{l} & \textcolor{gregoriocolor}{m} & \textcolor{gregoriocolor}{n} & \textcolor{gregoriocolor}{p} & \textcolor{gregoriocolor}{q} & \textcolor{gregoriocolor}{r} & \textcolor{gregoriocolor}{s} & \textcolor{gregoriocolor}{t} & \textcolor{gregoriocolor}{u} \\
 12 & 13 & 14 & 15 & 16 & 17 & 18 & 19 & 20 & 21 & 22 & 23 & 24 & 25 & 26 & 27 & 28 & 29 & 30 \\
\end{tabularx}
\vspace{0.5\baselineskip}
\noindent\begin{tabularx}{\linewidth}{*{12}{>{\centering\arraybackslash}X}}
 \textcolor{gregoriocolor}{A} & \textcolor{gregoriocolor}{B} & \textcolor{gregoriocolor}{C} & \textcolor{gregoriocolor}{D} & \textcolor{gregoriocolor}{E} & F & \textcolor{gregoriocolor}{F} & \textcolor{gregoriocolor}{G} & \textcolor{gregoriocolor}{H} & \textcolor{gregoriocolor}{M} & \textcolor{gregoriocolor}{N} & \textcolor{gregoriocolor}{P} \\
 1 & 2 & 3 & 4 & 5 & 6 & 6 & 7 & 8 & 9 & 10 & 11 \\
\end{tabularx}

\begin{paracol}{2}
\selectlanguage{latin}
\lettrine[lines=2]{I}{n} monte Nebo, terræ 
 Moab, sancti Móysis, legislatóris et Prophétæ.
\switchcolumn
\selectlanguage{english}
\lettrine[lines=2]{O}{n} Mount Nebo, in the land of Moab, 
 the holy lawgiver and prophet Moses.
\switchcolumn*
\selectlanguage{latin}
Neápoli, in Campánia, natális sanctæ Cándidæ, quæ sancto Petro Apóstolo, ad eam urbem veniénti, 
 prima occúrrit, atque, ab eo baptizáta, póstea sancto fine quiévit.
\switchcolumn
\selectlanguage{english}
At Naples in Campania, the birthday 
 of St. Candida, who was the first to meet St. Peter when he came to that 
 city, and being baptized by him afterwards ended her holy life in peace.
\switchcolumn*
\selectlanguage{latin}
Tréviris sancti 
 Marcélli, Epíscopi et Mártyris.
\switchcolumn
\selectlanguage{english}
At Treves, St. Marcellus, bishop and 
 martyr.
\switchcolumn*
\selectlanguage{latin}
Ancyræ, in Galátia, natális sanctórum trium puerórum Mártyrum, id est Rufíni, Silváni et 
 Vitálici.
\switchcolumn
\selectlanguage{english}
At Ancyra in Galatia, the birthday 
 of three saintly boys, Rufinus, Silvanus, and Vitalicus, martyrs.
\switchcolumn*
\selectlanguage{latin}
Eódem die sanctórum 
 Mártyrum Magni, Casti et Máximi.
\switchcolumn
\selectlanguage{english}
On the same day, the holy martyrs 
 Magnus, Castus and Maximus.
\switchcolumn*
\selectlanguage{latin}
Cabillóne, in Gálliis, 
 sancti Marcélli Mártyris, qui, sub Antoníno Imperatóre, cum a Præside Prisco 
 ad profánum convívium fuísset invitátus, et, hujúsmodi épulas éxsecrans, 
 omnes qui áderant, cur idólis deservírent, líbera increpatióne corríperet, 
 ab eódem Præside, inaudíto crudelitátis génere, cíngulo tenus defóssus est 
 in terra; sicque, cum in Dei láudibus tríduo perseverásset, incontaminátum 
 spíritum réddidit.
\switchcolumn
\selectlanguage{english}
At Chalons in France, under Emperor 
 Antoninus, St. Marcellus, martyr. Being invited to a profane banquet 
 by the governor Priscus, he scorned to partake of the meats that were 
 served, and reproved with great freedom all persons present for worshipping 
 idols. For this, with unheard-of cruelty, the same governor had him 
 buried alive up to the waist. After persevering for three days in 
 praising God, he yielded up his undefiled spirit.
\switchcolumn*
\selectlanguage{latin}
Eódem die sanctórum 
 Thamélis, ántea idolórum sacerdótis, et Sociórum Mártyrum, sub Hadriáno 
 Imperatóre.
\switchcolumn
\selectlanguage{english}
On the same day, St. Thamel, 
 previously a pagan priest, and his companions, martyrs under Emperor 
 Hadrian.
\switchcolumn*
\selectlanguage{latin}
Item sanctórum Mártyrum 
 Theódori, Océani, Ammiáni et Juliáni; qui sub Maximiáno Imperatóre, 
 disséctis pédibus in ignem conjécti, martyrium consummárunt.
\switchcolumn
\selectlanguage{english}
Also, the holy martyrs Theodore, 
 Oceanus, Ammian, and Julian, who had their feet cut off, and completed their 
 martyrdom by being thrown into the fire, in the time of Emperor Maximian.
\switchcolumn*
\selectlanguage{latin}
Romæ sancti Bonifátii 
 Primi, Papæ et Confessóris.
\switchcolumn
\selectlanguage{english}
At Rome, St. Boniface I, pope and 
 confessor.
\switchcolumn*
\selectlanguage{latin}
Arímini sancti Maríni 
 Diáconi.
\switchcolumn
\selectlanguage{english}
At Rimini, St. Marinus, deacon.
\switchcolumn*
\selectlanguage{latin}
Panórmi natális sanctæ 
 Rosáliæ, Vírginis Panormitánæ, ex régio Cároli Magni sánguine ortæ; quæ, pro 
 Christi amóre, patérnum principátum aulámque profúgit, et, in móntibus ac 
 spelúncis solitária, cæléstem vitam duxit.
\switchcolumn
\selectlanguage{english}
At Palermo, the birthday of St. 
 Rosalia, virgin, a native of that city, born of the royal blood of 
 Charlemagne. For the love of Christ, she forsook the princely court of 
 her father, and led a saintly life alone in mountains and caverns.
\switchcolumn*
\selectlanguage{latin}
Vitérbii Translátio 
 beátæ Rosæ Vírginis, ex tértio Ordine sancti Francísci, témpore Alexándri 
 Papæ Quarti.
\switchcolumn
\selectlanguage{english}
At Viterbo, the translation of St. 
 Rose the Virgin, of the Third Order of St. Francis, during the pontificate 
 of Pope Alexander IV.
\switchcolumn*
\selectlanguage{latin}
\end{paracol}


% ---- martyrology/mart09/mart0905.htm
\needspace{10\baselineskip}
\begin{paracol}{2}
\selectlanguage{latin}
\begin{center}{\color{gregoriocolor} Nonis Septémbris. 
 Luna\dots\ }\end{center}
\switchcolumn
\selectlanguage{english}
\begin{center}{\color{gregoriocolor} The   5\textsuperscript{th} Day of 
 September. The\dots\ Day of the Moon.}\end{center}
\end{paracol}

\noindent\begin{tabularx}{\linewidth}{*{19}{>{\centering\arraybackslash}X}}
 \textcolor{gregoriocolor}{a} & \textcolor{gregoriocolor}{b} & \textcolor{gregoriocolor}{c} & \textcolor{gregoriocolor}{d} & \textcolor{gregoriocolor}{e} & \textcolor{gregoriocolor}{f} & \textcolor{gregoriocolor}{g} & \textcolor{gregoriocolor}{h} & \textcolor{gregoriocolor}{i} & \textcolor{gregoriocolor}{k} & \textcolor{gregoriocolor}{l} & \textcolor{gregoriocolor}{m} & \textcolor{gregoriocolor}{n} & \textcolor{gregoriocolor}{p} & \textcolor{gregoriocolor}{q} & \textcolor{gregoriocolor}{r} & \textcolor{gregoriocolor}{s} & \textcolor{gregoriocolor}{t} & \textcolor{gregoriocolor}{u} \\
 13 & 14 & 15 & 16 & 17 & 18 & 19 & 20 & 21 & 22 & 23 & 24 & 25 & 26 & 27 & 28 & 29 & 30 & 1 \\
\end{tabularx}
\vspace{0.5\baselineskip}
\noindent\begin{tabularx}{\linewidth}{*{12}{>{\centering\arraybackslash}X}}
 \textcolor{gregoriocolor}{A} & \textcolor{gregoriocolor}{B} & \textcolor{gregoriocolor}{C} & \textcolor{gregoriocolor}{D} & \textcolor{gregoriocolor}{E} & F & \textcolor{gregoriocolor}{F} & \textcolor{gregoriocolor}{G} & \textcolor{gregoriocolor}{H} & \textcolor{gregoriocolor}{M} & \textcolor{gregoriocolor}{N} & \textcolor{gregoriocolor}{P} \\
 2 & 3 & 4 & 5 & 6 & 7 & 7 & 8 & 9 & 10 & 11 & 12 \\
\end{tabularx}

\begin{paracol}{2}
\selectlanguage{latin}
\lettrine[lines=2]{S}{ancti} Lauréntii 
 Justiniáni, primi Patriárchæ Venetiárum et Confessóris, qui pontificálem 
 Cáthedram hac die invítus ascéndit, et sexto Idus Januárii obdormívit in 
 Dómino.
\switchcolumn
\selectlanguage{english}
\lettrine[lines=2]{S}{aint} Lawrence Justinian, first 
 patriarch of Venice and confessor, who on this day unwillingly ascended the 
 episcopal throne. His birthday is the 8th of January.
\switchcolumn*
\selectlanguage{latin}
Romæ, in suburbáno, 
 beáti Victoríni, Epíscopi et Mártyris; qui, sanctitáte et miráculis clarus, 
 sacerdótium Amiternínæ urbis, totíus pópuli electióne, adéptus est. 
 Póstmodum, sub Nerva Trajáno, apud Cutílias, ubi puténtes et sulphúreæ 
 emánant aquæ, cum áliis Dei servis, relegátus, ab Aureliáno Júdice jussus 
 est suspéndi cápite deórsum; idque cum per tríduum pro nómine Christi passus 
 fuísset, tandem, glorióse coronátus, victor migrávit ad Dóminum. Ejus 
 corpus rapuérunt Christiáni, et honorífica sepultúra Amitérni, in Vestínis, 
 condidérunt.
\switchcolumn
\selectlanguage{english}
In the suburbs of Rome, blessed 
 Victorinus, bishop and martyr, in the time of Nerva Trajan. Being 
 renowned for sanctity and miracles, he was elected bishop of Amiterno by the 
 whole populace, but afterwards he was banished, with other servants of God, 
 to Contigliano, where fetid sulphurous waters spring forth, and was 
 suspended with his head downward by order of the judge Aurelian. 
 Having for the name of Christ endured this torment for three days, he was 
 gloriously crowned and went victoriously to our Lord. His body was 
 taken away by the Christians and buried with due honours at Amiterno.
\switchcolumn*
\selectlanguage{latin}
Constantinópoli 
 sanctórum Mártyrum Urbáni, Theodóri, Menedémi, et Sociórum septuagínta 
 septem ex órdine ecclesiástico; qui a Valénte Imperatóre, pro fide cathólica, 
 in navígio impósiti, jussi sunt in mari combúri.
\switchcolumn
\selectlanguage{english}
At Constantinople, the holy martyrs 
 Urbanus, Theodore, Menedemus, and their companions of ecclesiastical rank, 
 seventy-seven in number, who were put in a ship by the command of Emperor 
 Valens, and burned on the sea for the Catholic faith.
\switchcolumn*
\selectlanguage{latin}
In Portu Románo pássio 
 sancti Herculáni mílitis, qui, sub Gallo Imperatóre, ob Christi fidem cæsus 
 flagris et cápite obtruncátus est.
\switchcolumn
\selectlanguage{english}
At Porto, the birthday of St. 
 Herculanus, martyr, who was scourged and beheaded in the reign of Emperor 
 Gallus because of the Christian faith.
\switchcolumn*
\selectlanguage{latin}
Cápuæ sanctórum 
 Mártyrum Quínctii, Arcóntii et Donáti.
\switchcolumn
\selectlanguage{english}
At Capua, the holy martyrs Quinctus, 
 Arcontius, and Donatus.
\switchcolumn*
\selectlanguage{latin}
Eódem die sancti Rómuli, 
 qui, cum esset Trajáni aulæ præféctus ac sævítiam Imperatóris in Christiános 
 detestarétur, cæsus est virgis, et cápite truncátus.
\switchcolumn
\selectlanguage{english}
On the same day, St. Romulus, 
 prefect of Trajan's court. For reproving the cruelty of the emperor 
 towards Christians, he was scourged with rods and beheaded.
\switchcolumn*
\selectlanguage{latin}
Melitínæ, in Arménia, pássio sanctórum mílitum Eudóxii, Zenónis, Macárii, et Sociórum mille centum 
 et quátuor; qui, cum abjecíssent milítiæ cíngulum, in persecutióne 
 Diocletiáni, pro Christi confessióne necáti sunt.
\switchcolumn
\selectlanguage{english}
At Melitine in Armenia, during the 
 persecution of Diocletian, the martyrdom of the holy soldiers Eudoxius, 
 Zeno, Macarius, and their companions to the number of eleven hundred and 
 four, who threw away their military belts and were put to death for the 
 confession of Christ.
\switchcolumn*
\selectlanguage{latin}
In pago Tarvanénsi, 
 monastério Sithinénsi, in Gállia, sancti Bertíni Abbátis.
\switchcolumn
\selectlanguage{english}
In the neighbourhood of Terouanne, 
 in the monastery of Sithiu, in France, St. Bertinus, abbot.
\switchcolumn*
\selectlanguage{latin}
Toléti, in Hispánia, 
 sanctæ Obdúliæ Vírginis.
\switchcolumn
\selectlanguage{english}
At Toledo in Spain, St. Obdulia, 
 virgin.
\switchcolumn*
\selectlanguage{latin}
\end{paracol}


% ---- martyrology/mart09/mart0906.htm
\needspace{10\baselineskip}
\begin{paracol}{2}
\selectlanguage{latin}
\begin{center}{\color{gregoriocolor} Octávo Idus Septémbris. 
 Luna\dots\ }\end{center}
\switchcolumn
\selectlanguage{english}
\begin{center}{\color{gregoriocolor} The   6\textsuperscript{th} Day of 
 September. The\dots\ Day of the Moon.}\end{center}
\end{paracol}

\noindent\begin{tabularx}{\linewidth}{*{19}{>{\centering\arraybackslash}X}}
 \textcolor{gregoriocolor}{a} & \textcolor{gregoriocolor}{b} & \textcolor{gregoriocolor}{c} & \textcolor{gregoriocolor}{d} & \textcolor{gregoriocolor}{e} & \textcolor{gregoriocolor}{f} & \textcolor{gregoriocolor}{g} & \textcolor{gregoriocolor}{h} & \textcolor{gregoriocolor}{i} & \textcolor{gregoriocolor}{k} & \textcolor{gregoriocolor}{l} & \textcolor{gregoriocolor}{m} & \textcolor{gregoriocolor}{n} & \textcolor{gregoriocolor}{p} & \textcolor{gregoriocolor}{q} & \textcolor{gregoriocolor}{r} & \textcolor{gregoriocolor}{s} & \textcolor{gregoriocolor}{t} & \textcolor{gregoriocolor}{u} \\
 14 & 15 & 16 & 17 & 18 & 19 & 20 & 21 & 22 & 23 & 24 & 25 & 26 & 27 & 28 & 29 & 30 & 1 & 2 \\
\end{tabularx}
\vspace{0.5\baselineskip}
\noindent\begin{tabularx}{\linewidth}{*{12}{>{\centering\arraybackslash}X}}
 \textcolor{gregoriocolor}{A} & \textcolor{gregoriocolor}{B} & \textcolor{gregoriocolor}{C} & \textcolor{gregoriocolor}{D} & \textcolor{gregoriocolor}{E} & F & \textcolor{gregoriocolor}{F} & \textcolor{gregoriocolor}{G} & \textcolor{gregoriocolor}{H} & \textcolor{gregoriocolor}{M} & \textcolor{gregoriocolor}{N} & \textcolor{gregoriocolor}{P} \\
 3 & 4 & 5 & 6 & 7 & 8 & 8 & 9 & 10 & 11 & 12 & 13 \\
\end{tabularx}

\begin{paracol}{2}
\selectlanguage{latin}
\lettrine[lines=2]{I}{n} Palæstína sancti 
 Zacharíæ Prophétæ, qui, de Chaldǽa senex in pátriam revérsus, ibíque 
 defúnctus, juxta Aggǽum Prophétam cónditus jacet.
\switchcolumn
\selectlanguage{english}
\lettrine[lines=2]{I}{n} Palestine, the prophet Zachary, 
 who returned in his old age from Chaldea to his own country, and lies buried 
 near the prophet Aggeus.
\switchcolumn*
\selectlanguage{latin}
In Hellespónto sancti 
 Onesíphori, Apostolórum discípuli, cujus méminit sanctus Paulus ad Timótheum 
 scribens. Ipse autem Onesíphorus ibídem, una cum sancto Porphyrio, 
 jussu Hadriáni Procónsulis ácriter verberátus et a ferócibus raptátus equis, 
 spíritum Deo réddidit.
\switchcolumn
\selectlanguage{english}
In the Hellespont, St. Onesiphorus, 
 disciple of the apostles, of whom St. Paul speaks in his Letter to Timothy. 
 He was severely scourged with St. Porphyry, by order of the proconsul 
 Adrian, and being dragged by wild horses, gave up his soul unto God.
\switchcolumn*
\selectlanguage{latin}
In Africa sanctórum 
 Episcopórum Donatiáni, Præsidii, Mansuéti, Germáni et Fúsculi; qui, in 
 persecutióne Wandálica, jussu Ariáni Regis Hunneríci, pro assertióne 
 cathólicæ veritátis, fústibus diríssime cæsi et exsílio elimináti sunt. 
 Inter eos étiam erat Epíscopus, nómine Lætus, strénuus atque doctíssimus vir, 
 qui, post diutúrnos cárceris squalóres, incéndio concremátus est.
\switchcolumn
\selectlanguage{english}
In Africa, in the persecution of the 
 Vandals, the holy bishops Donatian, Praesidius, Mansuetus, Germanus, and 
 Fusculus, who were most cruelly scourged and sent into exile by order of the 
 Arian king Hunneric, because they proclaimed the Catholic truth. Among 
 them was one named Laetus, also a bishop, a courageous and very learned man, 
 who was burned alive after a long imprisonment in a loathsome dungeon.
\switchcolumn*
\selectlanguage{latin}
Alexandríæ pássio 
 sanctórum Mártyrum Fausti Presbyteri, Macárii et Sociórum decem; qui, sub 
 Décio Imperatóre et Valério Præside, pro Christi nómine, abscíssis 
 cervícibus, martyrium complevérunt.
\switchcolumn
\selectlanguage{english}
At Alexandria, in the time of 
 Emperor Decius and the governor Valerius, the holy martyrs Faustus, a 
 priest, Macarius, and ten companions, who received the martyr's crown by 
 being beheaded for the name of Christ.
\switchcolumn*
\selectlanguage{latin}
In Cappadócia sanctórum 
 Mártyrum Cóttidi Diáconi, Eugénii et Sociórum.
\switchcolumn
\selectlanguage{english}
In Cappadocia, the holy martyrs 
 Cottidus, deacon, Eugene, and their companions.
\switchcolumn*
\selectlanguage{latin}
Verónæ sancti Petrónii, 
 Epíscopi et Confessóris.
\switchcolumn
\selectlanguage{english}
At Verona, St. Petronius, bishop and 
 confessor.
\switchcolumn*
\selectlanguage{latin}
Romæ sancti Eleuthérii 
 Abbátis, qui Dei servus fuit, atque (ut sanctus Gregórius Papa scribit) 
 oratióne et lácrimis mórtuum suscitávit.
\switchcolumn
\selectlanguage{english}
At Rome, the holy abbot Eleutherius, 
 a servant of God, who, according to the testimony of Pope St. Gregory, 
 raised a dead man to life by his prayers and tears.
\switchcolumn*
\selectlanguage{latin}
\end{paracol}


% ---- martyrology/mart09/mart0907.htm
\needspace{10\baselineskip}
\begin{paracol}{2}
\selectlanguage{latin}
\begin{center}{\color{gregoriocolor} Séptimo Idus Septémbris. 
 Luna\dots\ }\end{center}
\switchcolumn
\selectlanguage{english}
\begin{center}{\color{gregoriocolor} The   7\textsuperscript{th} Day of 
 September. The\dots\ Day of the Moon.}\end{center}
\end{paracol}

\noindent\begin{tabularx}{\linewidth}{*{19}{>{\centering\arraybackslash}X}}
 \textcolor{gregoriocolor}{a} & \textcolor{gregoriocolor}{b} & \textcolor{gregoriocolor}{c} & \textcolor{gregoriocolor}{d} & \textcolor{gregoriocolor}{e} & \textcolor{gregoriocolor}{f} & \textcolor{gregoriocolor}{g} & \textcolor{gregoriocolor}{h} & \textcolor{gregoriocolor}{i} & \textcolor{gregoriocolor}{k} & \textcolor{gregoriocolor}{l} & \textcolor{gregoriocolor}{m} & \textcolor{gregoriocolor}{n} & \textcolor{gregoriocolor}{p} & \textcolor{gregoriocolor}{q} & \textcolor{gregoriocolor}{r} & \textcolor{gregoriocolor}{s} & \textcolor{gregoriocolor}{t} & \textcolor{gregoriocolor}{u} \\
 15 & 16 & 17 & 18 & 19 & 20 & 21 & 22 & 23 & 24 & 25 & 26 & 27 & 28 & 29 & 30 & 1 & 2 & 3 \\
\end{tabularx}
\vspace{0.5\baselineskip}
\noindent\begin{tabularx}{\linewidth}{*{12}{>{\centering\arraybackslash}X}}
 \textcolor{gregoriocolor}{A} & \textcolor{gregoriocolor}{B} & \textcolor{gregoriocolor}{C} & \textcolor{gregoriocolor}{D} & \textcolor{gregoriocolor}{E} & F & \textcolor{gregoriocolor}{F} & \textcolor{gregoriocolor}{G} & \textcolor{gregoriocolor}{H} & \textcolor{gregoriocolor}{M} & \textcolor{gregoriocolor}{N} & \textcolor{gregoriocolor}{P} \\
 4 & 5 & 6 & 7 & 8 & 9 & 9 & 10 & 11 & 12 & 13 & 14 \\
\end{tabularx}

\begin{paracol}{2}
\selectlanguage{latin}
\lettrine[lines=2]{T}{recis,} in Gállia, 
 sancti Nemórii Diáconi, et Sociórum Mártyrum, quos Attila, Rex Hunnórum, 
 interfécit.
\switchcolumn
\selectlanguage{english}
\lettrine[lines=2]{A}{t} Troyes, St. Nemorius, deacon, and 
 his companions, all martyrs, who were slain by Attila, king of the Huns.
\switchcolumn*
\selectlanguage{latin}
Nicomedíæ natális beáti 
 Joánnis Mártyris, qui, cum vidéret crudélia advérsus Christiános edícta in 
 foro pendére, hinc fídei ardóre accénsus, injécta manu, illa detráxit atque 
 discérpsit. Cumque hoc relátum esset Diocletiáno et Maximiáno Augústis, 
 in eádem urbe constitútis, ómnia suppliciórum génera in eum experíri 
 jussérunt; quæ vir nobilíssimus tanta vultus ac spíritus alacritáte pértulit, 
 ut ne tristis quidem pro his vidéri potúerit.
\switchcolumn
\selectlanguage{english}
At Nicomedia, the birthday of the 
 blessed martyr John, who upon seeing the cruel edicts against Christians, 
 posted in the public square, and being inflamed with an ardent faith, 
 reached out his hand, took them away and tore them up. This was 
 related to Emperors Diocletian and Maximian, then residing in the city, who 
 gave orders that he should be subjected to many kinds of torments. The 
 noble champion bore them with such cheerfulness of spirit as not to shew on 
 his countenance the least trace of pain or grief.
\switchcolumn*
\selectlanguage{latin}
Cæsaréæ, in Cappadócia, sancti Eupsychii Mártyris, qui, sub Hadriáno Imperatóre, accusátus quod 
 Christiánus esset, in cárcerem conjéctus est; et, paulo post inde emíssus, 
 patrimónium statim véndidit, et prétium partim paupéribus, partim 
 accusatóribus, tamquam benefactóribus, distríbuit. Sed, íterum 
 comprehénsus, atque, cum sacrificáre nollet idólis, sævíssime dilaniátus et 
 gládio confóssus, sub Saprítio Júdice martyrium consummávit.
\switchcolumn
\selectlanguage{english}
At Caesarea in Cappadocia, in the 
 time of Emperor Adrian, St. Eupsychius, martyr, who was accused of 
 professing Christianity and who was cast into prison. Having been 
 released shortly after, he immediately sold his inheritance, and distributed 
 the price of it partly to his accusers, whom he regarded as his benefactors. 
 But being again arrested, under the judge Sapritius, he was tortured, 
 pierced through with a sword, and thus completed his martyrdom.
\switchcolumn*
\selectlanguage{latin}
Pompejópoli, in Cilícia, sancti Sozóntis Mártyris, qui, sub Maximiáno Imperatóre, in ignem injéctus, 
 réddidit spíritum.
\switchcolumn
\selectlanguage{english}
At Pompeiopolis in Cilicia, in the 
 time of Emperor Maximian, St. Sozon, a martyr who was thrown into the fire 
 and yielded up his spirit.
\switchcolumn*
\selectlanguage{latin}
Aquiléjæ sancti 
 Anastásii Mártyris.
\switchcolumn
\selectlanguage{english}
At Aquileia, St. Anastasius, martyr.
\switchcolumn*
\selectlanguage{latin}
Apud Aléxiam véterem, 
 in território Augustodunénsi, sanctæ Regínæ, Vírginis et Mártyris; quæ, sub 
 Procónsule Olybrio, cárceris et equúlei ac lampadárum afflícta supplíciis, 
 demum, cápite damnáta, migrávit ad Sponsum.
\switchcolumn
\selectlanguage{english}
In the diocese of Autun, under the 
 proconsul Olybrius, St. Regina, virgin and martyr. After having 
 suffered imprisonment, the rack, and burning with torches, she was finally 
 condemned to capital punishment, and so went to her spouse.
\switchcolumn*
\selectlanguage{latin}
Aureliánis, in Gállia, 
 deposítio sancti Evótii Epíscopi, qui, primo Románæ Ecclésiæ Subdiáconus, 
 dehinc divíno múnere per colúmbam designátus est Póntifex præfátæ urbis.
\switchcolumn
\selectlanguage{english}
At Orleans in France, the departure 
 from this life of the holy bishop Evortius, who was first a subdeacon of the 
 Roman Church, and afterwards, through a divine favour, was designated by a 
 dove as bishop of that city.
\switchcolumn*
\selectlanguage{latin}
In Gállia sancti 
 Augustális, Epíscopi et Confessóris.
\switchcolumn
\selectlanguage{english}
In France, St. Augustalis, bishop 
 and confessor.
\switchcolumn*
\selectlanguage{latin}
Cápuæ sancti Pámphili 
 Epíscopi.
\switchcolumn
\selectlanguage{english}
At Capua, St. Pamphilus, bishop.
\switchcolumn*
\selectlanguage{latin}
In território Parisiénsi sancti Clodoáldi, Presbyteri et Confessóris.
\switchcolumn
\selectlanguage{english}
In the territory of Paris, St. 
 Cloud, priest and confessor.
\switchcolumn*
\selectlanguage{latin}
\end{paracol}


% ---- martyrology/mart09/mart0908.htm
\needspace{10\baselineskip}
\begin{paracol}{2}
\selectlanguage{latin}
\begin{center}{\color{gregoriocolor} Sexto Idus Septémbris. 
 Luna\dots\ }\end{center}
\switchcolumn
\selectlanguage{english}
\begin{center}{\color{gregoriocolor} The   8\textsuperscript{th} Day of 
 September. The\dots\ Day of the Moon.}\end{center}
\end{paracol}

\noindent\begin{tabularx}{\linewidth}{*{19}{>{\centering\arraybackslash}X}}
 \textcolor{gregoriocolor}{a} & \textcolor{gregoriocolor}{b} & \textcolor{gregoriocolor}{c} & \textcolor{gregoriocolor}{d} & \textcolor{gregoriocolor}{e} & \textcolor{gregoriocolor}{f} & \textcolor{gregoriocolor}{g} & \textcolor{gregoriocolor}{h} & \textcolor{gregoriocolor}{i} & \textcolor{gregoriocolor}{k} & \textcolor{gregoriocolor}{l} & \textcolor{gregoriocolor}{m} & \textcolor{gregoriocolor}{n} & \textcolor{gregoriocolor}{p} & \textcolor{gregoriocolor}{q} & \textcolor{gregoriocolor}{r} & \textcolor{gregoriocolor}{s} & \textcolor{gregoriocolor}{t} & \textcolor{gregoriocolor}{u} \\
 16 & 17 & 18 & 19 & 20 & 21 & 22 & 23 & 24 & 25 & 26 & 27 & 28 & 29 & 30 & 1 & 2 & 3 & 4 \\
\end{tabularx}
\vspace{0.5\baselineskip}
\noindent\begin{tabularx}{\linewidth}{*{12}{>{\centering\arraybackslash}X}}
 \textcolor{gregoriocolor}{A} & \textcolor{gregoriocolor}{B} & \textcolor{gregoriocolor}{C} & \textcolor{gregoriocolor}{D} & \textcolor{gregoriocolor}{E} & F & \textcolor{gregoriocolor}{F} & \textcolor{gregoriocolor}{G} & \textcolor{gregoriocolor}{H} & \textcolor{gregoriocolor}{M} & \textcolor{gregoriocolor}{N} & \textcolor{gregoriocolor}{P} \\
 5 & 6 & 7 & 8 & 9 & 10 & 10 & 11 & 12 & 13 & 14 & 15 \\
\end{tabularx}

\begin{paracol}{2}
\selectlanguage{latin}
\lettrine[lines=2]{N}{atívitas} beatíssimæ 
 semper Vírginis Genitrícis Dei Maríæ.
\switchcolumn
\selectlanguage{english}
\lettrine[lines=2]{T}{he} Nativity of the most Blessed and 
 ever Virgin Mary, Mother of God.
\switchcolumn*
\selectlanguage{latin}
Sancti Hadriáni 
 Mártyris; cujus dies natális quarto Nonas Mártii recensétur, sed festívitas 
 hac die, qua sacrum ejus corpus Romam translátum fuit, potíssime celebrátur.
\switchcolumn
\selectlanguage{english}
St. Hadrian, martyr, whose birthday 
 is on the 4th of March. His feast, however, is observed today, the day 
 on which his holy body was translated to Rome.
\switchcolumn*
\selectlanguage{latin}
Valéntiæ, in Hispánia 
 Tarraconénsi, natális sancti Thomæ a Villa Nova, ex Eremitárum sancti 
 Augustíni Ordine, Epíscopi et Confessóris, insígnis propter flagrántem in 
 páuperes caritátem; qui ab Alexándro Papa Séptimo in Sanctórum númerum 
 adscríptus est. Ipsíus autem festum décimo Kaléndas Octóbris 
 celebrátur.
\switchcolumn
\selectlanguage{english}
At Valencia in Spain, the birthday 
 of St. Thomas of Villanova, bishop and confessor, of the order of the 
 Hermits of St. Augustine, distinguished by his ardent love for the poor. 
 He was inscribed among the saints by Pope Alexander VII, and his festival is 
 observed on the 22nd of this month.
\switchcolumn*
\selectlanguage{latin}
Alexandríæ sanctórum 
 Mártyrum Ammónis, Theóphili, Neotérii et aliórum vigínti duórum.
\switchcolumn
\selectlanguage{english}
At Alexandria, the holy martyrs 
 Ammon, Theophilus, Neoterius, and twenty-two others.
\switchcolumn*
\selectlanguage{latin}
Antiochíæ sanctórum 
 Timóthei et Fausti Mártyrum.
\switchcolumn
\selectlanguage{english}
At Antioch, the Saints Timothy and 
 Faustus, martyrs.
\switchcolumn*
\selectlanguage{latin}
Gazæ, in Palæstína, 
 sanctórum Mártyrum fratrum Eusébii, Néstabi et Zenónis; qui, témpore Juliáni 
 Apóstatæ, irruénte in eos turba Gentílium, discérpti atque necáti sunt.
\switchcolumn
\selectlanguage{english}
At Gaza in Palestine, in the time of 
 Julian the Apostate, the holy martyrs Eusebius, Nestabus, and Zeno, 
 brothers, who were torn to pieces by a multitude of pagans that rushed upon 
 them.
\switchcolumn*
\selectlanguage{latin}
Ibídem sancti Néstoris 
 Mártyris, qui sub eódem Juliáno, ab iísdem Gentílibus furéntibus sævíssime 
 cruciátus, emísit spíritum.
\switchcolumn
\selectlanguage{english}
In the same place, and under the 
 same Julian, St. Nestor, martyr, who breathed his last after being most 
 cruelly tortured by the same furious heathen.
\switchcolumn*
\selectlanguage{latin}
Romæ sancti Sérgii 
 Primi, Papæ et Confessóris.
\switchcolumn
\selectlanguage{english}
At Rome, St. Sergius I, pope and 
 confessor.
\switchcolumn*
\selectlanguage{latin}
Frisíngæ sancti 
 Corbiniáni, qui fuit primus ejúsdem civitátis Epíscopus. Hic, a sancto 
 Gregório Secúndo Pontífice ordinátus et ad prædicándum Evangélium missus, 
 úberes fructus in Gállia et Germánia rétulit, ac demum, virtútibus et 
 miráculis clarus, in pace quiévit.
\switchcolumn
\selectlanguage{english}
At Freisingen, St. Corbinian, first 
 bishop of that city. Being consecrated by Pope Gregory II and sent to 
 preach the Gospel, he reaped abundant fruits in France and Germany, and 
 finally rested in peace, renowned for virtues and miracles.
\switchcolumn*
\selectlanguage{latin}
Carthágine nova, in 
 Ameríca meridionáli, sancti Petri Claver, Sacerdótis e Societáte Jesu et 
 Confessóris; qui mira sui abnegatióne et exímia caritáte Nigrítis in 
 servitútem abdúctis, annos ámplius quadragínta, óperam impéndens, tercénta 
 fere eórum míllia Christo sua ipse manu regenerávit; et a Leóne Décimo 
 tértio, Pontífice Máximo, in Sanctórum númerum relátus est, ad dein étiam 
 cæléstis Patrónus peculiáris sacrárum ad Nigrítas Missiónum constitútus et 
 declarátus.
\switchcolumn
\selectlanguage{english}
In New Carthage in South America, 
 St. Peter Claver, priest of the Society of Jesus and confessor. He 
 devoted more than forty years with wonderful mortification and exceeding 
 charity to the service of the Negroes who had been enslaved, and with his 
 own hand baptized in Christ almost three hundred thousand of them. 
 Pope Leo XIII added him to the list of the saints, and then declared him to 
 be the special heavenly patron of all missions for the Negroes.
\switchcolumn*
\selectlanguage{latin}
\end{paracol}


% ---- martyrology/mart09/mart0909.htm
\needspace{10\baselineskip}
\begin{paracol}{2}
\selectlanguage{latin}
\begin{center}{\color{gregoriocolor} Quinto Idus Septémbris. 
 Luna\dots\ }\end{center}
\switchcolumn
\selectlanguage{english}
\begin{center}{\color{gregoriocolor} The   9\textsuperscript{th} Day of 
 September. The\dots\ Day of the Moon.}\end{center}
\end{paracol}

\noindent\begin{tabularx}{\linewidth}{*{19}{>{\centering\arraybackslash}X}}
 \textcolor{gregoriocolor}{a} & \textcolor{gregoriocolor}{b} & \textcolor{gregoriocolor}{c} & \textcolor{gregoriocolor}{d} & \textcolor{gregoriocolor}{e} & \textcolor{gregoriocolor}{f} & \textcolor{gregoriocolor}{g} & \textcolor{gregoriocolor}{h} & \textcolor{gregoriocolor}{i} & \textcolor{gregoriocolor}{k} & \textcolor{gregoriocolor}{l} & \textcolor{gregoriocolor}{m} & \textcolor{gregoriocolor}{n} & \textcolor{gregoriocolor}{p} & \textcolor{gregoriocolor}{q} & \textcolor{gregoriocolor}{r} & \textcolor{gregoriocolor}{s} & \textcolor{gregoriocolor}{t} & \textcolor{gregoriocolor}{u} \\
 17 & 18 & 19 & 20 & 21 & 22 & 23 & 24 & 25 & 26 & 27 & 28 & 29 & 30 & 1 & 2 & 3 & 4 & 5 \\
\end{tabularx}
\vspace{0.5\baselineskip}
\noindent\begin{tabularx}{\linewidth}{*{12}{>{\centering\arraybackslash}X}}
 \textcolor{gregoriocolor}{A} & \textcolor{gregoriocolor}{B} & \textcolor{gregoriocolor}{C} & \textcolor{gregoriocolor}{D} & \textcolor{gregoriocolor}{E} & F & \textcolor{gregoriocolor}{F} & \textcolor{gregoriocolor}{G} & \textcolor{gregoriocolor}{H} & \textcolor{gregoriocolor}{M} & \textcolor{gregoriocolor}{N} & \textcolor{gregoriocolor}{P} \\
 6 & 7 & 8 & 9 & 10 & 11 & 11 & 12 & 13 & 14 & 15 & 16 \\
\end{tabularx}

\begin{paracol}{2}
\selectlanguage{latin}
\lettrine[lines=2]{N}{icomedíæ} pássio 
 sanctórum Mártyrum Doróthei et Gorgónii, qui, cum essent apud Diocletiánum 
 Augústum honóres amplíssimos consecúti, et persecutiónem, quam ille 
 Christiánis inferébat, detestaréntur, præsénte eo, jussi sunt primo appéndi, 
 et flagris toto córpore laniári; deínde, viscéribus pelle nudátis, acéto et 
 sale perfúndi, sicque assári in cratícula; atque ad últimum, láqueo necári. 
 Interjécto autem témpore, beáti Gorgónii corpus Romam delátus fuit, ac via 
 Latína pósitum, et inde ad Basílicam sancti Petri translátum.
\switchcolumn
\selectlanguage{english}
\lettrine[lines=2]{A}{t} Nicomedia, the holy martyrs 
 Dorothy and Gorgonius. The greatest honours had been conferred on them 
 by Emperor Diocletian, but as they detested the cruelty which he exercised 
 against the Christians, they were by his order hung up in his presence and 
 lacerated with whips. Then, having the skin torn off from their bodies 
 and vinegar and salt poured over them, they were burned on a gridiron, and 
 finally strangled. After some time the body of blessed Gorgonius was 
 brought to Rome and deposited on the Latin Way. From there it was 
 transferred to the basilica of St. Peter.
\switchcolumn*
\selectlanguage{latin}
In Sabínis, trigésimo 
 ab Urbe milliário, sanctórum Mártyrum Hyacínthi, Alexándri et Tibúrtii.
\switchcolumn
\selectlanguage{english}
Among the Sabines, thirty miles from 
 Rome, the holy martyrs Hyacinth, Alexander, and Tiburtius.
\switchcolumn*
\selectlanguage{latin}
Sebáste, in Arménia, sancti Severiáni, qui, cum Licínii Imperatóris esset miles, et Quadragínta 
 Mártyres in cárcere deténtos frequens visitáret, hinc, jussu Lysiæ Præsidis, 
 saxo ad pedes ligáto suspénsus est, ac, verbéribus cæsus et flagris laniátus, 
 in torméntis réddidit spíritum.
\switchcolumn
\selectlanguage{english}
At Sebaste in Armenia, St. Severian, 
 a soldier of Emperor Licinius. For frequently visiting the Forty 
 Martyrs in prison, he was suspended in the air with a stone tied to his feet 
 by order of the governor Lysias, and being scourged and torn with whips, 
 yielded up his soul in the midst of his torments.
\switchcolumn*
\selectlanguage{latin}
Eódem die pássio sancti 
 Stratónis, qui pro Christo, ad duas árbores ligátus atque discérptus, 
 martyrium consummávit.
\switchcolumn
\selectlanguage{english}
On the same day, St. Strato, who 
 ended his martyrdom for Christ by being tied to two trees and torn asunder.
\switchcolumn*
\selectlanguage{latin}
Item sanctórum Mártyrum 
 Rufíni et Rufiniáni fratrum.
\switchcolumn
\selectlanguage{english}
Also, the holy martyrs Rufinus and 
 Rufinian, brothers.
\switchcolumn*
\selectlanguage{latin}
In territorio 
 Tarvanénsi, in Gállia, sancti Audomári Epíscopi.
\switchcolumn
\selectlanguage{english}
In the territory of Terouanne, St. 
 Omer, bishop.
\switchcolumn*
\selectlanguage{latin}
In monastério Cluanénsi, 
 in Hibérnia, sancti Queráni, Presbyteri et Abbátis.
\switchcolumn
\selectlanguage{english}
In the monastery of Clonmacnoise in 
 Ireland, St. Kiaran, priest and abbot.
\switchcolumn*
\selectlanguage{latin}
\end{paracol}


% ---- martyrology/mart09/mart0910.htm
\needspace{10\baselineskip}
\begin{paracol}{2}
\selectlanguage{latin}
\begin{center}{\color{gregoriocolor} Quarto Idus Septémbris. 
 Luna\dots\ }\end{center}
\switchcolumn
\selectlanguage{english}
\begin{center}{\color{gregoriocolor} The   10\textsuperscript{th} Day of 
 September. The\dots\ Day of the Moon.}\end{center}
\end{paracol}

\noindent\begin{tabularx}{\linewidth}{*{19}{>{\centering\arraybackslash}X}}
 \textcolor{gregoriocolor}{a} & \textcolor{gregoriocolor}{b} & \textcolor{gregoriocolor}{c} & \textcolor{gregoriocolor}{d} & \textcolor{gregoriocolor}{e} & \textcolor{gregoriocolor}{f} & \textcolor{gregoriocolor}{g} & \textcolor{gregoriocolor}{h} & \textcolor{gregoriocolor}{i} & \textcolor{gregoriocolor}{k} & \textcolor{gregoriocolor}{l} & \textcolor{gregoriocolor}{m} & \textcolor{gregoriocolor}{n} & \textcolor{gregoriocolor}{p} & \textcolor{gregoriocolor}{q} & \textcolor{gregoriocolor}{r} & \textcolor{gregoriocolor}{s} & \textcolor{gregoriocolor}{t} & \textcolor{gregoriocolor}{u} \\
 18 & 19 & 20 & 21 & 22 & 23 & 24 & 25 & 26 & 27 & 28 & 29 & 30 & 1 & 2 & 3 & 4 & 5 & 6 \\
\end{tabularx}
\vspace{0.5\baselineskip}
\noindent\begin{tabularx}{\linewidth}{*{12}{>{\centering\arraybackslash}X}}
 \textcolor{gregoriocolor}{A} & \textcolor{gregoriocolor}{B} & \textcolor{gregoriocolor}{C} & \textcolor{gregoriocolor}{D} & \textcolor{gregoriocolor}{E} & F & \textcolor{gregoriocolor}{F} & \textcolor{gregoriocolor}{G} & \textcolor{gregoriocolor}{H} & \textcolor{gregoriocolor}{M} & \textcolor{gregoriocolor}{N} & \textcolor{gregoriocolor}{P} \\
 7 & 8 & 9 & 10 & 11 & 12 & 12 & 13 & 14 & 15 & 16 & 17 \\
\end{tabularx}

\begin{paracol}{2}
\selectlanguage{latin}
\lettrine[lines=2]{T}{olentíni,} in Picéno, deposítio sancti Nicolái Confessóris, ex Ordine Eremitárum sancti Augustíni.
\switchcolumn
\selectlanguage{english}
\lettrine[lines=2]{A}{t} Tolentino in Piceno, the 
 departure from this life of St. Nicholas, confessor, of the order of the 
 Hermits of St. Augustine.
\switchcolumn*
\selectlanguage{latin}
In Africa natális 
 sanctórum Episcopórum Nemesiáni, Felícis, Lúcii, altérius item Felícis, 
 Littéi, Polyáni, Victóris, Jadéris, Datívi et aliórum; qui, sub Valeriáno et 
 Galliéno, exsurgénte persecutiónis rábie, ad primam confessiónis Christi 
 constántiam gráviter fústibus cæsi sunt, deínde, compédibus vincti et ad 
 fodiénda metálla deputáti, gloriósæ confessiónis agónem consummárunt.
\switchcolumn
\selectlanguage{english}
In Africa, the birthday of the holy 
 bishops Nemesian, Felix, Lucius, another Felix, Litteus, Polyanus, Victor, 
 Jader, Dativus, and others. Because a violent persecution was breaking 
 out under Valerian and Gallienus, they were at their first courageous 
 confession of Christ beaten with rods, placed in irons, and sent to dig in 
 the metal mines where they completed their combat with a glorious 
 confession.
\switchcolumn*
\selectlanguage{latin}
Leódii, in Bélgio, 
 sancti Theodárdi, Epíscopi et Mártyris; qui ánimam suam pósuit pro óvibus 
 suis, et miraculórum signis post mortem illúxit.
\switchcolumn
\selectlanguage{english}
At Liege in Belgium, St. Theodard, 
 bishop and martyr, who laid down his life for his flock, and after his death 
 was renowned for the gift of miracles.
\switchcolumn*
\selectlanguage{latin}
Chalcédone sanctórum 
 Mártyrum Sósthenis et Victóris, qui, in persecutióne Diocletiáni, sub Asiæ 
 Procónsule Prisco, post víncula et béstias superátas, jussi sunt incéndi; at 
 illi, salutántes se ínvicem in ósculo sancto, in oratióne pósiti emisérunt 
 spíritum.
\switchcolumn
\selectlanguage{english}
At Chalcedon, in the persecution of 
 Diocletian, the holy martyrs Sosthenes and Victor. Under Priscus, 
 proconsul of Asia, they were loaded with fetters and exposed to the beasts, 
 after which they were condemned to be burned. But while they were 
 saluting each other with a holy kiss and praying, they expired.
\switchcolumn*
\selectlanguage{latin}
Item sanctórum Mártyrum 
 Apéllii, Lucæ et Cleméntis
\switchcolumn
\selectlanguage{english}
Also the holy martyrs Apellius, 
 Luke, and Clement.
\switchcolumn*
\selectlanguage{latin}
In Bithynia sanctárum Vírginum sorórum Menodóræ, Metrodóræ et Nymphodóræ; quæ sub Maximiáno Imperatóre et Frontóne Præside, ob intrépidam in Christi fide constántiam martyrio coronátæ, pervenérunt ad glóriam.
\switchcolumn
\selectlanguage{english}
In Bithynia, the holy virgins 
 Menodora, Metrodora, and Nymphodora, sisters. Under Emperor Maximian 
 and the governor Fronto, they were crowned with martyrdom, and went to 
 eternal glory.
\switchcolumn*
\selectlanguage{latin}
Compostéllæ sancti 
 Petri Epíscopi, qui multis virtútibus et miráculis cláruit.
\switchcolumn
\selectlanguage{english}
At Compostella, St. Peter, bishop, who was celebrated for his many virtues and miracles.
\switchcolumn*
\selectlanguage{latin}
In civitáte Albigénsi, in Gállia, sancti Sálvii, Epíscopi et Confessóris.
\switchcolumn
\selectlanguage{english}
In the city of Albi in France, St. Salvius, bishop and confessor.
\switchcolumn*
\selectlanguage{latin}
Nováriæ sancti Agápii 
 Epíscopi.
\switchcolumn
\selectlanguage{english}
At Novara, St. Agapius, bishop.
\switchcolumn*
\selectlanguage{latin}
Constantinópoli sanctæ 
 Pulchériæ Augústæ, Vírginis, religióne et pietáte insígnis.
\switchcolumn
\selectlanguage{english}
At Constantinople, St. Pulcheria, 
 empress and virgin, distinguished by her piety and zeal for religion.
\switchcolumn*
\selectlanguage{latin}
Neápoli, in Campánia, sanctæ Cándidæ junióris, miráculis claræ.
\switchcolumn
\selectlanguage{english}
At Naples in Campania, St. Candida 
 the Younger, famed for miracles.
\switchcolumn*
\selectlanguage{latin}
\end{paracol}


% ---- martyrology/mart09/mart0911.htm
\needspace{10\baselineskip}
\begin{paracol}{2}
\selectlanguage{latin}
\begin{center}{\color{gregoriocolor} Tértio Idus Septémbris. 
 Luna\dots\ }\end{center}
\switchcolumn
\selectlanguage{english}
\begin{center}{\color{gregoriocolor} The   11\textsuperscript{th} Day of 
 September. The\dots\ Day of the Moon.}\end{center}
\end{paracol}

\noindent\begin{tabularx}{\linewidth}{*{19}{>{\centering\arraybackslash}X}}
 \textcolor{gregoriocolor}{a} & \textcolor{gregoriocolor}{b} & \textcolor{gregoriocolor}{c} & \textcolor{gregoriocolor}{d} & \textcolor{gregoriocolor}{e} & \textcolor{gregoriocolor}{f} & \textcolor{gregoriocolor}{g} & \textcolor{gregoriocolor}{h} & \textcolor{gregoriocolor}{i} & \textcolor{gregoriocolor}{k} & \textcolor{gregoriocolor}{l} & \textcolor{gregoriocolor}{m} & \textcolor{gregoriocolor}{n} & \textcolor{gregoriocolor}{p} & \textcolor{gregoriocolor}{q} & \textcolor{gregoriocolor}{r} & \textcolor{gregoriocolor}{s} & \textcolor{gregoriocolor}{t} & \textcolor{gregoriocolor}{u} \\
 19 & 20 & 21 & 22 & 23 & 24 & 25 & 26 & 27 & 28 & 29 & 30 & 1 & 2 & 3 & 4 & 5 & 6 & 7 \\
\end{tabularx}
\vspace{0.5\baselineskip}
\noindent\begin{tabularx}{\linewidth}{*{12}{>{\centering\arraybackslash}X}}
 \textcolor{gregoriocolor}{A} & \textcolor{gregoriocolor}{B} & \textcolor{gregoriocolor}{C} & \textcolor{gregoriocolor}{D} & \textcolor{gregoriocolor}{E} & F & \textcolor{gregoriocolor}{F} & \textcolor{gregoriocolor}{G} & \textcolor{gregoriocolor}{H} & \textcolor{gregoriocolor}{M} & \textcolor{gregoriocolor}{N} & \textcolor{gregoriocolor}{P} \\
 8 & 9 & 10 & 11 & 12 & 13 & 13 & 14 & 15 & 16 & 17 & 18 \\
\end{tabularx}

\begin{paracol}{2}
\selectlanguage{latin}
\lettrine[lines=2]{R}{omæ,} via Salária 
 véteri, in cœmetério Basíllæ, natális sanctórum Mártyrum Proti et Hyacínthi 
 fratrum, eunuchórum beátæ Eugéniæ. Hi, sub Galliéno Imperatóre, 
 deprehénsi quod essent Christiáni, sacrificáre cogúntur; sed non 
 consentiéntes, primo duríssime verberáti sunt, ac tandem páriter decolláti.
\switchcolumn
\selectlanguage{english}
\lettrine[lines=2]{A}{t} Rome, on the old Salarian Way in 
 the cemetery of Basilla, the birthday of the holy martyrs Protus and 
 Hyacinth, brothers, and eunuchs in the service of blessed Eugenia. 
 They were arrested in the time of Emperor Gallienus on the charge of being 
 Christians, and urged to offer sacrifice to the gods. Because they 
 refused, they were most severely scourged and finally beheaded.
\switchcolumn*
\selectlanguage{latin}
Legióne, in Hispánia, 
 sancti Vincéntii, Abbátis et Mártyris.
\switchcolumn
\selectlanguage{english}
At Leon in Spain, St. Vincent, abbot 
 and martyr.
\switchcolumn*
\selectlanguage{latin}
Laodicéæ, in Syria, 
 pássio sanctórum Diodóri, Diomédis et Dídymi.
\switchcolumn
\selectlanguage{english}
At Laodicea in Syria, the martyrdom 
 of Saints Diodorus, Diomedes, and Didymus.
\switchcolumn*
\selectlanguage{latin}
In Ægypto sancti 
 Paphnútii Epíscopi, qui unus fuit ex iis Confessóribus, qui, sub Galério 
 Maximiáno Imperatóre, dextro óculo effósso et sinístro póplite excíso, ad 
 metálla damnáti fuérunt; deínde, sub Constantíno Magno, advérsus Ariános, 
 pro fide cathólica, strénue decertávit; et demum, multis corónis auctus, in 
 pace quiévit.
\switchcolumn
\selectlanguage{english}
In Egypt, the holy bishop Paphnutius, 
 one of those confessors who, under Emperor Galerius Maximinus, having the 
 right eye plucked out and the joint of the left knee cut, were condemned to 
 work in the metal mines. Afterwards, under Constantine the Great, he 
 courageously strove for the Catholic faith against the Arians, and at 
 length, adorned with many crowns, rested in peace.
\switchcolumn*
\selectlanguage{latin}
Lugdúni, in Gállia, 
 deposítio sancti Patiéntis Epíscopi.
\switchcolumn
\selectlanguage{english}
At Lyons in France, the death of St. 
 Patiens, bishop.
\switchcolumn*
\selectlanguage{latin}
Vercéllis sancti 
 Æmiliáni Epíscopi.
\switchcolumn
\selectlanguage{english}
At Vercelli, St. Aemilian, bishop.
\switchcolumn*
\selectlanguage{latin}
Alexandríæ sanctæ 
 Theodóræ, quæ, cum incáute deliquísset, inde, facti pænitens, mirábili 
 abstinéntia et patiéntia in hábitu sancto perseverávit incógnita usque ad 
 mortem.
\switchcolumn
\selectlanguage{english}
At Alexandria, St. Theodora, who 
 having committed a fault through imprudence and repenting of it, remained 
 unknown in a religious habit, and persevered until her death in practices of 
 extraordinary abstinence and patience.
\switchcolumn*
\selectlanguage{latin}
\end{paracol}


% ---- martyrology/mart09/mart0912.htm
\needspace{10\baselineskip}
\begin{paracol}{2}
\selectlanguage{latin}
\begin{center}{\color{gregoriocolor} Prídie Idus Septémbris. 
 Luna\dots\ }\end{center}
\switchcolumn
\selectlanguage{english}
\begin{center}{\color{gregoriocolor} The   12\textsuperscript{th} Day of 
 September. The\dots\ Day of the Moon.}\end{center}
\end{paracol}

\noindent\begin{tabularx}{\linewidth}{*{19}{>{\centering\arraybackslash}X}}
 \textcolor{gregoriocolor}{a} & \textcolor{gregoriocolor}{b} & \textcolor{gregoriocolor}{c} & \textcolor{gregoriocolor}{d} & \textcolor{gregoriocolor}{e} & \textcolor{gregoriocolor}{f} & \textcolor{gregoriocolor}{g} & \textcolor{gregoriocolor}{h} & \textcolor{gregoriocolor}{i} & \textcolor{gregoriocolor}{k} & \textcolor{gregoriocolor}{l} & \textcolor{gregoriocolor}{m} & \textcolor{gregoriocolor}{n} & \textcolor{gregoriocolor}{p} & \textcolor{gregoriocolor}{q} & \textcolor{gregoriocolor}{r} & \textcolor{gregoriocolor}{s} & \textcolor{gregoriocolor}{t} & \textcolor{gregoriocolor}{u} \\
 20 & 21 & 22 & 23 & 24 & 25 & 26 & 27 & 28 & 29 & 30 & 1 & 2 & 3 & 4 & 5 & 6 & 7 & 8 \\
\end{tabularx}
\vspace{0.5\baselineskip}
\noindent\begin{tabularx}{\linewidth}{*{12}{>{\centering\arraybackslash}X}}
 \textcolor{gregoriocolor}{A} & \textcolor{gregoriocolor}{B} & \textcolor{gregoriocolor}{C} & \textcolor{gregoriocolor}{D} & \textcolor{gregoriocolor}{E} & F & \textcolor{gregoriocolor}{F} & \textcolor{gregoriocolor}{G} & \textcolor{gregoriocolor}{H} & \textcolor{gregoriocolor}{M} & \textcolor{gregoriocolor}{N} & \textcolor{gregoriocolor}{P} \\
 9 & 10 & 11 & 12 & 13 & 14 & 14 & 15 & 16 & 17 & 18 & 19 \\
\end{tabularx}

\begin{paracol}{2}
\selectlanguage{latin}
\lettrine[lines=2]{F}{estum} sanctíssimi 
 Nóminis beátæ Maríæ, quod Innocéntius Undécimus, Póntifex Máximus, ob 
 insígnem victóriam de Turcis, ipsíus Vírginis præsídio, Vindobónæ in Austria 
 reportátam, celebrári jussit.
\switchcolumn
\selectlanguage{english}
\lettrine[lines=2]{T}{he} feast of the most holy Name of 
 the Blessed Virgin Mary, celebrated by order of the Sovereign Pontiff, 
 Innocent XI, on account of the signal victory gained over the Turks at 
 Vienna in Austria through her protection.
\switchcolumn*
\selectlanguage{latin}
In Bithynia sancti 
 Autónomi, Epíscopi et Mártyris; qui ex Itália, Diocletiáni Imperatóris 
 persecutiónem declínans, illuc proféctus, ibi, cum plúrimos convertísset ad 
 fidem, a furéntibus Gentílibus, dum sacra Mystéria perágeret, ad altáre 
 mactátus est, et hóstia Christi efféctus.
\switchcolumn
\selectlanguage{english}
In Bithynia, St. Autonomus, bishop 
 and martyr, who went to that country from Italy to avoid the persecution of 
 Diocletian. After he had converted many to the faith, he was killed at 
 the altar by the furious heathen while celebrating the sacred mysteries, and 
 thus he became a victim for Christ.
\switchcolumn*
\selectlanguage{latin}
Icónii, in Lycaónia, 
 sancti Curónoti Epíscopi, qui sub Perénnio Præside, cápite truncátus, 
 martyrii palmam accépit.
\switchcolumn
\selectlanguage{english}
At Iconium in Lycaonia, the holy 
 bishop Curonotus, who received the crown of martyrdom by being beheaded 
 under the governor Perennius.
\switchcolumn*
\selectlanguage{latin}
Alexandríæ natális 
 sanctórum Mártyrum Hierónidis, Leóntii, Serapiónis, Selésii, Valeriáni et 
 Stratónis, qui, sub Maximíno Imperatóre, ob confessiónem nóminis Christi, in 
 mare sunt demérsi.
\switchcolumn
\selectlanguage{english}
At Alexandria, in the time of 
 Emperor Maximinus, the birthday of the holy martyrs Hieronides, Leontius, 
 Serapion, Selesius, Valerian, and Strato, who were drowned in the sea for 
 the confession of the name of Christ.
\switchcolumn*
\selectlanguage{latin}
Meri, in Phrygia, 
 pássio sanctórum Mártyrum Macedónii, Theodúli et Tatiáni, qui, sub Juliáno Apóstata, ab Almáchio Præside, post ália torménta, super crates férreas 
 ignítas pósiti, exsultántes martyrium complevérunt.
\switchcolumn
\selectlanguage{english}
At Merum in Phrygia, the holy 
 martyrs Macedonius, Theodulus, and Tatian, under Julian the Apostate. 
 After other torments, they joyfully completed their martyrdom by being laid 
 on burning gridirons by order of the governor Almachius.
\switchcolumn*
\selectlanguage{latin}
Apud Papíam sancti 
 Juvéntii Epíscopi, de quo ágitur sexto Idus Februárii. Ipse, a beáto 
 Hermágora, discípulo sancti Marci Evangelístæ, ad eam urbem, una cum sancto 
 Syro, cujus memória recólitur quinto Idus Decémbris, diréctus est; et ambo, 
 prædicántes illic Christi Evangélium et magnis virtútibus ac miráculis 
 coruscántes, étiam vicínas urbes divínis opéribus illustrárunt, sicque in 
 pontificáli honóre, glorióso fine, quievérunt in pace.
\switchcolumn
\selectlanguage{english}
At Pavia, St. Juventius, bishop, 
 mentioned on the 8th of February. The blessed Hermagoras, disciple of 
 the evangelist St. Mark, sent him to that city along with St. Cyrus, who is 
 mentioned on the 9th of December. They both preached the Gospel of 
 Christ there, and being renowned for great virtues and miracles, enlightened 
 the neighbouring cities by divine works. They closed their glorious 
 careers in peace, invested with the episcopal office.
\switchcolumn*
\selectlanguage{latin}
Lugdúni, in Gállia, 
 deposítio sancti Sacerdótis Epíscopi.
\switchcolumn
\selectlanguage{english}
At Lyons in France, the death of St. 
 Sacerdos, bishop.
\switchcolumn*
\selectlanguage{latin}
Verónæ sancti Silvíni 
 Epíscopi.
\switchcolumn
\selectlanguage{english}
At Verona, St. Silvinus, bishop.
\switchcolumn*
\selectlanguage{latin}
Anderláci, prope 
 Bruxéllas, in Brabántia, sancti Guidónis Confessóris.
\switchcolumn
\selectlanguage{english}
At Anderlecht, near Brussels in 
 Belgium, St. Guy, confessor.
\switchcolumn*
\selectlanguage{latin}
\end{paracol}


% ---- martyrology/mart09/mart0913.htm
\needspace{10\baselineskip}
\begin{paracol}{2}
\selectlanguage{latin}
\begin{center}{\color{gregoriocolor} Idibus Septémbris. 
 Luna\dots\ }\end{center}
\switchcolumn
\selectlanguage{english}
\begin{center}{\color{gregoriocolor} The   13\textsuperscript{th} Day of 
 September. The\dots\ Day of the Moon.}\end{center}
\end{paracol}

\noindent\begin{tabularx}{\linewidth}{*{19}{>{\centering\arraybackslash}X}}
 \textcolor{gregoriocolor}{a} & \textcolor{gregoriocolor}{b} & \textcolor{gregoriocolor}{c} & \textcolor{gregoriocolor}{d} & \textcolor{gregoriocolor}{e} & \textcolor{gregoriocolor}{f} & \textcolor{gregoriocolor}{g} & \textcolor{gregoriocolor}{h} & \textcolor{gregoriocolor}{i} & \textcolor{gregoriocolor}{k} & \textcolor{gregoriocolor}{l} & \textcolor{gregoriocolor}{m} & \textcolor{gregoriocolor}{n} & \textcolor{gregoriocolor}{p} & \textcolor{gregoriocolor}{q} & \textcolor{gregoriocolor}{r} & \textcolor{gregoriocolor}{s} & \textcolor{gregoriocolor}{t} & \textcolor{gregoriocolor}{u} \\
 21 & 22 & 23 & 24 & 25 & 26 & 27 & 28 & 29 & 30 & 1 & 2 & 3 & 4 & 5 & 6 & 7 & 8 & 9 \\
\end{tabularx}
\vspace{0.5\baselineskip}
\noindent\begin{tabularx}{\linewidth}{*{12}{>{\centering\arraybackslash}X}}
 \textcolor{gregoriocolor}{A} & \textcolor{gregoriocolor}{B} & \textcolor{gregoriocolor}{C} & \textcolor{gregoriocolor}{D} & \textcolor{gregoriocolor}{E} & F & \textcolor{gregoriocolor}{F} & \textcolor{gregoriocolor}{G} & \textcolor{gregoriocolor}{H} & \textcolor{gregoriocolor}{M} & \textcolor{gregoriocolor}{N} & \textcolor{gregoriocolor}{P} \\
 10 & 11 & 12 & 13 & 14 & 15 & 15 & 16 & 17 & 18 & 19 & 20 \\
\end{tabularx}

\begin{paracol}{2}
\selectlanguage{latin}
\lettrine[lines=2]{A}{lexandríæ} natális 
 beáti Philíppi, patris sanctæ Eugéniæ Vírginis. Hic, dignitátem 
 Præfectúræ Ægypti déserens, Baptísmatis grátiam assecútus est; quem, in 
 oratióne constitútum, jussit Teréntius Præféctus, ejus succéssor, gládio 
 jugulári.
\switchcolumn
\selectlanguage{english}
\lettrine[lines=2]{A}{t} Alexandria, the birthday of 
 blessed Philip, father of the virgin St. Eugenia. Resigning the 
 dignity of prefect of Egypt, he received the grace of baptism. His 
 successor, the prefect Terentius, had him pierced through the throat with a 
 sword while he was praying.
\switchcolumn*
\selectlanguage{latin}
Item sanctórum Mártyrum 
 Macróbii et Juliáni, qui sub Licínio passi sunt.
\switchcolumn
\selectlanguage{english}
Also, the holy martyrs Macrobius and 
 Julian, who suffered under Licinius.
\switchcolumn*
\selectlanguage{latin}
Eódem die sancti 
 Ligórii Mártyris, qui a Gentílibus, in erémo degens, ob Christi fidem 
 necátus est.
\switchcolumn
\selectlanguage{english}
On the same day, St. Ligorius, 
 martyr. While living in the desert, he was murdered by heathens for 
 the faith of Christ.
\switchcolumn*
\selectlanguage{latin}
Alexandríæ sancti 
 Eulógii Epíscopi, doctrína et sanctitáte célebris.
\switchcolumn
\selectlanguage{english}
At Alexandria, St. Eulogius, a 
 bishop celebrated for learning and sanctity.
\switchcolumn*
\selectlanguage{latin}
Andégavi, in Gállia, 
 sancti Maurílii Epíscopi, qui innúmeris miráculis cláruit.
\switchcolumn
\selectlanguage{english}
At Angers in France, St. Maurilius, 
 a bishop renowned for numberless miracles.
\switchcolumn*
\selectlanguage{latin}
Apud Sénonas sancti 
 Amáti, Epíscopi et Confessóris.
\switchcolumn
\selectlanguage{english}
At Sens, St. Amatus, bishop and 
 confessor.
\switchcolumn*
\selectlanguage{latin}
In monastério Romárico, 
 in Gállia, sancti Amáti, Presbyteri et Abbátis, abstinéntia et miraculórum 
 dono illústris.
\switchcolumn
\selectlanguage{english}
In the monastery of Remiremont in 
 France, St. Amatus, priest and abbot, illustrious for the virtue of 
 abstinence and the gift of miracles.
\switchcolumn*
\selectlanguage{latin}
Ipso die sancti Venérii 
 Confessóris, admirándæ sanctitátis viri, qui in ínsula Palmária vitam 
 eremíticam duxit.
\switchcolumn
\selectlanguage{english}
The same day, St. Venerius, 
 confessor, a man of admirable sanctity who led the life of a hermit on the 
 island of Palmaria.
\switchcolumn*
\selectlanguage{latin}
\end{paracol}


% ---- martyrology/mart09/mart0914.htm
\needspace{10\baselineskip}
\begin{paracol}{2}
\selectlanguage{latin}
\begin{center}{\color{gregoriocolor} Décimo octávo Kaléndas Octóbris. 
 Luna\dots\ }\end{center}
\switchcolumn
\selectlanguage{english}
\begin{center}{\color{gregoriocolor} The   14\textsuperscript{th} Day of 
 September. The\dots\ Day of the Moon.}\end{center}
\end{paracol}

\noindent\begin{tabularx}{\linewidth}{*{19}{>{\centering\arraybackslash}X}}
 \textcolor{gregoriocolor}{a} & \textcolor{gregoriocolor}{b} & \textcolor{gregoriocolor}{c} & \textcolor{gregoriocolor}{d} & \textcolor{gregoriocolor}{e} & \textcolor{gregoriocolor}{f} & \textcolor{gregoriocolor}{g} & \textcolor{gregoriocolor}{h} & \textcolor{gregoriocolor}{i} & \textcolor{gregoriocolor}{k} & \textcolor{gregoriocolor}{l} & \textcolor{gregoriocolor}{m} & \textcolor{gregoriocolor}{n} & \textcolor{gregoriocolor}{p} & \textcolor{gregoriocolor}{q} & \textcolor{gregoriocolor}{r} & \textcolor{gregoriocolor}{s} & \textcolor{gregoriocolor}{t} & \textcolor{gregoriocolor}{u} \\
 22 & 23 & 24 & 25 & 26 & 27 & 28 & 29 & 30 & 1 & 2 & 3 & 4 & 5 & 6 & 7 & 8 & 9 & 10 \\
\end{tabularx}
\vspace{0.5\baselineskip}
\noindent\begin{tabularx}{\linewidth}{*{12}{>{\centering\arraybackslash}X}}
 \textcolor{gregoriocolor}{A} & \textcolor{gregoriocolor}{B} & \textcolor{gregoriocolor}{C} & \textcolor{gregoriocolor}{D} & \textcolor{gregoriocolor}{E} & F & \textcolor{gregoriocolor}{F} & \textcolor{gregoriocolor}{G} & \textcolor{gregoriocolor}{H} & \textcolor{gregoriocolor}{M} & \textcolor{gregoriocolor}{N} & \textcolor{gregoriocolor}{P} \\
 11 & 12 & 13 & 14 & 15 & 16 & 16 & 17 & 18 & 19 & 20 & 21 \\
\end{tabularx}

\begin{paracol}{2}
\selectlanguage{latin}
\lettrine[lines=2]{E}{xaltátio} sanctæ Crucis, quando Heracl\'ius Imperátor, Chósroa Rege devícto, eam de Pérside Hierosólymam reportávit.
\switchcolumn
\selectlanguage{english}
\lettrine[lines=2]{T}{he} Exaltation of the Holy Cross, 
 when Emperor Heraclius, after defeating King Chosroes, brought it back to 
 Jerusalem from Persia.
\switchcolumn*
\selectlanguage{latin}
Romæ, via Appia, beáti 
 Cornélii, Papæ et Mártyris; qui, in persecutióne Décii, post exsílii 
 relegatiónem, jussus est plumbátis cædi, et sic, cum áliis vigínti et uno 
 promíscui sexus, decollári. Sed et Cæreális, miles cum Sallústia uxóre, 
 quos idem Cornélius in fide instrúxerat, eódem die sunt cápite plexi.
\switchcolumn
\selectlanguage{english}
At Rome, on the Appian Way, during 
 the persecution of Decius, blessed Cornelius, pope and martyr. After 
 being banished, he was scourged with leaded whips and then beheaded with 
 twenty-one others of both sexes. On the same day were condemned to 
 capital punishment Caerealis, a soldier, and his wife Sallustia, who had 
 been instructed in the faith by the same Cornelius.
\switchcolumn*
\selectlanguage{latin}
In Africa pássio sancti 
 Cypriáni, Epíscopi Carthaginénsis, sanctitáte et doctrína claríssimi; qui, 
 sub Valeriáno et Galliéno Princípibus, post durum exsílium, cápitis 
 detruncatióne martyrium consummávit, sexto milliário a Carthágine, juxta 
 mare. Eorúndem vero sanctórum Cornélii et Cypriáni memória sextodécimo 
 Kaléndas Octóbris festíve celebrátur.
\switchcolumn
\selectlanguage{english}
In Africa, in the time of Emperors 
 Valerian and Gallienus, St. Cyprian, bishop of Carthage, most renowned for 
 holiness and learning. It was near the seashore, six miles from the 
 city, that he completed his martyrdom by beheading, after enduring a most 
 painful exile. The feast of the Saints Cornelius and Cyprian is kept 
 on the 16th of this month.
\switchcolumn*
\selectlanguage{latin}
Apud Cománam, in Ponto, 
 natális sancti Joánnis, Epíscopi Constantinopolitáni, Confessóris et 
 Ecclésiæ Doctóris, propter áureum eloquéntiæ flumen cognoménto Chrysóstomi; 
 qui, ab inimicórum factióne in exsílium ejéctus, et, cum e sancti Innocéntii 
 Primi, Summi Pontíficis, decréto inde revocarétur, in itínere, a 
 custodiéntibus milítibus multa mala perpéssus, ánimam Deo réddidit. 
 Ejus autem festívitas sexto Kaléndas Februárii celebrátur, quo die sacrum 
 ipsíus corpus a Theodósio junióre Constantinópolim fuit translátum. 
 Hunc vero præclaríssimum divíni verbi præcónem Pius Papa Décimus cæléstem 
 Oratórum sacrórum Patrónum declarávit atque constítuit.
\switchcolumn
\selectlanguage{english}
At Comana in Pontus, the birthday of 
 St. John, bishop of Constantinople, confessor and doctor of the Church, 
 surnamed Chrysostom because of his golden eloquence. He was cast into exile by a faction of his enemies, but was recalled by a decree of Pope 
 Innocent I. However, he suffered many evils on the journey at the 
 hands of the soldiers who guarded him, and he rendered up his soul unto God. 
 His feast is kept on the 27th of January, on which day his holy body was 
 translated to Constantinople by Theodosius the Younger. Pope Pius X 
 declared and appointed this glorious preacher of the divine Word as heavenly 
 patron of those preaching of holy things.
\switchcolumn*
\selectlanguage{latin}
Tréviris sancti Matérni Epíscopi, qui fuit discípulus beáti Petri Apóstoli; ac Tungrénses, Coloniénses et Trevirénses, aliósque finítimos pópulos ad Christi fidem perdúxit.
\switchcolumn
\selectlanguage{english}
At Treves, the holy bishop Maternus, 
 a disciple of the blessed apostle Peter, who brought to the faith of Christ 
 the inhabitants of Tongres, Cologne, Treves, and of the neighbouring 
 country.
\switchcolumn*
\selectlanguage{latin}
Romæ sancti Crescéntii 
 púeri, qui sancti Euthymii fílius éxstitit; atque, in persecutióne 
 Diocletiáni, sub Turpílio Júdice, via Salária, gládio percússus, occúbuit.
\switchcolumn
\selectlanguage{english}
On the Salarian Way at Rome, during 
 the persecution of Diocletian, St. Crescentius, the young son of St. 
 Euthymius, whose life was ended by the sword, under the judge Turpilius.
\switchcolumn*
\selectlanguage{latin}
In Africa pássio 
 sanctórum Mártyrum Crescentiáni, Victóris, Rósulæ et Generális.
\switchcolumn
\selectlanguage{english}
In Africa, the passion of the holy 
 martyrs Crescentian, Victor, Rosula, and Generalis.
\switchcolumn*
\selectlanguage{latin}
\end{paracol}


% ---- martyrology/mart09/mart0915.htm
\needspace{10\baselineskip}
\begin{paracol}{2}
\selectlanguage{latin}
\begin{center}{\color{gregoriocolor} Décimo séptimo Kaléndas Octóbris. 
 Luna\dots\ }\end{center}
\switchcolumn
\selectlanguage{english}
\begin{center}{\color{gregoriocolor} The   15\textsuperscript{th} Day of 
 September. The\dots\ Day of the Moon.}\end{center}
\end{paracol}

\noindent\begin{tabularx}{\linewidth}{*{19}{>{\centering\arraybackslash}X}}
 \textcolor{gregoriocolor}{a} & \textcolor{gregoriocolor}{b} & \textcolor{gregoriocolor}{c} & \textcolor{gregoriocolor}{d} & \textcolor{gregoriocolor}{e} & \textcolor{gregoriocolor}{f} & \textcolor{gregoriocolor}{g} & \textcolor{gregoriocolor}{h} & \textcolor{gregoriocolor}{i} & \textcolor{gregoriocolor}{k} & \textcolor{gregoriocolor}{l} & \textcolor{gregoriocolor}{m} & \textcolor{gregoriocolor}{n} & \textcolor{gregoriocolor}{p} & \textcolor{gregoriocolor}{q} & \textcolor{gregoriocolor}{r} & \textcolor{gregoriocolor}{s} & \textcolor{gregoriocolor}{t} & \textcolor{gregoriocolor}{u} \\
 23 & 24 & 25 & 26 & 27 & 28 & 29 & 30 & 1 & 2 & 3 & 4 & 5 & 6 & 7 & 8 & 9 & 10 & 11 \\
\end{tabularx}
\vspace{0.5\baselineskip}
\noindent\begin{tabularx}{\linewidth}{*{12}{>{\centering\arraybackslash}X}}
 \textcolor{gregoriocolor}{A} & \textcolor{gregoriocolor}{B} & \textcolor{gregoriocolor}{C} & \textcolor{gregoriocolor}{D} & \textcolor{gregoriocolor}{E} & F & \textcolor{gregoriocolor}{F} & \textcolor{gregoriocolor}{G} & \textcolor{gregoriocolor}{H} & \textcolor{gregoriocolor}{M} & \textcolor{gregoriocolor}{N} & \textcolor{gregoriocolor}{P} \\
 12 & 13 & 14 & 15 & 16 & 17 & 17 & 18 & 19 & 20 & 21 & 22 \\
\end{tabularx}

\begin{paracol}{2}
\selectlanguage{latin}
\lettrine[lines=2]{O}{ctáva} Nativitátis 
 beátæ Maríæ Vírginis.
\switchcolumn
\selectlanguage{english}
\lettrine[lines=2]{T}{he} Octave of the Nativity of the 
 Blessed Virgin Mary.
\switchcolumn*
\selectlanguage{latin}
Festum Septem Dolórum 
 ejúsdem beatíssimæ Vírginis Maríæ.
\switchcolumn
\selectlanguage{english}
The feast of the Seven Sorrows of 
 the same most Blessed Virgin Mary.
\switchcolumn*
\selectlanguage{latin}
Romæ, via Nomentána, 
 natális beáti Nicomédis, Presbyteri et Mártyris, qui, cum díceret 
 compelléntibus se sacrificáre: « Ego non sacrífico nisi Deo omnipoténti, qui 
 regnat in cælis », plumbátis diutíssime cæsus est, atque in eo torménto 
 migrávit ad Dóminum.
\switchcolumn
\selectlanguage{english}
At Rome, on the Via Nomentana, the 
 birthday of blessed Nicomedes, priest and martyr. Because he said to 
 those who would compel him to sacrifice: ``I offer sacrifice only to the 
 omnipotent God who reigneth in heaven,'' he was for a long time scourged with 
 leaded whips, and thus passed to the Lord.
\switchcolumn*
\selectlanguage{latin}
Córdubæ, in Hispánia, 
 sanctórum Mártyrum Emilæ Diáconi, et Jeremíæ, qui, in persecutióne Arábica, 
 post longam cárceris maceratiónem, demum, cervícibus pro Christo abscíssis, 
 martyrium complevérunt.
\switchcolumn
\selectlanguage{english}
At Cordova in Spain, the holy 
 martyrs Emilas, deacon, and Jeremias, who ended their martyrdom in the 
 persecution of the Arabs by being beheaded after a long stay in prison.
\switchcolumn*
\selectlanguage{latin}
In território 
 Cabillonénsi sancti Valeriáni Mártyris, quem Priscus Præses, suspénsum et 
 gravi ungulárum laceratióne cruciátum, tandem, cum in Christi confessióne 
 vidéret immóbilem ac læto ánimo in ejus láudibus permanéntem, gládio 
 animadvérti præcépit.
\switchcolumn
\selectlanguage{english}
In the diocese of Chalons, St. 
 Valerian, martyr, who was suspended on high by the governor Priscus, and 
 tortured with iron hooks. Remaining immovable in the confession of 
 Christ, and continuing joyfully to praise him, he was struck with the sword 
 by order of the same magistrate.
\switchcolumn*
\selectlanguage{latin}
Hadrianópoli, in 
 Thrácia, sanctórum Mártyrum Máximi, Theodóri et Asclepiódoti; qui sub 
 Maximiáno Imperatóre coronáti sunt.
\switchcolumn
\selectlanguage{english}
At Adrianople in Thrace, the holy martyrs 
 Maximus, Theodore, and Asclepiodotus, who were crowned under Emperor 
 Maximian.
\switchcolumn*
\selectlanguage{latin}
Item sancti Porphyrii 
 mimi, qui, coram Juliáno Apóstata, per jocum suscípiens baptísmum, Dei 
 virtúte derepénte mutátur, et Christiánum se esse profitétur; ac mox, ipsíus 
 Imperatóris mandáto secúri percússus, martyrio coronátur.
\switchcolumn
\selectlanguage{english}
Also, St. Porphyry, a comedian, who 
 was baptized in jest in the presence of Julian the Apostate, but was 
 suddenly converted by the power of God and declared himself a Christian. 
 By order of the emperor he was thereupon struck with an axe, and thus 
 crowned with martyrdom.
\switchcolumn*
\selectlanguage{latin}
Eódem die sancti Nicétæ 
 Gothi, qui ab Athanaríco Rege, ob cathólicam fidem, jussus est igne combúri.
\switchcolumn
\selectlanguage{english}
On the same day, St. Nicetas, a 
 Goth, who was burned alive for the Catholic faith by order of King Athanaric.
\switchcolumn*
\selectlanguage{latin}
Marcianópoli, in 
 Thrácia, sanctæ Melitínæ Mártyris, quæ, sub Antoníno Imperatóre et Antíocho 
 Præside, cum ad Gentílium fana semel et íterum ducta esset, atque idóla 
 semper corrúerent, ídeo suspénsa et laniáta est, ac demum cápite plexa.
\switchcolumn
\selectlanguage{english}
At Marcianapolis in Thrace, St. 
 Melitina, a martyr, in the time of Emperor Antoninus and the governor 
 Antiochus. She was twice led to the temples of the heathens, and since 
 the idols fell to the ground each time, she was hanged and torn, and finally 
 beheaded.
\switchcolumn*
\selectlanguage{latin}
Tulli, in Gállia, 
 sancti Apri Epíscopi.
\switchcolumn
\selectlanguage{english}
At Toul in France, St. Aper, bishop.
\switchcolumn*
\selectlanguage{latin}
Item sancti Leobíni, 
 Epíscopi Carnuténsis.
\switchcolumn
\selectlanguage{english}
Also, St. Leobinus, bishop of 
 Chartres.
\switchcolumn*
\selectlanguage{latin}
Lugdúni, in Gállia, 
 sancti Albíni Epíscopi.
\switchcolumn
\selectlanguage{english}
At Lyons in France, St. Albinus, 
 bishop.
\switchcolumn*
\selectlanguage{latin}
Eódem die deposítio 
 sancti Aichárdi Abbátis.
\switchcolumn
\selectlanguage{english}
On the same day, the death of St. 
 Aichard, abbot.
\switchcolumn*
\selectlanguage{latin}
In Gállia sanctæ 
 Eutrópiæ Víduæ.
\switchcolumn
\selectlanguage{english}
In France, St. Eutropia, widow.
\switchcolumn*
\selectlanguage{latin}
Génuæ sanctæ Catharínæ 
 Víduæ, contémptu mundi et caritáte in Deum insígnis.
\switchcolumn
\selectlanguage{english}
In Genoa, St. Catherine, a widow, 
 renowned for her contempt of the world and her love of God.
\switchcolumn*
\selectlanguage{latin}
\end{paracol}


% ---- martyrology/mart09/mart0916.htm
\needspace{10\baselineskip}
\begin{paracol}{2}
\selectlanguage{latin}
\begin{center}{\color{gregoriocolor} Sextodécimo Kaléndas Octóbris. 
 Luna\dots\ }\end{center}
\switchcolumn
\selectlanguage{english}
\begin{center}{\color{gregoriocolor} The   16\textsuperscript{th} Day of 
 September. The\dots\ Day of the Moon.}\end{center}
\end{paracol}

\noindent\begin{tabularx}{\linewidth}{*{19}{>{\centering\arraybackslash}X}}
 \textcolor{gregoriocolor}{a} & \textcolor{gregoriocolor}{b} & \textcolor{gregoriocolor}{c} & \textcolor{gregoriocolor}{d} & \textcolor{gregoriocolor}{e} & \textcolor{gregoriocolor}{f} & \textcolor{gregoriocolor}{g} & \textcolor{gregoriocolor}{h} & \textcolor{gregoriocolor}{i} & \textcolor{gregoriocolor}{k} & \textcolor{gregoriocolor}{l} & \textcolor{gregoriocolor}{m} & \textcolor{gregoriocolor}{n} & \textcolor{gregoriocolor}{p} & \textcolor{gregoriocolor}{q} & \textcolor{gregoriocolor}{r} & \textcolor{gregoriocolor}{s} & \textcolor{gregoriocolor}{t} & \textcolor{gregoriocolor}{u} \\
 24 & 25 & 26 & 27 & 28 & 29 & 30 & 1 & 2 & 3 & 4 & 5 & 6 & 7 & 8 & 9 & 10 & 11 & 12 \\
\end{tabularx}
\vspace{0.5\baselineskip}
\noindent\begin{tabularx}{\linewidth}{*{12}{>{\centering\arraybackslash}X}}
 \textcolor{gregoriocolor}{A} & \textcolor{gregoriocolor}{B} & \textcolor{gregoriocolor}{C} & \textcolor{gregoriocolor}{D} & \textcolor{gregoriocolor}{E} & F & \textcolor{gregoriocolor}{F} & \textcolor{gregoriocolor}{G} & \textcolor{gregoriocolor}{H} & \textcolor{gregoriocolor}{M} & \textcolor{gregoriocolor}{N} & \textcolor{gregoriocolor}{P} \\
 13 & 14 & 15 & 16 & 17 & 18 & 18 & 19 & 20 & 21 & 22 & 23 \\
\end{tabularx}

\begin{paracol}{2}
\selectlanguage{latin}
\lettrine[lines=2]{S}{anctórum} Mártyrum 
 Cornélii Papæ, et Cypriáni, Carthaginénsis Epíscopi, quorum memória décimo 
 octávo Kaléndas Octóbris recólitur.
\switchcolumn
\selectlanguage{english}
\lettrine[lines=2]{T}{he} holy martyrs Cornelius, pope, 
 and Cyprian, bishop of Carthage, who were mentioned on the14th of September.
\switchcolumn*
\selectlanguage{latin}
Chalcédone natális 
 sanctæ Euphémiæ, Vírginis et Mártyris; quæ, sub Diocletiáno Imperatóre et 
 Prisco Procónsule, torménta, cárceres, vérbera, arguménta rotárum, ignes, 
 póndera lápidum, béstias, plagas virgárum, serras acútas, sartágines ignítas 
 pro Christo superávit. Sed, rursus in theátrum ad béstias ducta, ibi, 
 cum orásset ad Dóminum ut jam spíritum suum suscíperet, una ex iis morsum 
 sancto córpori infigénte, céteris pedes ejus lambéntibus, immaculátum 
 spíritum Deo réddidit.
\switchcolumn
\selectlanguage{english}
At Chalcedon, the birthday of St. 
 Euphemia, virgin and martyr, under Emperor Diocletian and the proconsul 
 Priscus. For her faith in our Lord she was subjected to tortures, 
 imprisonment, blows, the torment of the wheel, fire, the crushing weight of 
 stones, the teeth of the beasts, scourging with rods, the cutting of sharp 
 saws, and burning pans, all of which she survived. But when she was 
 again exposed to the beasts in the amphitheatre, praying to our Lord to 
 receive her spirit, one of the animals inflicted a bite on her holy body 
 although the rest of them licked her feet, and she yielded her unspotted 
 soul unto God.
\switchcolumn*
\selectlanguage{latin}
Romæ sanctórum Mártyrum 
 Lúciæ, nóbilis matrónæ, et Geminiáni; quos ambos Diocletiánus Imperátor, 
 pœnis gravíssimis afflíctos diúque tortos, tandem, post laudábilem martyrii 
 victóriam, gládio animadvérti præcépit.
\switchcolumn
\selectlanguage{english}
At Rome, the holy martyrs Lucy, a 
 noble matron, and Geminian, who were subjected to grievous afflictions and 
 were for a long time tortured by the command of Emperor Diocletian. 
 Finally, being put to the sword, they obtained the glorious victory of 
 martyrdom.
\switchcolumn*
\selectlanguage{latin}
Natális sancti Martíni 
 Primi, Papæ et Mártyris, qui, cum Sérgium ac Paulum et Pyrrhum hæréticos, 
 Romæ coácta Synodo, condemnásset, ídeo, jussu Constántis, Imperatóris 
 hærétici, per fraudem captus et Constantinópolim perdúctus, in Chersonésum 
 relegátus est, ibíque, ob cathólicam fidem ærúmnis conféctus, vitam finívit, 
 multísque miráculis cláruit. Ejus corpus, póstea Romam translátum, in 
 Ecclésia sanctórum Silvéstri et Martíni cónditum fuit. Ipsíus tamen 
 festívitas prídie Idus Novémbris celebrátur.
\switchcolumn
\selectlanguage{english}
The birthday of St. Martin I, pope 
 and martyr. He had called together a council at Rome and condemned the 
 heretics Sergius, Paul and Pyrrhus. By order of the heretical Emperor 
 Constantius he was taken prisoner through a deceit, brought to 
 Constantinople, and exiled to the Chersonese. There he ended his life, 
 worn out with his labours for the Catholic faith and favoured with many 
 virtues. His body was afterwards brought to Rome and buried in the 
 church of Saints Sylvester and Martin. His feast, however, is observed 
 on the 12th of November.
\switchcolumn*
\selectlanguage{latin}
Romæ item natális 
 sanctæ Cæcíliæ, Vírginis et Mártyris, quæ sponsum suum Valeriánum et fratrem 
 ejus Tibúrtium ad credéndum in Christum perdúxit, et ad martyrium incitávit. 
 Hanc Almáchius, Urbis Præféctus, post eórum martyrium tenéri, atque illústri 
 passióne post ignem superátum, fecit gládio consummári, témpore Marci 
 Aurélii Sevéri Alexándri Imperatóris. Ejus vero festum recólitur décimo Kaléndas Decémbris.
\switchcolumn
\selectlanguage{english}
Also at Rome, the birthday of St. 
 Cecilia, virgin and martyr. She brought her husband and brother 
 Tiburtius to the faith of Christ and afterwards encouraged them on to 
 martyrdom. Almachius, prefect of the city, after their martyrdom, had 
 her arrested and slain by the sword, after she had endured many trials and 
 had passed through fire unhurt. This was in the reign of Emperor 
 Marcus Aurelius Severus Alexander. Her feast is celebrated on the 22nd 
 of November.
\switchcolumn*
\selectlanguage{latin}
Heracléæ, in Thrácia, 
 sanctæ Sebastiánæ Mártyris, quæ a beáto Paulo Apóstolo ad Christi fidem est 
 perdúcta; atque, sub Domitiáno Imperatóre et Sérgio Præside, váriis modis 
 tentáta, gládio tandem cæsa est.
\switchcolumn
\selectlanguage{english}
At Heraclea in Thrace, under Emperor 
 Domitian and the governor Sergius, St. Sebastiana, martyr. Being 
 brought to the faith of Christ by the blessed apostle Paul, she was 
 tormented in various ways and finally beheaded.
\switchcolumn*
\selectlanguage{latin}
Romæ, via Flamínia, 
 sanctórum Mártyrum Abúndii Presbyteri, et Abundántii Diáconi, quos 
 Diocletiánus Imperátor, una cum illústri viro Marciáno et ejus fílio Joánne, 
 quem illi a mórtuis suscitáverant, gládio feríri jussit, décimo ab Urbe 
 lápide.
\switchcolumn
\selectlanguage{english}
At Rome, at a place on the Flaminian 
 Way ten miles from the city, the holy martyrs Abundius, a priest, and 
 Abundantius, a deacon, whom Emperor Diocletian ordered to be struck with the 
 sword, together with Marcian, an illustrious man, and his son John, whom 
 they raised from the dead.
\switchcolumn*
\selectlanguage{latin}
Córdubæ, in Hispánia, 
 sanctórum Mártyrum Rogélli et Servidéi, qui, mánibus pedibúsque abscíssis, 
 ad últimum decolláti sunt.
\switchcolumn
\selectlanguage{english}
At Cordova in Spain, the holy 
 martyrs Rogellus and Servusdeus, who were beheaded after their hands and 
 feet had been cut off.
\switchcolumn*
\selectlanguage{latin}
Apud Cándidam Casam, in 
 Scótia, sancti Niniáni, Epíscopi et Confessóris.
\switchcolumn
\selectlanguage{english}
At Whithorn in Scotland, St. Ninian, 
 bishop and confessor.
\switchcolumn*
\selectlanguage{latin}
In Anglia sanctæ Edíthæ 
 Vírginis, Regis Anglórum Edgári fíliæ, quæ, in monastério a téneris annis 
 Deo dicáta, sæculum hoc ignorávit pótius quam relíquit.
\switchcolumn
\selectlanguage{english}
In England, St. Edith, virgin, 
 daughter of the English King Edgar. She was consecrated to God in a 
 monastery from her earliest years, and it may be said rather that she never 
 knew the world than that she forsook it.
\switchcolumn*
\selectlanguage{latin}
In monte Cassíno beáti 
 Victóris Papæ Tértii, qui, sancti Gregórii Séptimi succéssor, Apostólicam 
 Sedem novo splendóre illustrávit, insígnem de Saracénis triúmphum divína ope 
 consecútus. Cultum, ab immemorábili témpore eídem exhíbitum, Leo 
 Décimus tértius, Póntifex Máximus, ratum hábuit et confirmávit.
\switchcolumn
\selectlanguage{english}
At Monte Cassino, the blessed Pope 
 Victor III, successor of Pope St. Gregory VII, who shed a fresh lustre on 
 the Apostolic See, and by God's help gained a famous victory over the 
 Saracens. Pope Leo XIII approved and confirmed the veneration given 
 him from time immemorial.
\switchcolumn*
\selectlanguage{latin}
\end{paracol}


% ---- martyrology/mart09/mart0917.htm
\needspace{10\baselineskip}
\begin{paracol}{2}
\selectlanguage{latin}
\begin{center}{\color{gregoriocolor} Quintodécimo Kaléndas Octóbris. 
 Luna\dots\ }\end{center}
\switchcolumn
\selectlanguage{english}
\begin{center}{\color{gregoriocolor} The   17\textsuperscript{th} Day of 
 September. The\dots\ Day of the Moon.}\end{center}
\end{paracol}

\noindent\begin{tabularx}{\linewidth}{*{19}{>{\centering\arraybackslash}X}}
 \textcolor{gregoriocolor}{a} & \textcolor{gregoriocolor}{b} & \textcolor{gregoriocolor}{c} & \textcolor{gregoriocolor}{d} & \textcolor{gregoriocolor}{e} & \textcolor{gregoriocolor}{f} & \textcolor{gregoriocolor}{g} & \textcolor{gregoriocolor}{h} & \textcolor{gregoriocolor}{i} & \textcolor{gregoriocolor}{k} & \textcolor{gregoriocolor}{l} & \textcolor{gregoriocolor}{m} & \textcolor{gregoriocolor}{n} & \textcolor{gregoriocolor}{p} & \textcolor{gregoriocolor}{q} & \textcolor{gregoriocolor}{r} & \textcolor{gregoriocolor}{s} & \textcolor{gregoriocolor}{t} & \textcolor{gregoriocolor}{u} \\
 25 & 26 & 27 & 28 & 29 & 30 & 1 & 2 & 3 & 4 & 5 & 6 & 7 & 8 & 9 & 10 & 11 & 12 & 13 \\
\end{tabularx}
\vspace{0.5\baselineskip}
\noindent\begin{tabularx}{\linewidth}{*{12}{>{\centering\arraybackslash}X}}
 \textcolor{gregoriocolor}{A} & \textcolor{gregoriocolor}{B} & \textcolor{gregoriocolor}{C} & \textcolor{gregoriocolor}{D} & \textcolor{gregoriocolor}{E} & F & \textcolor{gregoriocolor}{F} & \textcolor{gregoriocolor}{G} & \textcolor{gregoriocolor}{H} & \textcolor{gregoriocolor}{M} & \textcolor{gregoriocolor}{N} & \textcolor{gregoriocolor}{P} \\
 14 & 15 & 16 & 17 & 18 & 19 & 19 & 20 & 21 & 22 & 23 & 24 \\
\end{tabularx}

\begin{paracol}{2}
\selectlanguage{latin}
\lettrine[lines=2]{I}{n} monte Alvérniæ, in 
 Etrúria, commemorátio Impressiónis sacrórum Stígmatum, quibus sanctus 
 Francíscus, Ordinis Minórum Institútor, in suis mánibus, pédibus et látere, 
 mirábili Dei grátia, impréssus fuit.
\switchcolumn
\selectlanguage{english}
\lettrine[lines=2]{T}{he} commemoration of the Impression 
 of the Sacred Stigmata which St. Francis, founder of the Order of Friars 
 Minor, received through a wonderful favour of God in his hands, feet, and 
 side, at Mount Alverina in Etruria.
\switchcolumn*
\selectlanguage{latin}
Romæ natális sancti 
 Robérti Bellarmíno, Confessóris, e Societáte Jesu, atque Cardinális et 
 Capuáni olim Epíscopi, sanctitáte, doctrína, et plúrimis ad cathólicæ fídei 
 et Apostólicæ Sedis defensiónem suscéptis labóribus claríssimi; quem Pius 
 Undécimus, Póntifex Máximus, Sanctórum honóribus auxit et universális 
 Ecclésiæ Doctórem declarávit, ejúsque festum tértio Idus Maji recoléndum 
 indíxit.
\switchcolumn
\selectlanguage{english}
At Rome, the birthday of St. Robert 
 Bellarmine of the Society of Jesus, confessor and cardinal, and also 
 formerly bishop of Capua. He is noted for his holiness, learning, and 
 the many great tasks he performed in defence of the Catholic faith and the 
 Apostolic See. Pope Pius XI bestowed on him the honours of the saints, 
 declared him to be a doctor of the universal Church, and appointed the 13th 
 of May as his feast day.
\switchcolumn*
\selectlanguage{latin}
Romæ, via Tiburtína, 
 natális sancti Justíni, Presbyteri et Mártyris; qui, in persecutióne 
 Valeriáni et Galliéni, ob confessiónis glóriam fuit insígnis. Hic 
 beáti Pontíficis Xysti Secúndi, Lauréntii, Hippólyti aliorúmque plurimórum 
 Sanctórum córpora sepelívit, ac demum, sub Cláudio, martyrium consummávit.
\switchcolumn
\selectlanguage{english}
At Rome, on the road to Tivoli, the 
 birthday of St. Justin, priest and martyr, who distinguished himself by a 
 glorious confession of the faith during the persecution of Valerian and 
 Gallienus. He buried the bodies of the blessed Pontiff Sixtus II, of 
 Lawrence, Hippolytus, and many other saints, and finally completed his 
 martyrdom under Claudius.
\switchcolumn*
\selectlanguage{latin}
Item Romæ sanctórum 
 Mártyrum Narcíssi et Crescentiónis.
\switchcolumn
\selectlanguage{english}
Also at Rome, the holy martyrs 
 Narcissus and Crescentio.
\switchcolumn*
\selectlanguage{latin}
Apud Leódium, in Bélgio, 
 beáti Lambérti, Epíscopi Trajecténsis, qui, cum régiam domum zelo religiónis 
 increpásset, a nocéntibus ínnocens occísus est, sicque aulam regni cæléstis 
 perpétuo victúrus intrávit.
\switchcolumn
\selectlanguage{english}
At Liege in Belgium, blessed 
 Lambert, bishop of Maestricht. Through his zeal for religion he 
 rebuked the royal family, and was undeservedly put to death by the guilty, 
 and thus he entered the court of the heavenly kingdom, to enjoy it forever.
\switchcolumn*
\selectlanguage{latin}
Cæsaraugústæ, in 
 Hispánia, sancti Petri de Arbues, primi in Aragóniæ regno Quæsitóris fídei; 
 quem, a relápsis Judæis ob eándem, quam pro múnere suo strénue tuebátur, 
 cathólicam fidem, immániter trucidátum, sanctórum Mártyrum catálogo Pius 
 Papa Nonus adjúnxit.
\switchcolumn
\selectlanguage{english}
At Saragossa in Spain, St. Peter of 
 Arbues, first inquisitor of the faith in the kingdom of Aragon, who received 
 the palm of martyrdom by being barbarously massacred by apostate Jews for 
 courageously defending the Catholic faith, according to the duties of his 
 office. He was added to the list of martyr saints by Pius IX.
\switchcolumn*
\selectlanguage{latin}
In Británnia sanctórum 
 Mártyrum Sócratis et Stéphani.
\switchcolumn
\selectlanguage{english}
In England, the holy martyrs 
 Socrates and Stephen.
\switchcolumn*
\selectlanguage{latin}
Noviodúni, in Gálliis, 
 sanctórum Mártyrum Valeriáni, Macríni et Gordiáni.
\switchcolumn
\selectlanguage{english}
At Noyon in France, the holy martyrs 
 Valerian, Macrinus, and Gordian.
\switchcolumn*
\selectlanguage{latin}
Augustodúni sancti 
 Flocélli púeri, qui, sub Antoníno Imperatóre et Valeriáno Præside, multa 
 passus, demum, a feris discérptus, martyrii corónam adéptus est.
\switchcolumn
\selectlanguage{english}
At Autun, under Emperor Antoninus 
 and the governor Valerian, St. Flocellus, a boy, who, after many sufferings, 
 was torn to pieces by wild beasts, and thus won the crown of martyrs.
\switchcolumn*
\selectlanguage{latin}
Córdubæ, in Hispánia, 
 sanctæ Colúmbæ, Vírginis et Mártyris.
\switchcolumn
\selectlanguage{english}
At Cordova in Spain, St. Columba, 
 virgin and martyr.
\switchcolumn*
\selectlanguage{latin}
In Phrygia sanctæ 
 Ariádne Mártyris, sub Hadriáno Imperatóre.
\switchcolumn
\selectlanguage{english}
In Phrygia, St. Ariadne, martyr, 
 under Emperor Hadrian.
\switchcolumn*
\selectlanguage{latin}
Eódem die sanctæ 
 Agathoclíæ, quæ, cum esset ancílla cujúsdam mulíeris infidélis, a dómina sua 
 longo témpore verbéribus aliísque ærúmnis est vexáta ut Christum negáret; 
 demum, obláta Júdici ac sævius laniáta, et nihilóminus in confessióne fídei 
 persístens, ídeo, post excisiónem linguæ, in ignem projécta est.
\switchcolumn
\selectlanguage{english}
On the same day, St. Agathoclia, 
 servant of an infidel woman, who was for a long time subjected by her to 
 blows and other afflictions that she might deny Christ. She was 
 finally presented to the judge and cruelly lacerated, but since she 
 persisted in confessing the faith, they cut off her tongue and threw her 
 into the flames.
\switchcolumn*
\selectlanguage{latin}
Medioláni deposítio 
 sancti Sátyri Confessóris, cujus insígnia mérita sanctus Ambrósius, ejus 
 frater, commémorat.
\switchcolumn
\selectlanguage{english}
At Milan, the death of St. Satyrus, 
 confessor, whose distinguished merits are mentioned by his brother, St. 
 Ambrose.
\switchcolumn*
\selectlanguage{latin}
Apud Bíngiam, in 
 diœcési Moguntinénsi, sanctæ Hildegárdis Vírginis.
\switchcolumn
\selectlanguage{english}
At Bingen, in the diocese of Mainz, 
 St. Hildegard, virgin.
\switchcolumn*
\selectlanguage{latin}
Romæ sanctæ Theodóræ 
 matrónæ, quæ, in persecutióne Diocletiáni, sanctis Martyribus sédulo 
 ministrábat.
\switchcolumn
\selectlanguage{english}
At Rome, St. Theodora, a matron who 
 zealously ministered to the martyrs in the persecution of Diocletian.
\switchcolumn*
\selectlanguage{latin}
\end{paracol}


% ---- martyrology/mart09/mart0918.htm
\needspace{10\baselineskip}
\begin{paracol}{2}
\selectlanguage{latin}
\begin{center}{\color{gregoriocolor} Quartodécimo Kaléndas Octóbris. 
 Luna\dots\ }\end{center}
\switchcolumn
\selectlanguage{english}
\begin{center}{\color{gregoriocolor} The   18\textsuperscript{th} Day of 
 September. The\dots\ Day of the Moon.}\end{center}
\end{paracol}

\noindent\begin{tabularx}{\linewidth}{*{19}{>{\centering\arraybackslash}X}}
 \textcolor{gregoriocolor}{a} & \textcolor{gregoriocolor}{b} & \textcolor{gregoriocolor}{c} & \textcolor{gregoriocolor}{d} & \textcolor{gregoriocolor}{e} & \textcolor{gregoriocolor}{f} & \textcolor{gregoriocolor}{g} & \textcolor{gregoriocolor}{h} & \textcolor{gregoriocolor}{i} & \textcolor{gregoriocolor}{k} & \textcolor{gregoriocolor}{l} & \textcolor{gregoriocolor}{m} & \textcolor{gregoriocolor}{n} & \textcolor{gregoriocolor}{p} & \textcolor{gregoriocolor}{q} & \textcolor{gregoriocolor}{r} & \textcolor{gregoriocolor}{s} & \textcolor{gregoriocolor}{t} & \textcolor{gregoriocolor}{u} \\
 26 & 27 & 28 & 29 & 30 & 1 & 2 & 3 & 4 & 5 & 6 & 7 & 8 & 9 & 10 & 11 & 12 & 13 & 14 \\
\end{tabularx}
\vspace{0.5\baselineskip}
\noindent\begin{tabularx}{\linewidth}{*{12}{>{\centering\arraybackslash}X}}
 \textcolor{gregoriocolor}{A} & \textcolor{gregoriocolor}{B} & \textcolor{gregoriocolor}{C} & \textcolor{gregoriocolor}{D} & \textcolor{gregoriocolor}{E} & F & \textcolor{gregoriocolor}{F} & \textcolor{gregoriocolor}{G} & \textcolor{gregoriocolor}{H} & \textcolor{gregoriocolor}{M} & \textcolor{gregoriocolor}{N} & \textcolor{gregoriocolor}{P} \\
 15 & 16 & 17 & 18 & 19 & 20 & 20 & 21 & 22 & 23 & 24 & 25 \\
\end{tabularx}

\begin{paracol}{2}
\selectlanguage{latin}
\lettrine[lines=2]{A}{uximi,} in Picéno, sancti Joséphi a Cupertíno, Sacerdótis ex Ordine Minórum Conventuálium et 
 Confessóris; quem Clemens Papa Décimus tértius in Sanctórum númerum rétulit.
\switchcolumn
\selectlanguage{english}
\lettrine[lines=2]{A}{t} Osimo in Piceno, St. Joseph of 
 Cupertino, priest and confessor of the Order of Friars Minor Conventual, who was placed 
 among the saints by Clement XIII.
\switchcolumn*
\selectlanguage{latin}
In Chálcide Græciæ 
 natális sancti Methódii, qui prius Olympii, in Lycia, et póstea Tyri, in 
 Phœnícia, éxstitit Epíscopus, sermónis nitóre ac doctrína claríssimus; atque, 
 ad extrémum novíssimæ persecutiónis (ut scribit sanctus Hierónymus), 
 martyrio coronátus est.
\switchcolumn
\selectlanguage{english}
In Chalcis of Greece, the birthday 
 of St. Methodius, bishop of Olympius in Lycia and afterwards of Tyre in 
 Phoenicia, most renowned for eloquence and learning. St. Jerome says 
 that he won the martyr's crown at the end of the last persecution.
\switchcolumn*
\selectlanguage{latin}
In território Viennénsi 
 sancti Ferréoli Mártyris, qui, cum esset tribuníciæ potestátis, jussu 
 impiíssimi Præsidis Crispíni tentus, et primo crudelíssime verberátus, 
 deínde, gravi catenárum póndere onústus, in tetérrimum cárcerem trusus est; 
 unde, solútis Dei nutu vínculis et jánuis cárceris patefáctis, éxiens, ab 
 insequéntibus íterum est captus, ac martyrii palmam obtruncatióne cápitis 
 percépit.
\switchcolumn
\selectlanguage{english}
In the diocese of Vienne, the holy 
 martyr Ferreol, a tribune, who was arrested by order of the impious governor 
 Crispinus, most cruelly scourged, loaded with heavy chains, and cast into a 
 dark dungeon. A miracle broke his bonds and opened the doors of the 
 prison, from which he made his escape, but he was taken again by his 
 pursuers and received the palm of martyrdom by being beheaded.
\switchcolumn*
\selectlanguage{latin}
Item sanctárum Mártyrum 
 Sophíæ et Irénes.
\switchcolumn
\selectlanguage{english}
Also, the Saints Sophia and Irene, 
 martyrs.
\switchcolumn*
\selectlanguage{latin}
Medioláni sancti 
 Eustórgii Primi, ejúsdem civitátis Epíscopi, beáti Ambrósii testimónio 
 célebris.
\switchcolumn
\selectlanguage{english}
At Milan, St. Eustorgius, first 
 bishop of that city, highly praised by blessed Ambrose.
\switchcolumn*
\selectlanguage{latin}
Gortynæ, in Creta, 
 sancti Euménii, Epíscopi et Confessóris.
\switchcolumn
\selectlanguage{english}
At Gortyna in Crete, St. Eumenius, 
 bishop and confessor.
\switchcolumn*
\selectlanguage{latin}
\end{paracol}


% ---- martyrology/mart09/mart0919.htm
\needspace{10\baselineskip}
\begin{paracol}{2}
\selectlanguage{latin}
\begin{center}{\color{gregoriocolor} Tertiodécimo Kaléndas Octóbris. 
 Luna\dots\ }\end{center}
\switchcolumn
\selectlanguage{english}
\begin{center}{\color{gregoriocolor} The   19\textsuperscript{th} Day of 
 September. The\dots\ Day of the Moon.}\end{center}
\end{paracol}

\noindent\begin{tabularx}{\linewidth}{*{19}{>{\centering\arraybackslash}X}}
 \textcolor{gregoriocolor}{a} & \textcolor{gregoriocolor}{b} & \textcolor{gregoriocolor}{c} & \textcolor{gregoriocolor}{d} & \textcolor{gregoriocolor}{e} & \textcolor{gregoriocolor}{f} & \textcolor{gregoriocolor}{g} & \textcolor{gregoriocolor}{h} & \textcolor{gregoriocolor}{i} & \textcolor{gregoriocolor}{k} & \textcolor{gregoriocolor}{l} & \textcolor{gregoriocolor}{m} & \textcolor{gregoriocolor}{n} & \textcolor{gregoriocolor}{p} & \textcolor{gregoriocolor}{q} & \textcolor{gregoriocolor}{r} & \textcolor{gregoriocolor}{s} & \textcolor{gregoriocolor}{t} & \textcolor{gregoriocolor}{u} \\
 27 & 28 & 29 & 30 & 1 & 2 & 3 & 4 & 5 & 6 & 7 & 8 & 9 & 10 & 11 & 12 & 13 & 14 & 15 \\
\end{tabularx}
\vspace{0.5\baselineskip}
\noindent\begin{tabularx}{\linewidth}{*{12}{>{\centering\arraybackslash}X}}
 \textcolor{gregoriocolor}{A} & \textcolor{gregoriocolor}{B} & \textcolor{gregoriocolor}{C} & \textcolor{gregoriocolor}{D} & \textcolor{gregoriocolor}{E} & F & \textcolor{gregoriocolor}{F} & \textcolor{gregoriocolor}{G} & \textcolor{gregoriocolor}{H} & \textcolor{gregoriocolor}{M} & \textcolor{gregoriocolor}{N} & \textcolor{gregoriocolor}{P} \\
 16 & 17 & 18 & 19 & 20 & 21 & 21 & 22 & 23 & 24 & 25 & 26 \\
\end{tabularx}

\begin{paracol}{2}
\selectlanguage{latin}
\lettrine[lines=2]{P}{utéolis,} in Campánia, sanctórum Mártyrum Januárii, Beneventánæ civitátis Epíscopi, ejúsque Diáconi 
 Festi, et Desidérii Lectóris, Sósii, Diáconi Ecclésiæ Misenátis; Próculi, 
 Diáconi Puteoláni; Eutychii et Acútii. Hi omnes, post víncula et cárceres, cápite cæsi sunt, sub Diocletiáno Príncipe. Corpus sancti 
 Januárii delátum fuit Neápolim, atque honorífice in Ecclésia tumulátum; ubi 
 étiam beatíssimi Mártyris sanguis in ampúlla vítrea adhuc servátur, qui, in 
 conspéctu cápitis illíus pósitus, velut recens liquéscere et ebullíre 
 conspícitur.
\switchcolumn
\selectlanguage{english}
\lettrine[lines=2]{A}{t} Pozzuoli in Campania, the holy 
 martyrs Januarius, bishop of Benevento; Festus, his deacon, and Desiderius, 
 a lector, together with Sosius, a deacon of the Church of Miseno; Proculus, 
 deacon of Pozzuoli; Eutychius, and Acutius, who were bound and imprisoned 
 and then beheaded during the reign of Diocletian. The body of St. Januarius 
 was brought to Naples and buried in the church with due honours, where even 
 now the blood of the blessed martyr is kept in a vial, and when placed close 
 to his head is seen to become liquid and bubble up as if it were just taken 
 from his veins.
\switchcolumn*
\selectlanguage{latin}
In Palæstína sanctórum 
 Mártyrum et Ægypti Episcopórum Pélei, Nili et Elíæ; qui, témpore 
 persecutiónis Diocletiáni, cum plúrimis Cléricis, pro Christo sunt igne 
 consúmpti.
\switchcolumn
\selectlanguage{english}
In Palestine, the holy martyrs 
 Peleus, Nilus, and Elias, bishops in Egypt, with many others of the clergy, 
 who were consumed by fire for the sake of Christ during the persecution of 
 Diocletian.
\switchcolumn*
\selectlanguage{latin}
Nucériæ natális 
 sanctórum Mártyrum Felícis et Constántiæ, qui passi sunt sub Neróne.
\switchcolumn
\selectlanguage{english}
At Nocera, the birthday of the holy 
 martyrs Felix and Constantia, who suffered under Nero.
\switchcolumn*
\selectlanguage{latin}
Eódem die sanctórum 
 Mártyrum Tróphimi, Sabbátii et Dorymedóntis, sub Probo Imperatóre. Ex 
 eis Sabbátius Antiochíæ, jussu Attici Præsidis, támdiu flagris cæsus est, 
 donec emítteret spíritum; Tróphimus vero, Synnadam, in Phrygia, ad Perénnium 
 Præsidem missus, ibi, post multos cruciátus, cum Dorymedónte Senatóre, 
 cápitis decollatióne martyrium consummávit.
\switchcolumn
\selectlanguage{english}
Also, the holy martyrs Trophimus, 
 Sabbatius, and Dorymedon, senator, under Emperor Probus. By command of 
 the governor Atticus at Antioch, Sabbatius was scourged until he expired. 
 Trophimus was sent to the governor Perennius at Synnada, where he and the 
 senator Dorymedon completed their martyrdom by being beheaded after enduring 
 many torments.
\switchcolumn*
\selectlanguage{latin}
Eleutherópoli, in 
 Palæstína, sanctæ Susánnæ, Vírginis et Mártyris; quæ, orta ex idolórum 
 sacerdóte Arthémio et Judæa mulíere Martha, ad Christiánam fidem, paréntibus 
 mórtuis, est convérsa, atque ob eándem fidem ab Alexándro Præfécto várie 
 torta et in cárcerem trusa, ibi orans migrávit ad Sponsum.
\switchcolumn
\selectlanguage{english}
At Eleutheropolis in Palestine, St. 
 Susanna, virgin and martyr. She was the daughter of Arthemius, a pagan 
 priest, and of Martha, a Jewish woman, and after the death of her parents 
 she was converted to the Christian faith. For this she was tortured in 
 various ways, and cast in prison by the prefect Alexander, and there gave up 
 her spirit while at prayer.
\switchcolumn*
\selectlanguage{latin}
Córdubæ, in Hispánia, 
 sanctæ Pompósæ, Vírginis et Mártyris; quæ in persecutióne Arábica, ob impávidam Christi confessiónem decolláta gládio, palmam consecúta est.
\switchcolumn
\selectlanguage{english}
At Cordova in Spain, St. Pomposa, 
 virgin and martyr. Because of her fearless witness to Christ she was 
 beheaded in the Arab persecution, and thus obtained the palm of martyrdom.
\switchcolumn*
\selectlanguage{latin}
Cantuáriæ sancti 
 Theodóri Epíscopi, qui, a beáto Vitaliáno Papa in Angliam missus, doctrína 
 et sanctitáte refúlsit.
\switchcolumn
\selectlanguage{english}
At Canterbury, the holy bishop 
 Theodore, who was sent to England by Pope Vitalian, and who was renowned for 
 learning and holiness.
\switchcolumn*
\selectlanguage{latin}
Turónis, in Gállia, 
 sancti Eustóchii Epíscopi, magnárum virtútum viri.
\switchcolumn
\selectlanguage{english}
At Tours in France, St. Eustochius, 
 bishop, a man of great virtue.
\switchcolumn*
\selectlanguage{latin}
In território 
 Lingoniénsi sancti Sequáni, Presbyteri et Confessóris.
\switchcolumn
\selectlanguage{english}
In the diocese of Langres, St. 
 Sequanus, priest and confessor.
\switchcolumn*
\selectlanguage{latin}
Barcinóne, in Hispánia, 
 beátæ Maríæ de Cervellióne, ex Ordine beátæ Maríæ de Mercéde redemptiónis 
 captivórum, Vírginis; quæ, ob præséntem quam invocántibus confert opem, 
 María de Subsídio vulgo nuncupátur.
\switchcolumn
\selectlanguage{english}
At Barcelona in Spain, blessed Mary 
 de Cervellione, virgin, of the Order of Our Lady of Ransom. She is 
 commonly called Mary of Help on account of the prompt assistance she renders 
 to those who invoke her.
\switchcolumn*
\selectlanguage{latin}
In vico Druelle, in 
 diœcési Ruthenénsi, in Gállia, sanctæ Maríæ Guliélmæ-Æmíliæ de Rodat, 
 Vírginis, Congregatiónis Sorórum a sancta Família Fundatrícis, puéllis 
 erudiéndis et egénis sublevándis addictíssimæ, quæ a Pio Duodécimo, 
 Pontífice Máximo, inter sanctas Vírgines reláta est.
\switchcolumn
\selectlanguage{english}
In the village of Druelle, in the 
 diocese of Rodez in France, St. Marie Guillemette Emilie de Rodat, virgin, 
 and foundress of the Congregation of Sisters of the Holy Family, which was 
 established to teach poor and orphaned girls. Pius XII added her name 
 to the number of holy virgins.
\switchcolumn*
\selectlanguage{latin}
\end{paracol}


% ---- martyrology/mart09/mart0920.htm
\needspace{10\baselineskip}
\begin{paracol}{2}
\selectlanguage{latin}
\begin{center}{\color{gregoriocolor} Duodécimo Kaléndas Octóbris. 
 Luna\dots\ }\end{center}
\switchcolumn
\selectlanguage{english}
\begin{center}{\color{gregoriocolor} The   20\textsuperscript{th} Day of 
 September. The\dots\ Day of the Moon.}\end{center}
\end{paracol}

\noindent\begin{tabularx}{\linewidth}{*{19}{>{\centering\arraybackslash}X}}
 \textcolor{gregoriocolor}{a} & \textcolor{gregoriocolor}{b} & \textcolor{gregoriocolor}{c} & \textcolor{gregoriocolor}{d} & \textcolor{gregoriocolor}{e} & \textcolor{gregoriocolor}{f} & \textcolor{gregoriocolor}{g} & \textcolor{gregoriocolor}{h} & \textcolor{gregoriocolor}{i} & \textcolor{gregoriocolor}{k} & \textcolor{gregoriocolor}{l} & \textcolor{gregoriocolor}{m} & \textcolor{gregoriocolor}{n} & \textcolor{gregoriocolor}{p} & \textcolor{gregoriocolor}{q} & \textcolor{gregoriocolor}{r} & \textcolor{gregoriocolor}{s} & \textcolor{gregoriocolor}{t} & \textcolor{gregoriocolor}{u} \\
 28 & 29 & 30 & 1 & 2 & 3 & 4 & 5 & 6 & 7 & 8 & 9 & 10 & 11 & 12 & 13 & 14 & 15 & 16 \\
\end{tabularx}
\vspace{0.5\baselineskip}
\noindent\begin{tabularx}{\linewidth}{*{12}{>{\centering\arraybackslash}X}}
 \textcolor{gregoriocolor}{A} & \textcolor{gregoriocolor}{B} & \textcolor{gregoriocolor}{C} & \textcolor{gregoriocolor}{D} & \textcolor{gregoriocolor}{E} & F & \textcolor{gregoriocolor}{F} & \textcolor{gregoriocolor}{G} & \textcolor{gregoriocolor}{H} & \textcolor{gregoriocolor}{M} & \textcolor{gregoriocolor}{N} & \textcolor{gregoriocolor}{P} \\
 17 & 18 & 19 & 20 & 21 & 22 & 22 & 23 & 24 & 25 & 26 & 27 \\
\end{tabularx}

\begin{paracol}{2}
\selectlanguage{latin}
\lettrine[lines=1]{V}{igília} sancti Matthæi, Apóstoli et Evangelístæ.
\switchcolumn
\selectlanguage{english}
\lettrine[lines=1]{T}{he} vigil of St. Matthew, apostle and evangelist.
\switchcolumn*
\selectlanguage{latin}
Romæ pássio sanctórum 
 Mártyrum Eustáchii et Theopístis uxóris, cum duóbus fíliis Agapíto et 
 Theopísto, qui, sub Hadriáno Imperatóre, damnáti ad béstias, sed Dei ope ab 
 iis nullátenus læsi, tandem, in bovem æneum candéntem inclúsi, martyrium 
 consummárunt.
\switchcolumn
\selectlanguage{english}
At Rome, the holy martyrs Eustace, 
 and Theopistes, his wife, with their two sons, Agapitus and Theopistus. 
 Under Emperor Hadrian they were condemned to be cast to the beasts, but by 
 the power of God they were uninjured by them, so they were shut up in a 
 heated brazen ox, and thus completed their martyrdom.
\switchcolumn*
\selectlanguage{latin}
Cyzici, in Propóntide, 
 natális sanctórum Mártyrum Faustæ Vírginis, et Evilásii, sub Maximiáno 
 Imperatóre; e quibus Fausta, ab eódem Evilásio, idolórum sacerdóte, 
 decalváta et ad turpitúdinem rasa, suspénsa et torta est. Deínde, cum 
 eam vellet médiam secáre, et carnífices lǽdere non valérent, stupens 
 crédidit in Christum Evilásius; et, dum ipse quoque, Imperatóris jussu, 
 fórtiter torquerétur, Fausta, cápite terebráta, clavis toto córpore confíxa 
 et sartágini ignítæ impósita, tandem, cum eódem Evilásio, illam voce de 
 cælis vocánte, transívit ad Dóminum.
\switchcolumn
\selectlanguage{english}
At Cyzicum, on the sea of Marmora, 
 the birthday of the holy martyrs Evilasius and the virgin Fausta, in the 
 time of Emperor Maximian. Fausta's head was shaved to shame her, and 
 she was hung up and tortured by Evilasius, then a pagan priest. But 
 when he wished to have her body cut in two, the executioners could not 
 inflict any injury upon her. Amazed at this prodigy, Evilasius 
 believed in Christ and was cruelly tortured by order of the emperor; at the 
 same time Fausta had her head bored through and her whole body pierced with 
 nails. She was then laid on a heated gridiron, and being called by a 
 celestial voice, went in company with Evilasius to enjoy the blessedness of 
 heaven.
\switchcolumn*
\selectlanguage{latin}
In Phrygia sanctórum 
 Mártyrum Dionysii et Priváti.
\switchcolumn
\selectlanguage{english}
In Phrygia, the holy martyrs Denis 
 and Privatus.
\switchcolumn*
\selectlanguage{latin}
Item sancti Prisci 
 Mártyris, qui, punctim pugiónibus transverberátus, cápite plexus est.
\switchcolumn
\selectlanguage{english}
Also St. Priscus, martyr, whose body 
 was pierced throughout with daggers, after which he was beheaded.
\switchcolumn*
\selectlanguage{latin}
Perge, in Pamphylia, 
 sanctórum Theodóri, et Philíppæ matris, ac Sociórum Mártyrum, sub Antoníno 
 Imperatóre.
\switchcolumn
\selectlanguage{english}
At Pergen in Pamphylia, the Saints 
 Theodore, his mother Philippa, and their fellow martyrs, in the time of 
 Emperor Antoninus.
\switchcolumn*
\selectlanguage{latin}
Carthágine sanctæ 
 Cándidæ, Vírginis et Mártyris; quæ, sub Maximiáno Imperatóre, plagis toto 
 córpore dilaceráta, martyrio coronátur.
\switchcolumn
\selectlanguage{english}
At Carthage, under Emperor Maximian, 
 St. Candida, virgin and martyr. After her body was lacerated by whips 
 she was crowned with martyrdom.
\switchcolumn*
\selectlanguage{latin}
Medioláni sancti 
 Clicérii, Epíscopi et Confessóris.
\switchcolumn
\selectlanguage{english}
At Milan, St. Clicerius, bishop and 
 confessor.
\switchcolumn*
\selectlanguage{latin}
Romæ Translátio 
 córporis sancti Agapíti Primi, Papæ et Confessóris, ex urbe Constantinópoli, 
 in qua Póntifex décimo Kaléndas Maji obdormíerat in Dómino.
\switchcolumn
\selectlanguage{english}
At Rome, the translation of the body 
 of St. Agapitus I, pope and confessor, from the city of Constantinople, in 
 which he died on the 22nd of April.
\switchcolumn*
\selectlanguage{latin}
\end{paracol}


% ---- martyrology/mart09/mart0921.htm
\needspace{10\baselineskip}
\begin{paracol}{2}
\selectlanguage{latin}
\begin{center}{\color{gregoriocolor} Undécimo Kaléndas Octóbris. 
 Luna\dots\ }\end{center}
\switchcolumn
\selectlanguage{english}
\begin{center}{\color{gregoriocolor} The   21\textsuperscript{st} Day of 
 September. The\dots\ Day of the Moon.}\end{center}
\end{paracol}

\noindent\begin{tabularx}{\linewidth}{*{19}{>{\centering\arraybackslash}X}}
 \textcolor{gregoriocolor}{a} & \textcolor{gregoriocolor}{b} & \textcolor{gregoriocolor}{c} & \textcolor{gregoriocolor}{d} & \textcolor{gregoriocolor}{e} & \textcolor{gregoriocolor}{f} & \textcolor{gregoriocolor}{g} & \textcolor{gregoriocolor}{h} & \textcolor{gregoriocolor}{i} & \textcolor{gregoriocolor}{k} & \textcolor{gregoriocolor}{l} & \textcolor{gregoriocolor}{m} & \textcolor{gregoriocolor}{n} & \textcolor{gregoriocolor}{p} & \textcolor{gregoriocolor}{q} & \textcolor{gregoriocolor}{r} & \textcolor{gregoriocolor}{s} & \textcolor{gregoriocolor}{t} & \textcolor{gregoriocolor}{u} \\
 29 & 30 & 1 & 2 & 3 & 4 & 5 & 6 & 7 & 8 & 9 & 10 & 11 & 12 & 13 & 14 & 15 & 16 & 17 \\
\end{tabularx}
\vspace{0.5\baselineskip}
\noindent\begin{tabularx}{\linewidth}{*{12}{>{\centering\arraybackslash}X}}
 \textcolor{gregoriocolor}{A} & \textcolor{gregoriocolor}{B} & \textcolor{gregoriocolor}{C} & \textcolor{gregoriocolor}{D} & \textcolor{gregoriocolor}{E} & F & \textcolor{gregoriocolor}{F} & \textcolor{gregoriocolor}{G} & \textcolor{gregoriocolor}{H} & \textcolor{gregoriocolor}{M} & \textcolor{gregoriocolor}{N} & \textcolor{gregoriocolor}{P} \\
 18 & 19 & 20 & 21 & 22 & 23 & 23 & 24 & 25 & 26 & 27 & 28 \\
\end{tabularx}

\begin{paracol}{2}
\selectlanguage{latin}
\lettrine[lines=2]{I}{n} Æthiópia natális 
 sancti Matthæi, Apóstoli et Evangelístæ; qui, in ea regióne prædicans, 
 martyrium passus est. Hujus Evangélium, Hebræo sermóne conscríptum, 
 ipso Matthæo revelánte, invéntum est, una cum córpore beáti Bárnabæ Apóstoli, 
 témpore Zenónis Imperatóris.
\switchcolumn
\selectlanguage{english}
\lettrine[lines=2]{T}{he} birthday of St. Matthew, apostle 
 and evangelist, who suffered martyrdom in Ethiopia while engaged in 
 preaching. The Gospel written by him in Hebrew was found by his own 
 revelation during the time of Emperor Zeno, together with the relics of the 
 blessed apostle Barnabas.
\switchcolumn*
\selectlanguage{latin}
In terra Saar sancti 
 Jonæ Prophétæ, qui sepúltus est in Geth.
\switchcolumn
\selectlanguage{english}
In the land of the Saar, the holy 
 prophet Jonas, who was buried in Geth.
\switchcolumn*
\selectlanguage{latin}
In Æthiópia sanctæ 
 Iphigéniæ Vírginis, quæ, baptizáta a beáto Matthæo Apóstolo et Deo dicáta, 
 sancto fine quiévit.
\switchcolumn
\selectlanguage{english}
In Ethiopia, St. Iphigenia, virgin, 
 who was baptized and consecrated to God by the blessed apostle Matthew, and 
 who ended her holy life in peace.
\switchcolumn*
\selectlanguage{latin}
Romæ sancti Pámphili 
 Mártyris.
\switchcolumn
\selectlanguage{english}
At Rome, St Pamphilius, martyr.
\switchcolumn*
\selectlanguage{latin}
Eódem die, via Cláudia, vigésimo ab Urbe milliário, pássio sancti Alexándri Epíscopi, qui, sub 
 Antoníno Imperatóre, pro Christi fide, víncula, fustes, equúleum, lámpades 
 ardéntes, ungulárum laniatiónem, béstias ac fornácis superávit flammas, ac 
 tandem, gládio cæsus, vitam adéptus est gloriósam. Ejus corpus póstea 
 beátus Dámasus Papa in Urbem tránstulit sexto Kaléndas Decémbris.
\switchcolumn
\selectlanguage{english}
On the Claudian Way, twenty miles 
 from Rome, in the time of Emperor Antoninus, the martyrdom of St. Alexander, 
 bishop. For the faith of Christ he was loaded with fetters, scourged, 
 tortured, burned with torches, torn with iron hooks, exposed to the beasts, 
 and cast into the flames, but having overcome all these torments, he was 
 finally beheaded, and thus attained the glory of eternal life. His 
 body was afterwards carried into the city by blessed Pope Damasus on the 
 26th of November.
\switchcolumn*
\selectlanguage{latin}
In Cypro sancti Isácii, 
 Epíscopi et Mártyris.
\switchcolumn
\selectlanguage{english}
In Cyprus, St. Isacius, bishop and 
 martyr.
\switchcolumn*
\selectlanguage{latin}
In Phœnícia sancti 
 Eusébii Mártyris, qui, cum ultro Præféctum adísset et se Christiánum esse 
 denuntiásset, ab eo, multis torméntis afflíctus, cápite cæsus est.
\switchcolumn
\selectlanguage{english}
In Phoenicia, St. Eusebius, martyr, 
 who of his own accord went to the prefect and declared himself a Christian. 
 He was subjected by him to many torments, and finally beheaded.
\switchcolumn*
\selectlanguage{latin}
In Cypro sancti Melétii 
 Epíscopi et Confessóris.
\switchcolumn
\selectlanguage{english}
In Cyprus, St. Meletius, bishop and 
 confessor.
\switchcolumn*
\selectlanguage{latin}
\end{paracol}


% ---- martyrology/mart09/mart0922.htm
\needspace{10\baselineskip}
\begin{paracol}{2}
\selectlanguage{latin}
\begin{center}{\color{gregoriocolor} Décimo Kaléndas Octóbris. 
 Luna\dots\ }\end{center}
\switchcolumn
\selectlanguage{english}
\begin{center}{\color{gregoriocolor} The   22\textsuperscript{nd} Day of 
 September. The\dots\ Day of the Moon.}\end{center}
\end{paracol}

\noindent\begin{tabularx}{\linewidth}{*{19}{>{\centering\arraybackslash}X}}
 \textcolor{gregoriocolor}{a} & \textcolor{gregoriocolor}{b} & \textcolor{gregoriocolor}{c} & \textcolor{gregoriocolor}{d} & \textcolor{gregoriocolor}{e} & \textcolor{gregoriocolor}{f} & \textcolor{gregoriocolor}{g} & \textcolor{gregoriocolor}{h} & \textcolor{gregoriocolor}{i} & \textcolor{gregoriocolor}{k} & \textcolor{gregoriocolor}{l} & \textcolor{gregoriocolor}{m} & \textcolor{gregoriocolor}{n} & \textcolor{gregoriocolor}{p} & \textcolor{gregoriocolor}{q} & \textcolor{gregoriocolor}{r} & \textcolor{gregoriocolor}{s} & \textcolor{gregoriocolor}{t} & \textcolor{gregoriocolor}{u} \\
 30 & 1 & 2 & 3 & 4 & 5 & 6 & 7 & 8 & 9 & 10 & 11 & 12 & 13 & 14 & 15 & 16 & 17 & 18 \\
\end{tabularx}
\vspace{0.5\baselineskip}
\noindent\begin{tabularx}{\linewidth}{*{12}{>{\centering\arraybackslash}X}}
 \textcolor{gregoriocolor}{A} & \textcolor{gregoriocolor}{B} & \textcolor{gregoriocolor}{C} & \textcolor{gregoriocolor}{D} & \textcolor{gregoriocolor}{E} & F & \textcolor{gregoriocolor}{F} & \textcolor{gregoriocolor}{G} & \textcolor{gregoriocolor}{H} & \textcolor{gregoriocolor}{M} & \textcolor{gregoriocolor}{N} & \textcolor{gregoriocolor}{P} \\
 19 & 20 & 21 & 22 & 23 & 24 & 24 & 25 & 26 & 27 & 28 & 29 \\
\end{tabularx}

\begin{paracol}{2}
\selectlanguage{latin}
\lettrine[lines=2]{S}{ancti} Thomæ a Villa 
 Nova, ex Eremitárum sancti Augustíni Ordine, Epíscopi Valentíni et 
 Confessóris; cujus dies natális recólitur sexto Idus Septémbris.
\switchcolumn
\selectlanguage{english}
\lettrine[lines=2]{S}{t.} Thomas of Villanova, of the 
 Order of Hermits of St. Augustine, archbishop of Valencia and confessor, 
 whose birthday is the 8th of September.
\switchcolumn*
\selectlanguage{latin}
Sedúni, in Gállia, in 
 loco Agáuno, natális sanctórum Mártyrum Thebæórum Maurítii, Exsupérii, 
 Cándidi, Victóris, Innocéntii et Vitális, cum Sóciis ejúsdem legiónis; qui, 
 sub Maximiáno, pro Christo necáti, gloriósa passióne mundum illustrárunt.
\switchcolumn
\selectlanguage{english}
At St. Maurice, near Sion in 
 Switzerland, the birthday of the holy Theban martyrs Maurice, Exuperius, 
 Candidus, Victor, Innocent, and Vitalis, with their companions of the same 
 legion, whose martyrdom for the faith during the time of Maximian filled the 
 world with the glory of their sufferings.
\switchcolumn*
\selectlanguage{latin}
Romæ pássio sanctárum 
 Vírginum et Mártyrum Dignæ et Eméritæ, sub Valeriáno et Gallieno; quarum 
 relíquiæ in Ecclésia sancti Marcélli asservántur.
\switchcolumn
\selectlanguage{english}
At Rome, the martyrdom of the holy 
 virgins and martyrs Digna and Emerita, under Valerian and Gallienus. 
 Their relics are kept in the church of St. Marcellus.
\switchcolumn*
\selectlanguage{latin}
Ratisbónæ, in Bavária, sancti Emmerámi, Epíscopi et Mártyris; qui, ut álios liberáret, mortem 
 sævíssimam, Christi causa, patiénter súbiit.
\switchcolumn
\selectlanguage{english}
At Ratisbon in Bavaria, St. 
 Emmeramus, bishop and martyr, who patiently endured a most cruel death for 
 the sake of our Lord, in order to set others free.
\switchcolumn*
\selectlanguage{latin}
Apud pagum Castrénsium 
 sancti Jonæ, Presbyteri et Mártyris; qui, cum sancto Dionysio proféctus in 
 Gálliam, ibídem, Juliáni Præfécti jussu verbéribus cæsus, gládio martyrium 
 consummávit.
\switchcolumn
\selectlanguage{english}
At Arpajon, near Paris, St. Jonas, 
 priest and martyr, who went to France along with St. Denis. After he 
 was scourged by the order of the prefect Julian, his martyrdom was ended by 
 the sword.
\switchcolumn*
\selectlanguage{latin}
Antinoópoli, in Ægypto, 
 sanctæ Iráidis, Vírginis Alexandrínæ, et Sociórum Mártyrum. Ipsa 
 Virgo, cum esset ad hauriéndam e próximo fonte aquam egréssa, et navim 
 vidísset Confessóribus Christi onústam, prótinus, relícta hydria, se illis 
 adjúnxit, ac, simul cum iis in urbem ducta, prima ómnium, post multa 
 supplícia, cápite cæsa est; deínde Presbyteri, Diáconi, et Vírgines, aliíque 
 omnes eódem mortis génere consúmpti sunt.
\switchcolumn
\selectlanguage{english}
At Antinopolis in Egypt, the holy 
 martyrs Irais, an Alexandrian virgin, and her companions. Having gone 
 out to draw water at a near-by fountain, and seeing a boat loaded with 
 Christian confessors, she immediately left her vessel and joined them. 
 She was conducted to the city with them, and after many torments she was the 
 first to have her head struck off. After her, priests, deacons, 
 virgins, and all others underwent the same kind of death.
\switchcolumn*
\selectlanguage{latin}
Romæ sancti Felícis 
 Papæ Quarti, qui pro fide cathólica plúrimum laborávit.
\switchcolumn
\selectlanguage{english}
At Rome, Pope St. Felix IV, who 
 laboured exceedingly for the Catholic faith.
\switchcolumn*
\selectlanguage{latin}
Apud civitátem 
 Meldénsem beáti Sanctíni Epíscopi, qui fuit discípulus sancti Dionysii 
 Areopaíitæ; et ejúsdem civitátis Epíscopus ab eo consecrátus, primus illic 
 Evangélium prædicávit.
\switchcolumn
\selectlanguage{english}
At Meaux, blessed Sanctinus, bishop, 
 a disciple of St. Denis the Areopagite, by whom he was consecrated bishop of 
 that city, and was the first to preach the Gospel there.
\switchcolumn*
\selectlanguage{latin}
In território 
 Constantiénsi, in Gállia, sancti Lautónis Epíscopi.
\switchcolumn
\selectlanguage{english}
In the territory of Coutances, St. 
 Lauto, bishop.
\switchcolumn*
\selectlanguage{latin}
In monte Glonna, ad 
 Lígerim flumen, in Gállia, sancti Floréntii Presbyteri.
\switchcolumn
\selectlanguage{english}
At Mount Glonna in France, the holy 
 priest Florentius.
\switchcolumn*
\selectlanguage{latin}
In óppido cui Leprósii 
 nomen, in território Bituricénsi, sancti Silváni Confessóris.
\switchcolumn
\selectlanguage{english}
In the territory of Bourges, St. 
 Sylvanus, confessor.
\switchcolumn*
\selectlanguage{latin}
Laudúni, in Gállia, 
 sanctæ Salabérgæ Abbatíssæ.
\switchcolumn
\selectlanguage{english}
At Laon in France, St. Salaberga, 
 abbess.
\switchcolumn*
\selectlanguage{latin}
\end{paracol}


% ---- martyrology/mart09/mart0923.htm
\needspace{10\baselineskip}
\begin{paracol}{2}
\selectlanguage{latin}
\begin{center}{\color{gregoriocolor} Nono Kaléndas Octóbris. 
 Luna\dots\ }\end{center}
\switchcolumn
\selectlanguage{english}
\begin{center}{\color{gregoriocolor} The   23\textsuperscript{rd} Day of 
 September. The\dots\ Day of the Moon.}\end{center}
\end{paracol}

\noindent\begin{tabularx}{\linewidth}{*{19}{>{\centering\arraybackslash}X}}
 \textcolor{gregoriocolor}{a} & \textcolor{gregoriocolor}{b} & \textcolor{gregoriocolor}{c} & \textcolor{gregoriocolor}{d} & \textcolor{gregoriocolor}{e} & \textcolor{gregoriocolor}{f} & \textcolor{gregoriocolor}{g} & \textcolor{gregoriocolor}{h} & \textcolor{gregoriocolor}{i} & \textcolor{gregoriocolor}{k} & \textcolor{gregoriocolor}{l} & \textcolor{gregoriocolor}{m} & \textcolor{gregoriocolor}{n} & \textcolor{gregoriocolor}{p} & \textcolor{gregoriocolor}{q} & \textcolor{gregoriocolor}{r} & \textcolor{gregoriocolor}{s} & \textcolor{gregoriocolor}{t} & \textcolor{gregoriocolor}{u} \\
 1 & 2 & 3 & 4 & 5 & 6 & 7 & 8 & 9 & 10 & 11 & 12 & 13 & 14 & 15 & 16 & 17 & 18 & 19 \\
\end{tabularx}
\vspace{0.5\baselineskip}
\noindent\begin{tabularx}{\linewidth}{*{12}{>{\centering\arraybackslash}X}}
 \textcolor{gregoriocolor}{A} & \textcolor{gregoriocolor}{B} & \textcolor{gregoriocolor}{C} & \textcolor{gregoriocolor}{D} & \textcolor{gregoriocolor}{E} & F & \textcolor{gregoriocolor}{F} & \textcolor{gregoriocolor}{G} & \textcolor{gregoriocolor}{H} & \textcolor{gregoriocolor}{M} & \textcolor{gregoriocolor}{N} & \textcolor{gregoriocolor}{P} \\
 20 & 21 & 22 & 23 & 24 & 25 & 25 & 26 & 27 & 28 & 29 & 30 \\
\end{tabularx}

\begin{paracol}{2}
\selectlanguage{latin}
\lettrine[lines=2]{R}{omæ} sancti Lini, Papæ 
 et Mártyris, qui, primus post beátum Petrum Apóstolum, Románum Ecclésiam 
 gubernávit, et, martyrio coronátus, sepúltus est in Vaticáno, prope eúndem 
 Apóstolum.
\switchcolumn
\selectlanguage{english}
\lettrine[lines=2]{A}{t} Rome, St. Linus, pope and martyr, 
 who governed the Roman Church next after the blessed apostle Peter. He 
 was crowned with martyrdom and was buried on the Vatican Hill beside the 
 same apostle.
\switchcolumn*
\selectlanguage{latin}
Icónii, in Lycaónia, 
 sanctæ Theclæ, Vírginis et Mártyris; quæ, a sancto Paulo Apóstolo ad fidem 
 perdúcta, ignes ac béstias, sub Neróne Imperatóre, in Christi confessióne 
 devícit; et, post plúrima ad multórum doctrínam superáta certámina, venit 
 Seleucíam, ibíque requiévit in pace. Ipsam vero sancti Patres summis 
 láudibus celebrárunt.
\switchcolumn
\selectlanguage{english}
At Iconium in Lycaonia, St. Thecla, 
 virgin and martyr, who was brought to the faith by the apostle St. Paul. 
 Under Emperor Nero, she was victorious over the flames and the beasts to 
 which she was exposed for the faith of Christ. After many combats 
 endured for the instruction of others, she went to Seleucia, where she ended 
 her days in peace. Her memory has been eulogized by the holy Fathers.
\switchcolumn*
\selectlanguage{latin}
In Hispánia sanctárum 
 mulíerum Xantíppæ et Polyxenæ, quæ fuérunt Apostolórum discípulæ.
\switchcolumn
\selectlanguage{english}
In Spain, the holy women Xantippa 
 and Polyxena, who were disciples of the apostles.
\switchcolumn*
\selectlanguage{latin}
In Africa sanctórum 
 Mártyrum Andréæ, Joánnis, Petri et Antónii.
\switchcolumn
\selectlanguage{english}
In Africa, the holy martyrs Andrew, 
 John, Peter and Anthony.
\switchcolumn*
\selectlanguage{latin}
Ancónæ sancti 
 Constántii, Ecclésiæ Mansionárii, miraculórum grátia conspícui.
\switchcolumn
\selectlanguage{english}
At Ancona, St. Constantius, 
 sacristan of the church, renowned for the gift of miracles.
\switchcolumn*
\selectlanguage{latin}
In Campánia 
 commemorátio beáti Sósii, Diáconi, Misenátis, quem sanctus Epíscopus 
 Januárius, cum de cápite illíus, Evangélium in Ecclésia legéntis, flammam 
 vidéret exsúrgere, Mártyrem futúrum præannuntiávit; et, non post multos 
 dies, ipse Sósius, cum esset annórum trigínta, cápitis detruncatióne 
 martyrium cum eódem Epíscopo suscépit.
\switchcolumn
\selectlanguage{english}
In Campania, the commemoration of 
 the blessed Sosius, deacon of the church of Miseno. The holy bishop 
 Januarius, upon seeing a flame arise from his head as he was reading the 
 Gospel in the church, foretold that he would be a martyr. Not many 
 days after, when he was thirty years of age, he and the holy bishop suffered 
 martyrdom by beheading.
\switchcolumn*
\selectlanguage{latin}
Sescíaci, in 
 Constantiénsi Gálliæ território, item Commemorátio sancti Patérni, Epíscopi 
 Abrincénsis et Confessóris; cujus dies natális sextodécimo Kaléndas Maji 
 recólitur.
\switchcolumn
\selectlanguage{english}
At Scicy in the district of 
 Coutances in France, the commemoration of St. Paternus, bishop of Avranches 
 and confessor, whose birthday is recalled on the 16th of April.
\switchcolumn*
\selectlanguage{latin}
\end{paracol}


% ---- martyrology/mart09/mart0924.htm
\needspace{10\baselineskip}
\begin{paracol}{2}
\selectlanguage{latin}
\begin{center}{\color{gregoriocolor} Octávo Kaléndas Octóbris. 
 Luna\dots\ }\end{center}
\switchcolumn
\selectlanguage{english}
\begin{center}{\color{gregoriocolor} The   24\textsuperscript{th} Day of 
 September. The\dots\ Day of the Moon.}\end{center}
\end{paracol}

\noindent\begin{tabularx}{\linewidth}{*{19}{>{\centering\arraybackslash}X}}
 \textcolor{gregoriocolor}{a} & \textcolor{gregoriocolor}{b} & \textcolor{gregoriocolor}{c} & \textcolor{gregoriocolor}{d} & \textcolor{gregoriocolor}{e} & \textcolor{gregoriocolor}{f} & \textcolor{gregoriocolor}{g} & \textcolor{gregoriocolor}{h} & \textcolor{gregoriocolor}{i} & \textcolor{gregoriocolor}{k} & \textcolor{gregoriocolor}{l} & \textcolor{gregoriocolor}{m} & \textcolor{gregoriocolor}{n} & \textcolor{gregoriocolor}{p} & \textcolor{gregoriocolor}{q} & \textcolor{gregoriocolor}{r} & \textcolor{gregoriocolor}{s} & \textcolor{gregoriocolor}{t} & \textcolor{gregoriocolor}{u} \\
 2 & 3 & 4 & 5 & 6 & 7 & 8 & 9 & 10 & 11 & 12 & 13 & 14 & 15 & 16 & 17 & 18 & 19 & 20 \\
\end{tabularx}
\vspace{0.5\baselineskip}
\noindent\begin{tabularx}{\linewidth}{*{12}{>{\centering\arraybackslash}X}}
 \textcolor{gregoriocolor}{A} & \textcolor{gregoriocolor}{B} & \textcolor{gregoriocolor}{C} & \textcolor{gregoriocolor}{D} & \textcolor{gregoriocolor}{E} & F & \textcolor{gregoriocolor}{F} & \textcolor{gregoriocolor}{G} & \textcolor{gregoriocolor}{H} & \textcolor{gregoriocolor}{M} & \textcolor{gregoriocolor}{N} & \textcolor{gregoriocolor}{P} \\
 21 & 22 & 23 & 24 & 25 & 26 & 26 & 27 & 28 & 29 & 30 & 1 \\
\end{tabularx}

\begin{paracol}{2}
\selectlanguage{latin}
\lettrine[lines=2]{F}{estum} beátæ Maríæ 
 Vírginis de Mercéde nuncupátæ, Ordinis redemptiónis captivórum sub ejus 
 nómine Institutrícis, de cujus Apparitióne ágitur quarto Idus Augústi.
\switchcolumn
\selectlanguage{english}
\lettrine[lines=2]{T}{he} feast of our Lady of Ransom, 
 Foundress of the Order for the Redemption of Captives. The apparition 
 of the same Blessed Virgin occurred on the 10th of August.
\switchcolumn*
\selectlanguage{latin}
Bríxiæ deposítio sancti 
 Anathalónis Epíscopi, qui, beáti Bárnabæ Apóstoli discípulus, in ejus locum 
 Epíscopus Ecclésiæ Mediolanénsis succéssit.
\switchcolumn
\selectlanguage{english}
At Brescia, the death of St. 
 Anathalo, bishop. He was a disciple of the blessed apostle Barnabas, 
 and succeeded him as bishop of the Milanese church.
\switchcolumn*
\selectlanguage{latin}
In Pannónia sancti 
 Gerárdi, Epíscopi Morisénæ sedis et Mártyris, Hungarórum Apóstoli nuncupáti, 
 patrícii Véneti; qui, cum e Chanadiénsi óppido Albam Regálem se conférret, 
 prope flumen Danúbium ab infidélibus impetitus, lapídibus óbrutus ac tandem 
 láncea transfíxus occúbuit, sicque primus pátriam nóbili martyrio 
 illustrávit.
\switchcolumn
\selectlanguage{english}
In Hungary, St. Gerard, bishop of 
 Chzonad and martyr, patrician of Venice, called the apostle of the 
 Hungarians. During a journey from the town of Chzonad to Alba Regalis 
 he was attacked by the pagans near the river Danube, stoned by them, and 
 then pierced with a lance. He was thus the first to adorn his native 
 land with a noble martyrdom.
\switchcolumn*
\selectlanguage{latin}
Augustodúni natális 
 sanctórum Mártyrum Andóchii Presbyteri, Thyrsi Diáconi, et Felícis. 
 Hi, a beáto Polycárpo, Smyrnénsi Epíscopo, ab Oriénte dirécti ad docéndam 
 Gálliam, ibídem flagéllis duríssime cæsi, ac tota die invérsis mánibus 
 suspénsi, et in ignem missi sunt, sed non combústi; tandem eórum colla 
 véctibus feriúntur, et ita Mártyres gloriosíssime coronántur.
\switchcolumn
\selectlanguage{english}
At Autun, the birthday of the holy 
 martyrs Andochius, a priest, Thyrsus, a deacon, and Felix, who were sent 
 from the East by blessed Polycarp, bishop of Smyrna, to preach in France. 
 There they were severely scourged, hanged by the hands for a whole day, and 
 cast into the fire. Remaining uninjured, they had their necks broken 
 with heavy bars, and thus won a most glorious crown.
\switchcolumn*
\selectlanguage{latin}
In Ægypto pássio 
 sanctórum Paphnútii et Sociórum Mártyrum. Ipse, vitam in solitúdine 
 agens, cum audíret multos Christiános in vínculis retinéri, sponte, divíno 
 Spíritu cóncitus, Præfécto se offert, et Christiánam religiónem líbere 
 profitétur; a quo primum caténis férreis constríngitur, et in equúleo 
 diutíssime torquétur, deínde cum áliis plúrimis ad Diocletiánum míttitur, 
 cujus jussu, ipse palmæ affígitur, céteri autem ferro necántur.
\switchcolumn
\selectlanguage{english}
In Egypt, the holy martyrs 
 Paphnutius and his companions. While leading a solitary life, St. 
 Paphnutius heard that many Christians were kept in bonds. Moved by the 
 spirit of God, he voluntarily offered himself to the prefect, and freely 
 confessed the Christian faith. He was bound by him with iron chains, 
 and for a long time tortured on the rack. Then, being sent with many 
 others to Diocletian, by his order he was fastened to a palm tree, and the 
 rest were struck with the sword.
\switchcolumn*
\selectlanguage{latin}
Chalcédone sanctórum 
 quadragínta novem Mártyrum, qui, post martyrium sanctæ Euphémiæ, sub 
 Diocletiáno Imperatóre, damnáti ad béstias, et, cum ab iis divínitus líberi 
 evasíssent, demum, gládio percússi, migravérunt in cælum.
\switchcolumn
\selectlanguage{english}
At Chalcedon, under Emperor 
 Diocletian, after the martyrdom of St. Euphemia, forty-nine holy martyrs who 
 were condemned to be devoured by the beasts, but being miraculously 
 delivered, were finally struck with the sword and went to heaven.
\switchcolumn*
\selectlanguage{latin}
Arvérnis, in Gállia, 
 deposítio sancti Rústici, Epíscopi et Confessóris.
\switchcolumn
\selectlanguage{english}
In Auvergne, the death of St. 
 Rusticus, bishop and confessor.
\switchcolumn*
\selectlanguage{latin}
Flavíaci, in território Bellovacénsi, sancti Geremari, Presbyteri et Abbátis.
\switchcolumn
\selectlanguage{english}
In the diocese of Beauvais, St. 
 Geremarus, priest and abbot.
\switchcolumn*
\selectlanguage{latin}
Septémpedæ, in Picéno, 
 deposítio sancti Pacífici, Sacerdótis ex Ordine Minórum et Confessóris, 
 exímiæ patiéntiæ viri et solitúdinis amóre præclári, quem Gregórius Papa 
 Décimus sextus in Sanctórum cánonem rétulit.
\switchcolumn
\selectlanguage{english}
At San Severino in Piceno, the death 
 of St. Pacificus, priest and confessor of the Order of Friars Minor of St. 
 Francis of the Reformed Observance. Illustrious for his great patience 
 and his love of solitude, he was enrolled in the canon of the saints by Pope 
 Gregory XVI.
\switchcolumn*
\selectlanguage{latin}
\end{paracol}


% ---- martyrology/mart09/mart0925.htm
\needspace{10\baselineskip}
\begin{paracol}{2}
\selectlanguage{latin}
\begin{center}{\color{gregoriocolor} Séptimo Kaléndas Octóbris. 
 Luna\dots\ }\end{center}
\switchcolumn
\selectlanguage{english}
\begin{center}{\color{gregoriocolor} The   25\textsuperscript{th} Day of 
 September. The\dots\ Day of the Moon.}\end{center}
\end{paracol}

\noindent\begin{tabularx}{\linewidth}{*{19}{>{\centering\arraybackslash}X}}
 \textcolor{gregoriocolor}{a} & \textcolor{gregoriocolor}{b} & \textcolor{gregoriocolor}{c} & \textcolor{gregoriocolor}{d} & \textcolor{gregoriocolor}{e} & \textcolor{gregoriocolor}{f} & \textcolor{gregoriocolor}{g} & \textcolor{gregoriocolor}{h} & \textcolor{gregoriocolor}{i} & \textcolor{gregoriocolor}{k} & \textcolor{gregoriocolor}{l} & \textcolor{gregoriocolor}{m} & \textcolor{gregoriocolor}{n} & \textcolor{gregoriocolor}{p} & \textcolor{gregoriocolor}{q} & \textcolor{gregoriocolor}{r} & \textcolor{gregoriocolor}{s} & \textcolor{gregoriocolor}{t} & \textcolor{gregoriocolor}{u} \\
 3 & 4 & 5 & 6 & 7 & 8 & 9 & 10 & 11 & 12 & 13 & 14 & 15 & 16 & 17 & 18 & 19 & 20 & 21 \\
\end{tabularx}
\vspace{0.5\baselineskip}
\noindent\begin{tabularx}{\linewidth}{*{12}{>{\centering\arraybackslash}X}}
 \textcolor{gregoriocolor}{A} & \textcolor{gregoriocolor}{B} & \textcolor{gregoriocolor}{C} & \textcolor{gregoriocolor}{D} & \textcolor{gregoriocolor}{E} & F & \textcolor{gregoriocolor}{F} & \textcolor{gregoriocolor}{G} & \textcolor{gregoriocolor}{H} & \textcolor{gregoriocolor}{M} & \textcolor{gregoriocolor}{N} & \textcolor{gregoriocolor}{P} \\
 22 & 23 & 24 & 25 & 26 & 27 & 27 & 28 & 29 & 30 & 1 & 2 \\
\end{tabularx}

\begin{paracol}{2}
\selectlanguage{latin}
\lettrine[lines=2]{A}{pud} castéllum Emmaus 
 natális beáti Cléophæ, qui fuit Christi discípulus, quem et in eádem domo in 
 qua mensam Dómino paráverat, pro confessióne illíus a Judæis occísum tradunt, 
 et gloriósa memória sepúltum.
\switchcolumn
\selectlanguage{english}
\lettrine[lines=2]{A}{t} Emmaus, the birthday of blessed 
 Cleophas, disciple of Christ. It is related that he was killed by the 
 Jews for the confession of our Lord, and honourably buried in the same house 
 in which he had entertained him.
\switchcolumn*
\selectlanguage{latin}
Ambiáni, in Gállia, 
 beáti Firmíni Epíscopi, qui, in persecutióne Diocletiáni, sub Rictiováro 
 Præside, post vária torménta, cápitis decollatióne martyrium sumpsit.
\switchcolumn
\selectlanguage{english}
At Amiens in France, in the 
 persecution of Diocletian, blessed Firminus, bishop. Under the 
 governor Rictiovarus, after many torments he suffered martyrdom by being 
 beheaded.
\switchcolumn*
\selectlanguage{latin}
Eódem die, via Cláudia, sancti Herculáni, mílitis et Mártyris; qui, sub Antoníno Imperatóre, 
 miráculis in passióne beáti Alexándri Epíscopi ad Christum convérsus, atque 
 ob fídei confessiónem, post multa torménta, gládio cæsus est.
\switchcolumn
\selectlanguage{english}
At Rome, on the Claudian Way, under 
 Emperor Antoninus, St. Herculanus, soldier and martyr, who was converted to 
 Christ by the miracle wrought during the martyrdom of the blessed bishop 
 Alexander. After enduring many torments he was put to the sword.
\switchcolumn*
\selectlanguage{latin}
Damásci sanctórum 
 Mártyrum Pauli, et Tattæ cónjugis, ac Sabiniáni, Máximi, Rufi et Eugénii 
 filiórum; qui, Christiánæ religiónis accusáti, verbéribus aliísque 
 supplíciis torti sunt, et in cruciátibus ánimas Deo reddidérunt.
\switchcolumn
\selectlanguage{english}
At Damascus, the holy martyrs Paul, 
 his wife Tatta, and their sons Sabinian, Maximus, Rufus, and Eugene. 
 Accused of professing the Christian religion, they were scourged and 
 tortured in other ways until they gave up their souls unto God.
\switchcolumn*
\selectlanguage{latin}
In Asia pássio 
 sanctórum Bardomiáni, Eucárpi et aliórum vigínti sex Mártyrum.
\switchcolumn
\selectlanguage{english}
In Asia, the holy martyrs Bardomian, 
 Eucarpus, and twenty-six others.
\switchcolumn*
\selectlanguage{latin}
Lugdúni, in Gállia, 
 deposítio sancti Lupi, qui ex Anachoréta factus est Epíscopus.
\switchcolumn
\selectlanguage{english}
At Lyons in France, the death of St. 
 Lupus, at one time an anchoret, but later a bishop.
\switchcolumn*
\selectlanguage{latin}
Antisiodóri sancti 
 Anachárii, Epíscopi et Confessóris.
\switchcolumn
\selectlanguage{english}
At Auxerre, St. Anacharius, bishop 
 and confessor.
\switchcolumn*
\selectlanguage{latin}
Blesis, in Gállia, 
 sancti Solémnii, Epíscopus Carnuténsis, miráculis clari.
\switchcolumn
\selectlanguage{english}
At Blois in France, St. Solemnius, 
 bishop of Chartres, renowned for miracles.
\switchcolumn*
\selectlanguage{latin}
Eódem die sancti 
 Princípii, qui fuit Epíscopus Suessionénsis et frater beáti Remígii Epíscopi.
\switchcolumn
\selectlanguage{english}
On the same day, St. Principius, 
 bishop of Soissons, brother of the blessed bishop Remigius.
\switchcolumn*
\selectlanguage{latin}
Anágniæ sanctárum 
 Vírginum Auréliæ et Neomísiæ.
\switchcolumn
\selectlanguage{english}
At Anagni, the holy virgins Aurelia 
 and Neomysia.
\switchcolumn*
\selectlanguage{latin}
\end{paracol}


% ---- martyrology/mart09/mart0926.htm
\needspace{10\baselineskip}
\begin{paracol}{2}
\selectlanguage{latin}
\begin{center}{\color{gregoriocolor} Sexto Kaléndas Octóbris. 
 Luna\dots\ }\end{center}
\switchcolumn
\selectlanguage{english}
\begin{center}{\color{gregoriocolor} The   26\textsuperscript{th} Day of 
 September. The\dots\ Day of the Moon.}\end{center}
\end{paracol}

\noindent\begin{tabularx}{\linewidth}{*{19}{>{\centering\arraybackslash}X}}
 \textcolor{gregoriocolor}{a} & \textcolor{gregoriocolor}{b} & \textcolor{gregoriocolor}{c} & \textcolor{gregoriocolor}{d} & \textcolor{gregoriocolor}{e} & \textcolor{gregoriocolor}{f} & \textcolor{gregoriocolor}{g} & \textcolor{gregoriocolor}{h} & \textcolor{gregoriocolor}{i} & \textcolor{gregoriocolor}{k} & \textcolor{gregoriocolor}{l} & \textcolor{gregoriocolor}{m} & \textcolor{gregoriocolor}{n} & \textcolor{gregoriocolor}{p} & \textcolor{gregoriocolor}{q} & \textcolor{gregoriocolor}{r} & \textcolor{gregoriocolor}{s} & \textcolor{gregoriocolor}{t} & \textcolor{gregoriocolor}{u} \\
 4 & 5 & 6 & 7 & 8 & 9 & 10 & 11 & 12 & 13 & 14 & 15 & 16 & 17 & 18 & 19 & 20 & 21 & 22 \\
\end{tabularx}
\vspace{0.5\baselineskip}
\noindent\begin{tabularx}{\linewidth}{*{12}{>{\centering\arraybackslash}X}}
 \textcolor{gregoriocolor}{A} & \textcolor{gregoriocolor}{B} & \textcolor{gregoriocolor}{C} & \textcolor{gregoriocolor}{D} & \textcolor{gregoriocolor}{E} & F & \textcolor{gregoriocolor}{F} & \textcolor{gregoriocolor}{G} & \textcolor{gregoriocolor}{H} & \textcolor{gregoriocolor}{M} & \textcolor{gregoriocolor}{N} & \textcolor{gregoriocolor}{P} \\
 23 & 24 & 25 & 26 & 27 & 28 & 28 & 29 & 30 & 1 & 2 & 3 \\
\end{tabularx}

\begin{paracol}{2}
\selectlanguage{latin}
\lettrine[lines=2]{N}{icomedíæ} natális 
 sanctórum Mártyrum Cypriáni et Justínæ Vírginis. Hæc, sub Diocletiáno 
 Imperatóre et Eutólmio Præside, cum multa pro Christo pertulísset, ipsum 
 quoque Cypriánum, qui erat magnus et suis mágicis ártibus eam dementáre 
 conabátur, ad Christiánam fidem convértit; cum quo póstea martyrium sumpsit. 
 Eórum córpora, feris objécta, rapuérunt noctu quidam nautæ Christiáni, et 
 Romam detulérunt; quæ, póstmodum in Basílicam Constantiniánam transláta, 
 prope Baptistérium cóndita sunt.
\switchcolumn
\selectlanguage{english}
\lettrine[lines=2]{A}{t} Nicomedia, the birthday of the 
 holy martyrs Cyprian and the virgin Justina. Under Emperor Diocletian 
 and the governor Eutholmius, Justina suffered greatly for the faith of 
 Christ, and thus converted Cyprian, who, while a magician, had endeavoured 
 to bring her under the influence of his magical practices. She 
 afterwards suffered martyrdom with him. Their bodies were exposed to 
 the beasts, but were taken away in the night by some Christian sailors, and 
 carried to Rome. They were subsequently taken into the Constantinian 
 basilica, and buried near the baptistry.
\switchcolumn*
\selectlanguage{latin}
Romæ sancti Callístrati 
 Mártyris, et aliórum quadragínta novem mílitum; qui mílites, in persecutióne 
 Diocletiáni Imperatóris, cum Callístratus, insútus cúleo et in mare demérsus, 
 divína ope evasísset incólumis, ad Christiánam religiónem convérsi sunt, et 
 cum eo páriter martyrium subiérunt.
\switchcolumn
\selectlanguage{english}
At Rome, in the persecution of 
 Diocletian, the holy martyr Callistratus and forty-nine other soldiers who 
 endured martyrdom together. The companions of Callistratus were 
 converted to Christ upon seeing him miraculously delivered from drowning in 
 the sea, although he had been sewn up in a bag and thrown in.
\switchcolumn*
\selectlanguage{latin}
Bonóniæ sancti Eusébii, 
 Epíscopi et Confessóris.
\switchcolumn
\selectlanguage{english}
At Bologna, St. Eusebius, bishop and 
 confessor.
\switchcolumn*
\selectlanguage{latin}
Bríxiæ sancti Vigílii 
 Epíscopi.
\switchcolumn
\selectlanguage{english}
At Brescia, St. Vigilius, bishop.
\switchcolumn*
\selectlanguage{latin}
In agro Tusculáno beáti 
 Nili Abbátis, qui fundátor monastérii Cryptæ Ferrátæ ac vir magnæ sanctitátis 
 éxstitit.
\switchcolumn
\selectlanguage{english}
In the Tuscan plain, the blessed 
 Abbot Nilus, founder of the monastery of Grottaferrata, a man of eminent 
 sanctity.
\switchcolumn*
\selectlanguage{latin}
Tiférni, in Umbria, 
 sancti Amántii Presbyteri, virtúte miraculórum illústris.
\switchcolumn
\selectlanguage{english}
At Tiferno in Umbria, St. Amantius, 
 a priest distinguished for the gift of miracles.
\switchcolumn*
\selectlanguage{latin}
Albáni sancti Senatóris.
\switchcolumn
\selectlanguage{english}
At Albano, St. Senatore.
\switchcolumn*
\selectlanguage{latin}
\end{paracol}


% ---- martyrology/mart09/mart0927.htm
\needspace{10\baselineskip}
\begin{paracol}{2}
\selectlanguage{latin}
\begin{center}{\color{gregoriocolor} Quinto Kaléndas Octóbris. 
 Luna\dots\ }\end{center}
\switchcolumn
\selectlanguage{english}
\begin{center}{\color{gregoriocolor} The   27\textsuperscript{th} Day of 
 September. The\dots\ Day of the Moon.}\end{center}
\end{paracol}

\noindent\begin{tabularx}{\linewidth}{*{19}{>{\centering\arraybackslash}X}}
 \textcolor{gregoriocolor}{a} & \textcolor{gregoriocolor}{b} & \textcolor{gregoriocolor}{c} & \textcolor{gregoriocolor}{d} & \textcolor{gregoriocolor}{e} & \textcolor{gregoriocolor}{f} & \textcolor{gregoriocolor}{g} & \textcolor{gregoriocolor}{h} & \textcolor{gregoriocolor}{i} & \textcolor{gregoriocolor}{k} & \textcolor{gregoriocolor}{l} & \textcolor{gregoriocolor}{m} & \textcolor{gregoriocolor}{n} & \textcolor{gregoriocolor}{p} & \textcolor{gregoriocolor}{q} & \textcolor{gregoriocolor}{r} & \textcolor{gregoriocolor}{s} & \textcolor{gregoriocolor}{t} & \textcolor{gregoriocolor}{u} \\
 5 & 6 & 7 & 8 & 9 & 10 & 11 & 12 & 13 & 14 & 15 & 16 & 17 & 18 & 19 & 20 & 21 & 22 & 23 \\
\end{tabularx}
\vspace{0.5\baselineskip}
\noindent\begin{tabularx}{\linewidth}{*{12}{>{\centering\arraybackslash}X}}
 \textcolor{gregoriocolor}{A} & \textcolor{gregoriocolor}{B} & \textcolor{gregoriocolor}{C} & \textcolor{gregoriocolor}{D} & \textcolor{gregoriocolor}{E} & F & \textcolor{gregoriocolor}{F} & \textcolor{gregoriocolor}{G} & \textcolor{gregoriocolor}{H} & \textcolor{gregoriocolor}{M} & \textcolor{gregoriocolor}{N} & \textcolor{gregoriocolor}{P} \\
 24 & 25 & 26 & 27 & 28 & 29 & 29 & 30 & 1 & 2 & 3 & 4 \\
\end{tabularx}

\begin{paracol}{2}
\selectlanguage{latin}
Ægéæ natális sanctórum 
 Mártyrum Cosmæ et Damiáni fratrum, qui, in persecutióne Diocletiáni, post 
 multa torménta, víncula et cárceres, post mare et ignes, cruces, 
 lapidatiónem et sagíttas divínitus superátas, cápite plectúntur; cum quibus 
 étiam referúntur passi tres eórum fratres germáni, id est Anthimus, Leóntius 
 et Euprépius.
\switchcolumn
\selectlanguage{english}
\lettrine[lines=2]{A}{t} Aegea, during the persecution of 
 Diocletian, the birthday of the holy martyrs Cosmas and Damian, brothers. 
 After miraculously overcoming many torments from bonds, imprisonment, fire, 
 crucifixion, stoning, arrows, and from being cast into the sea, they were 
 beheaded. With them are said to have suffered three brothers: Anthimus, 
 Leontius, and Euprepius.
\switchcolumn*
\selectlanguage{latin}
Lutétiæ Parisiórum item 
 natális sancti Vincéntii a Paulo, Presbyteri et Confessóris, Congregatiónis 
 Presbyterórum Missiónis et Puellárum Caritátis Fundatóris, viri apostólici 
 et páuperum patris; quem Leo Décimus tértius, Póntifex Máximus, ómnium 
 Societátum caritátis, in toto cathólico Orbe exsisténtium et ab eódem Sancto 
 quomodólibet promanántium, cæléstem Patrónum apud Deum constítuit. 
 Ipsíus tamen festívitas quartodécimo Kaléndas Augústi celebrátur.
\switchcolumn
\selectlanguage{english}
At Paris, the birthday of St. 
 Vincent de Paul, priest and confessor, founder of the Congregation of the 
 Mission and of the Sisters of Charity, an apostolic man and father to the 
 poor. Pope Leo XIII appointed this saint as the heavenly patron before 
 God of all charitable societies in the world which in any way whatever draw 
 their origin from him. His feast is celebrated on the 19th of July.
\switchcolumn*
\selectlanguage{latin}
Bybli, in Phœnícia, 
 sancti Marci Epíscopi, qui et Joánnes a beáto Luca nominátur, atque fílius éxstitit illíus beátæ Maríæ, cujus memória tértio Kaléndas Júlii recensétur.
\switchcolumn
\selectlanguage{english}
At Byblos in Phoenicia, Bishop St. 
 Mark, whom St. Luke calls John, and who was the son of that blessed Mary who 
 is commemorated on the 29th of July.
\switchcolumn*
\selectlanguage{latin}
Medioláni sancti Caji 
 Epíscopi, qui fuit discípulus beáti Bárnabæ Apóstoli; et, multa in Nerónis 
 persecutióne passus, quiévit in pace.
\switchcolumn
\selectlanguage{english}
At Milan, the holy bishop Caius, a 
 disciple of the blessed apostle Barnabas, who passed calmly to rest after 
 suffering severely in the persecution of Nero.
\switchcolumn*
\selectlanguage{latin}
Romæ sanctæ Epicháridis, 
 mulíeris Senatóriæ, quæ, in eádem Diocletiáni persecutióne, plumbátis cæsa 
 est atque gládio percússa.
\switchcolumn
\selectlanguage{english}
At Rome, St. Epicharis, wife of a 
 senator, who was scourged with leaded whips and then struck with the sword 
 in the persecution of Diocletian.
\switchcolumn*
\selectlanguage{latin}
Tudérti, in Umbria, 
 sanctórum Mártyrum Fidéntii et Teréntii, sub eódem Diocletiáno.
\switchcolumn
\selectlanguage{english}
At Todi in Umbria, under the same 
 Diocletian, the holy martyrs Fidentius and Terence.
\switchcolumn*
\selectlanguage{latin}
Córdubæ, in Hispánia, 
 sanctórum Mártyrum Adúlfi et Joánnis fratrum, qui, in persecutióne Arábica, 
 pro Christo coronáti sunt; eorúmque animáta exémplo beáta Virgo Aurea, 
 ipsórum soror, ad fidem redúcta, et ipsa póstmodum martyrium fórtiter súbiit 
 quartodécimo Kaléndas Augústi.
\switchcolumn
\selectlanguage{english}
At Cordova in Spain, the holy 
 martyrs Adolph and John, brothers, who won the martyrs' crown in the 
 Arabian persecution. Their sister, the blessed virgin Aurea, was 
 inspired by their example to return to the faith and later bravely suffered 
 martyrdom on the 19th of July.
\switchcolumn*
\selectlanguage{latin}
Sedúni, in Gállia, 
 sancti Florentíni Mártyris, qui, una cum beáto Hilário, post abscissiónem 
 linguæ, jussus est gládio feríri.
\switchcolumn
\selectlanguage{english}
At Sion in Switzerland, St. 
 Florentius, martyr. After his tongue had been cut out, he was put to 
 the sword with blessed Hilary.
\switchcolumn*
\selectlanguage{latin}
Ravénnæ sancti Aderíti, 
 Epíscopi et Confessóris.
\switchcolumn
\selectlanguage{english}
At Ravenna, St. Aderitus, bishop and 
 confessor.
\switchcolumn*
\selectlanguage{latin}
Lutétiæ Parisiórum 
 sancti Elzeárii Cómitis.
\switchcolumn
\selectlanguage{english}
At Paris, St. Eleazar, a count.
\switchcolumn*
\selectlanguage{latin}
In Hannónia sanctæ 
 Hiltrúdis Vírginis.
\switchcolumn
\selectlanguage{english}
In Hainault, St. Hiltrude, virgin.
\switchcolumn*
\selectlanguage{latin}
\end{paracol}


% ---- martyrology/mart09/mart0928.htm
\needspace{10\baselineskip}
\begin{paracol}{2}
\selectlanguage{latin}
\begin{center}{\color{gregoriocolor} Quarto Kaléndas Octóbris. 
 Luna\dots\ }\end{center}
\switchcolumn
\selectlanguage{english}
\begin{center}{\color{gregoriocolor} The   28\textsuperscript{th} Day of 
 September. The\dots\ Day of the Moon.}\end{center}
\end{paracol}

\noindent\begin{tabularx}{\linewidth}{*{19}{>{\centering\arraybackslash}X}}
 \textcolor{gregoriocolor}{a} & \textcolor{gregoriocolor}{b} & \textcolor{gregoriocolor}{c} & \textcolor{gregoriocolor}{d} & \textcolor{gregoriocolor}{e} & \textcolor{gregoriocolor}{f} & \textcolor{gregoriocolor}{g} & \textcolor{gregoriocolor}{h} & \textcolor{gregoriocolor}{i} & \textcolor{gregoriocolor}{k} & \textcolor{gregoriocolor}{l} & \textcolor{gregoriocolor}{m} & \textcolor{gregoriocolor}{n} & \textcolor{gregoriocolor}{p} & \textcolor{gregoriocolor}{q} & \textcolor{gregoriocolor}{r} & \textcolor{gregoriocolor}{s} & \textcolor{gregoriocolor}{t} & \textcolor{gregoriocolor}{u} \\
 6 & 7 & 8 & 9 & 10 & 11 & 12 & 13 & 14 & 15 & 16 & 17 & 18 & 19 & 20 & 21 & 22 & 23 & 24 \\
\end{tabularx}
\vspace{0.5\baselineskip}
\noindent\begin{tabularx}{\linewidth}{*{12}{>{\centering\arraybackslash}X}}
 \textcolor{gregoriocolor}{A} & \textcolor{gregoriocolor}{B} & \textcolor{gregoriocolor}{C} & \textcolor{gregoriocolor}{D} & \textcolor{gregoriocolor}{E} & F & \textcolor{gregoriocolor}{F} & \textcolor{gregoriocolor}{G} & \textcolor{gregoriocolor}{H} & \textcolor{gregoriocolor}{M} & \textcolor{gregoriocolor}{N} & \textcolor{gregoriocolor}{P} \\
 25 & 26 & 27 & 28 & 29 & 1 & 30 & 1 & 2 & 3 & 4 & 5 \\
\end{tabularx}

\begin{paracol}{2}
\selectlanguage{latin}
\lettrine[lines=2]{A}{pud} Bolesláviam 
 véterem, in Bohémia, sancti Wenceslái, Ducis Bohemórum et Mártyris, 
 sanctitáte et miráculis gloriósi, qui, dolo fratris sui necátus, victor 
 pervénit ad palmam.
\switchcolumn
\selectlanguage{english}
\lettrine[lines=2]{I}{n} Bohemia, St. Wenceslas, duke of 
 Bohemia and martyr, renowned for holiness and miracles. Being murdered 
 by the deceit of his brother, he went triumphantly to heaven.
\switchcolumn*
\selectlanguage{latin}
Romæ sancti Priváti 
 Mártyris, qui, ulcéribus plenus, a beáto Callísto Papa est sanátus; inde, 
 sub Alexándro Imperatóre, ob Christi fidem plumbátis cæsus est usque ad 
 mortem.
\switchcolumn
\selectlanguage{english}
At Rome, St. Privatus, martyr, who 
 was cured of ulcers by blessed Pope Callistus. In the time of Emperor 
 Alexander he was scourged to death with leaded whips for the faith of 
 Christ.
\switchcolumn*
\selectlanguage{latin}
Item Romæ sancti 
 Stáctei Mártyris.
\switchcolumn
\selectlanguage{english}
In the same place, St. Stacteus, 
 martyr.
\switchcolumn*
\selectlanguage{latin}
In Africa sanctórum 
 Mártyrum Martiális, Lauréntii et aliórum vigínti.
\switchcolumn
\selectlanguage{english}
In Africa, the Saints Martial, 
 Lawrence, and twenty other martyrs.
\switchcolumn*
\selectlanguage{latin}
Antiochíæ Pisídiæ 
 sancti Marci Mártyris, pastóris óvium; itémque commemorátio sanctórum Alphíi, 
 Alexándri et Zósimi fratrum, Nicónis, Neónis, Heliodóri et trigínta mílitum, 
 qui, cum ad mirácula beáti Marci credidíssent in Christum, divérsis in locis 
 ac diébus variísque modis martyrio coronáti sunt.
\switchcolumn
\selectlanguage{english}
At Antioch in Pisidia, the holy 
 martyrs Mark, a shepherd, Alphius, Alexander, and Zosimus, his brothers; 
 also Nicon, Neon, Heliodorus, and thirty soldiers, who were converted to 
 Christ upon seeing the miracles of blessed Mark, and were crowned with 
 martyrdom in different places and in diverse manners.
\switchcolumn*
\selectlanguage{latin}
Eódem die pássio sancti 
 Máximi, sub Décio Imperatóre.
\switchcolumn
\selectlanguage{english}
On the same day, under Emperor 
 Decius, the martyrdom of St. Maximus.
\switchcolumn*
\selectlanguage{latin}
Tolósæ sancti Exsupérii, 
 Epíscopi et Confessóris; qui beátus vir quantum sibi parcus exstíterit 
 quantúmque áliis largus, sanctus Hierónymus relátu prosecútus est memorábili.
\switchcolumn
\selectlanguage{english}
At Toulouse, St. Exuperius, bishop 
 and confessor. St. Jerome gives a memorable testimony of this blessed 
 man, relating how severe he was towards himself and how liberal towards 
 others.
\switchcolumn*
\selectlanguage{latin}
Génuæ sancti Salomónis, 
 Epíscopi et Confessóris.
\switchcolumn
\selectlanguage{english}
At Genoa, St. Solomon, bishop and 
 confessor.
\switchcolumn*
\selectlanguage{latin}
Bríxiæ sancti Silvíni 
 Epíscopi.
\switchcolumn
\selectlanguage{english}
At Brescia, St. Silvinus, bishop.
\switchcolumn*
\selectlanguage{latin}
In Béthlehem Judæ 
 sanctæ Eustóchii Vírginis, quæ cum beáta Paula, matre sua, ex urbe Roma in 
 Palæstínam profécta est; ibíque, ad Præsépe Dómini cum áliis Virgínibus 
 enutríta, præcláris méritis fulgens migrávit ad Dóminum.
\switchcolumn
\selectlanguage{english}
At Bethlehem of Juda, the holy 
 virgin Eustochium, daughter of blessed Paula, who was brought up at 
 the manger of our Lord with other virgins, and being celebrated for her 
 merits, went to our Lord.
\switchcolumn*
\selectlanguage{latin}
Schorneshémii, prope 
 Mogúntiam, sanctæ Líobæ Vírginis, miráculis claræ.
\switchcolumn
\selectlanguage{english}
At Fulda near Mayence, St. Lioba, 
 virgin, renowned for miracles.
\switchcolumn*
\selectlanguage{latin}
\end{paracol}


% ---- martyrology/mart09/mart0929.htm
\needspace{10\baselineskip}
\begin{paracol}{2}
\selectlanguage{latin}
\begin{center}{\color{gregoriocolor} Tértio Kaléndas Octóbris. 
 Luna\dots\ }\end{center}
\switchcolumn
\selectlanguage{english}
\begin{center}{\color{gregoriocolor} The   29\textsuperscript{th} Day of 
 September. The\dots\ Day of the Moon.}\end{center}
\end{paracol}

\noindent\begin{tabularx}{\linewidth}{*{19}{>{\centering\arraybackslash}X}}
 \textcolor{gregoriocolor}{a} & \textcolor{gregoriocolor}{b} & \textcolor{gregoriocolor}{c} & \textcolor{gregoriocolor}{d} & \textcolor{gregoriocolor}{e} & \textcolor{gregoriocolor}{f} & \textcolor{gregoriocolor}{g} & \textcolor{gregoriocolor}{h} & \textcolor{gregoriocolor}{i} & \textcolor{gregoriocolor}{k} & \textcolor{gregoriocolor}{l} & \textcolor{gregoriocolor}{m} & \textcolor{gregoriocolor}{n} & \textcolor{gregoriocolor}{p} & \textcolor{gregoriocolor}{q} & \textcolor{gregoriocolor}{r} & \textcolor{gregoriocolor}{s} & \textcolor{gregoriocolor}{t} & \textcolor{gregoriocolor}{u} \\
 7 & 8 & 9 & 10 & 11 & 12 & 13 & 14 & 15 & 16 & 17 & 18 & 19 & 20 & 21 & 22 & 23 & 24 & 25 \\
\end{tabularx}
\vspace{0.5\baselineskip}
\noindent\begin{tabularx}{\linewidth}{*{12}{>{\centering\arraybackslash}X}}
 \textcolor{gregoriocolor}{A} & \textcolor{gregoriocolor}{B} & \textcolor{gregoriocolor}{C} & \textcolor{gregoriocolor}{D} & \textcolor{gregoriocolor}{E} & F & \textcolor{gregoriocolor}{F} & \textcolor{gregoriocolor}{G} & \textcolor{gregoriocolor}{H} & \textcolor{gregoriocolor}{M} & \textcolor{gregoriocolor}{N} & \textcolor{gregoriocolor}{P} \\
 26 & 27 & 28 & 29 & 1 & 2 & 1 & 2 & 3 & 4 & 5 & 6 \\
\end{tabularx}

\begin{paracol}{2}
\selectlanguage{latin}
\lettrine[lines=2]{I}{n} monte Gargáno 
 venerábilis memória beáti Michaélis Archángeli, quando ipsíus nómine ibi 
 consecráta fuit Ecclésia, vili quidem facta schémate, sed cælésti 
 præstans virtúte.
\switchcolumn
\selectlanguage{english}
\lettrine[lines=2]{O}{n} Mount Gargano, the commemoration 
 of the blessed archangel Michael. This festival is kept in memory of 
 the day when, under his invocation, there was consecrated a church, 
 unpretending in its exterior, but endowed with celestial virtue.
\switchcolumn*
\selectlanguage{latin}
Antisiodóri sancti 
 Fratérni, Epíscopi et Mártyris.
\switchcolumn
\selectlanguage{english}
At Auxerre, St. Fraternus, bishop 
 and martyr.
\switchcolumn*
\selectlanguage{latin}
In Thrácia natális 
 sanctórum Mártyrum Eutychii, Plauti et Heracléæ.
\switchcolumn
\selectlanguage{english}
In Thrace, the birthday of the holy 
 martyrs Eutychius, Plautus, and Heracleas.
\switchcolumn*
\selectlanguage{latin}
In Pérside sanctórum 
 Mártyrum Dadæ, e Sáporis Regis consanguíneis, Cásdoæ uxóris, et Gabdélæ 
 fílii; qui, honóribus exúti ac váriis torméntis dilaniáti, tandem, post 
 longos cárceres, gládio sunt animadvérsi.
\switchcolumn
\selectlanguage{english}
In Persia, the holy martyrs Dadas, a 
 blood relative of King Sapor, Casdoa, his wife, and Gabdelas, his son. 
 After being deprived of their dignities, and subjected to various torments, 
 they were imprisoned for a long time and finally put to the sword.
\switchcolumn*
\selectlanguage{latin}
In Arménia sanctárum 
 Vírginum Rípsimis et Sociárum Mártyrum, sub Tiridáte Rege.
\switchcolumn
\selectlanguage{english}
In Armenia, under King Tiridates, 
 the holy virgin Ripsimis and her martyr companions.
\switchcolumn*
\selectlanguage{latin}
In Pérside sanctæ 
 Gudéliæ Mártyris, quæ, cum plúrimos convertísset ad Christum, ac Solem et 
 Ignem adoráre noluísset, ideo, sub Sápore Rege, post multa torménta, cute 
 cápitis detrácta, ligno affíxa, méruit obtinére triúmphum.
\switchcolumn
\selectlanguage{english}
In Persia, under King Sapor, the 
 holy martyr Gudelia. After converting many to the faith, and having 
 refused to adore the sun and the fire, she was subjected to numerous 
 torments. Having the skin torn off her head, and being fastened to a 
 post, she merited an eternal triumph.
\switchcolumn*
\selectlanguage{latin}
In Ponte Curvo, apud 
 Aquínum, sancti Grimoáldi, Presbyteri et Confessóris.
\switchcolumn
\selectlanguage{english}
At Pontecorvo near Aquino, St. 
 Grimoaldus, priest and confessor.
\switchcolumn*
\selectlanguage{latin}
In Palæstína sancti 
 Quiríaci Anachorétæ.
\switchcolumn
\selectlanguage{english}
In Palestine, St. Quiriacus, an 
 anchoret.
\switchcolumn*
\selectlanguage{latin}
\end{paracol}


% ---- martyrology/mart09/mart0930.htm
\needspace{10\baselineskip}
\begin{paracol}{2}
\selectlanguage{latin}
\begin{center}{\color{gregoriocolor} Prídie Kaléndas Octóbris. 
 Luna\dots\ }\end{center}
\switchcolumn
\selectlanguage{english}
\begin{center}{\color{gregoriocolor} The   30\textsuperscript{th} Day of 
 September. The\dots\ Day of the Moon.}\end{center}
\end{paracol}

\noindent\begin{tabularx}{\linewidth}{*{19}{>{\centering\arraybackslash}X}}
 \textcolor{gregoriocolor}{a} & \textcolor{gregoriocolor}{b} & \textcolor{gregoriocolor}{c} & \textcolor{gregoriocolor}{d} & \textcolor{gregoriocolor}{e} & \textcolor{gregoriocolor}{f} & \textcolor{gregoriocolor}{g} & \textcolor{gregoriocolor}{h} & \textcolor{gregoriocolor}{i} & \textcolor{gregoriocolor}{k} & \textcolor{gregoriocolor}{l} & \textcolor{gregoriocolor}{m} & \textcolor{gregoriocolor}{n} & \textcolor{gregoriocolor}{p} & \textcolor{gregoriocolor}{q} & \textcolor{gregoriocolor}{r} & \textcolor{gregoriocolor}{s} & \textcolor{gregoriocolor}{t} & \textcolor{gregoriocolor}{u} \\
 8 & 9 & 10 & 11 & 12 & 13 & 14 & 15 & 16 & 17 & 18 & 19 & 20 & 21 & 22 & 23 & 24 & 25 & 26 \\
\end{tabularx}
\vspace{0.5\baselineskip}
\noindent\begin{tabularx}{\linewidth}{*{12}{>{\centering\arraybackslash}X}}
 \textcolor{gregoriocolor}{A} & \textcolor{gregoriocolor}{B} & \textcolor{gregoriocolor}{C} & \textcolor{gregoriocolor}{D} & \textcolor{gregoriocolor}{E} & F & \textcolor{gregoriocolor}{F} & \textcolor{gregoriocolor}{G} & \textcolor{gregoriocolor}{H} & \textcolor{gregoriocolor}{M} & \textcolor{gregoriocolor}{N} & \textcolor{gregoriocolor}{P} \\
 27 & 28 & 29 & 1 & 2 & 3 & 2 & 3 & 4 & 5 & 6 & 7 \\
\end{tabularx}

\begin{paracol}{2}
\selectlanguage{latin}
\lettrine[lines=2]{I}{n} Béthlehem Judæ 
 deposítio sancti Hierónymi Presbyteri, Confessóris et Ecclésiæ Doctóris, 
 qui, ómnium stúdia litterárum adéptus ac probatórum Monachórum imitátor 
 factus, multa hæresum monstra gládio suæ doctrínæ confódit; demum, cum ad 
 decrépitam usque vixísset ætátem, in pace quiévit, sepultúsque est ad 
 Præsépe Dómini. Ejus corpus, póstea Romam delátum, in Basílica sanctæ 
 Maríæ Majóris cónditum fuit.
\switchcolumn
\selectlanguage{english}
\lettrine[lines=2]{I}{n} Bethlehem of Juda, the death of 
 St. Jerome, priest and doctor of the Church. Excelling in all kinds of 
 learning, he imitated the life of the most approved monks, and disposed of 
 many monstrous heresies with the sword of his doctrine. Having at 
 length reached a very advanced age, he rested in peace and was buried near 
 the manger of our Lord. His body was afterwards transferred to Rome, 
 and placed in the basilica of St. Mary Major.
\switchcolumn*
\selectlanguage{latin}
Romæ natális sancti 
 Francísci Bórgiæ, Sacerdótis et Confessóris; qui Præpósitus Generális fuit 
 Societátis Jesu, ac vitæ asperitáte, oratiónis dono, abdicátis sæculi et 
 recusátis Ecclésiæ dignitátibus, vir memorábilis éxstitit. Ipsíus 
 autem festum sexto Idus Octóbris celebrátur.
\switchcolumn
\selectlanguage{english}
At Rome, the birthday of St. Francis 
 Borgia, priest and confessor. He was the General of the Society of 
 Jesus, and is memorable for his mortification, gift of prayer, the forsaking 
 of the world, and the refusal of ecclesiastical dignities. His feast 
 is observed on the 10th of October.
\switchcolumn*
\selectlanguage{latin}
Lexóvii, in Gállia, 
 item natális sanctæ Terésiæ a Jesu Infánte, ex Ordine Carmelitárum 
 Excalceatórum; quam, vitæ innocéntia et simplicitáte claríssimam, Pius 
 Undécimus, Póntifex Máximus, sanctárum Vírginum albo adscrípsit, peculiárem 
 ómnium Missiónum Patrónam declarávit, ejúsque festum quinto Nonas Octóbris 
 recoléndum esse decrévit.
\switchcolumn
\selectlanguage{english}
At Lisieux in France, the birthday 
 of St. Theresa of the Child Jesus, of the Order of Discalced Carmelites. 
 Seeing her to be most wonderful for her innocence of life and simplicity, 
 Pope Pius XI placed her name among the holy virgins and appointed her as 
 special patron before God of all missions, decreeing that her feast should 
 be observed on the 3rd of October.
\switchcolumn*
\selectlanguage{latin}
Romæ sancti Leopárdi 
 Mártyris, qui ex Juliáni Apóstatæ domésticis fuit; cui caput amputátum est, 
 et corpus ejus Aquisgránum póstea translátum.
\switchcolumn
\selectlanguage{english}
At Rome, the holy martyr Leopardus, 
 of the household of Julian the Apostate. He was beheaded at Rome, and 
 his body afterwards taken to Aix-la-Chapelle.
\switchcolumn*
\selectlanguage{latin}
Solodóri, in Gállia, 
 pássio sanctórum Mártyrum Victóris et Ursi, ex gloriósa legióne Thebæórum, 
 qui, sub Maximiáno Imperatóre, primum diris supplíciis cruciáti, sed cælésti 
 super eos lúmine coruscánte, ruéntibus in terram minístris, erépti; deínde 
 in ignem missi, sed in nullo pénitus læsi; novíssime gládio consummáti sunt.
\switchcolumn
\selectlanguage{english}
At Soleure in Switzerland, in the 
 time of Emperor Maximian, the passion of the holy martyrs Victor and Ursus, 
 of the glorious Theban legion. They were subjected to horrible 
 tortures, but a heavenly light shone over them causing the executioners to 
 fall to the ground, and they were delivered. Being then cast into the 
 fire without sustaining any injury, they finally perished by the sword.
\switchcolumn*
\selectlanguage{latin}
Placéntiæ sancti 
 Antoníni Mártyris, ex eádem legióne.
\switchcolumn
\selectlanguage{english}
At Piacenza, the holy martyr 
 Antoninus, a soldier of the same legion.
\switchcolumn*
\selectlanguage{latin}
Eódem die sancti 
 Gregórii, Epíscopi magnæ Arméniæ, qui, sub Diocletiáno, multa passus est; ac 
 tandem, Constantíni Magni Imperatóris témpore, in pace quiévit.
\switchcolumn
\selectlanguage{english}
On the same day, St. Gregory, bishop 
 of Greater Armenia, who, after many sufferings under Diocletian, rested in 
 peace.
\switchcolumn*
\selectlanguage{latin}
Cantuáriæ, in Anglia, 
 sancti Honórii, Epíscopi et Confessóris.
\switchcolumn
\selectlanguage{english}
At Canterbury in England, St. 
 Honorius, bishop and confessor.
\switchcolumn*
\selectlanguage{latin}
Romæ sanctæ Sophíæ 
 Víduæ, matris sanctárum Vírginum et Mártyrum Fídei, Spei et Caritátis.
\switchcolumn
\selectlanguage{english}
At Rome, St. Sophia, widow, mother 
 of the holy virgin martyrs Faith, Hope, and Charity.
\switchcolumn*
\selectlanguage{latin}
\end{paracol}
